{\begin {center}\textbf{ 1}\end {center}}\par
Файл \code{'1.txt'} містить відомість з рядками, в яких через двокрапку записа\-но поряд\-ковий номер (ціле), назву предмету (знаків пунктуації не містить), прізвище студента, ім'я студента, номер заліковки, оцінку в балах за 100 бальною шкалою.

Білі символи навколо роздільників не допускаються.

Приклад рядка файлу



\begin{verbatim}
13:algebra:surname:name:123456:77
\end{verbatim}


Нижче наведено код, який має залишити у файлі в кожному рядку тільки номер заліковки, назву предмету та оцінку в балах за 100-бальною шкалою (у заданій послі\-дов\-ності).

Рядок з прикладу має бути замінений на такий



\begin{verbatim}
123456:algebra:77
\end{verbatim}




Код:



\begin{verbatim}
regex = re.compile(r'XXX')
with open('1.txt', encoding='utf8') as f:
    lst = f.readlines()
for i in range(len(lst)):
    lst[i] = re.sub(regex, r'YYY', lst[i])
with open('1.txt', 'w', encoding='utf8') as f:
    f.writelines(lst)
\end{verbatim}







Який рядок треба записати на місце \kw{XXX} ?\par
\vspace{5mm}


Який рядок треба записати на місце \kw{YYY} ?\par


\newpage

{\begin {center}\textbf{ 2}\end {center}}\par
Рядки відомості через двокрапку містять порядковий номер (ціле), пріз\-ви\-ще студен\-та, ім'я студента, номер заліковки, оцінку в балах за 100 бальною шкалою.

Білі символи навколо роздільників не допускаються.

Приклад такого рядка присвоєно змінній \kw{s}.



\begin{verbatim}
s = '13:surname:name:123456:77'
\end{verbatim}


Нижче наведено код, за виконання якого для рядка \kw{s} значенням змінної \kw{out} має стати номер заліковки.

Для наведеного прикладу значенням змінної \kw{out} має стати рядок \code{123456} .



\begin{verbatim}
res = re.fullmatch(r'XXX', s)
s = ''
out = r'YYY'
\end{verbatim}


Код має працювати не тільки для конкретного рядка з прикладу, але й для кожного рядка з відомості.




Який рядок треба записати на місце \kw{XXX} ?\par
\vspace{5mm}


Який рядок треба записати на місце \kw{YYY} ?\par


\newpage

{\begin {center}\textbf{ 3}\end {center}}\par
Файл \code{'1.txt'} містить журнал з рядками, в яких через двокрапку записано поряд\-ко\-вий номер запису (ціле), ім'я користувача (складається з латиниці), час події (фор\-мат як у прикладі), тип події (ERROR, WARNING, INFO, STATUS), код події (ціле).

Білі символи навколо роздільників не допускаються.

Приклад рядка журналу



\begin{verbatim}
12:username:2022-05-05 13-26:ERROR:404
\end{verbatim}


Нижче наведено код, який шукає усі події, що відбулись\\ для користувача \code{anonymous}, та зберігає їх типи в змінній \kw{codes}.



\begin{verbatim}
codes = set()
regex = re.compile(r'XXX')
with open('1.txt', encoding='utf8') as f:
    for s in f:
        res = re.fullmatch(regex, s)
        if not res: continue
        s = ''
        codes.add(r'YYY')
\end{verbatim}





Який рядок треба записати на місце \kw{XXX} ?\par
\vspace{5mm}


Який рядок треба записати на місце \kw{YYY} ?\par


\newpage

{\begin {center}\textbf{ 4}\end {center}}\par
Файл \code{'1.txt'} містить відомість з рядками, в яких через кому записано поряд\-ковий номер (ціле), назву предмету (знаків пунктуації не містить), прізвище студента, ім'я студента, номер заліковки, оцінку в балах за 100 бальною шкалою.

Білі символи навколо роздільників не допускаються.

Приклад рядка файлу



\begin{verbatim}
13,algebra,surname,name,123456,77
\end{verbatim}


Нижче наведено код, який має залишити у файлі в кожному рядку тільки номер заліковки та предмет.

Рядок з прикладу має бути замінений на такий



\begin{verbatim}
123456,algebra
\end{verbatim}




Код



\begin{verbatim}
regex = re.compile(r'XXX')
with open('1.txt', encoding='utf8') as f:
    lst = f.readlines()
for i in range(len(lst)):
    lst[i] = re.sub(regex, r'YYY', lst[i])
with open('1.txt', 'w', encoding='utf8') as f:
    f.writelines(lst)
\end{verbatim}









Який рядок треба записати на місце \kw{XXX} ?\par
\vspace{5mm}


Який рядок треба записати на місце \kw{YYY} ?\par


\newpage

{\begin {center}\textbf{ 5}\end {center}}\par
Файл \code{'1.txt'} містить відомість з рядками, в яких через кому записано поряд\-ковий номер (ціле), назву предмету (знаків пунктуації не містить), прізвище студента, ім'я студента, номер заліковки, оцінку в балах за 100 бальною шкалою.

Білі символи навколо роздільників не допускаються.

Приклад рядка файлу



\begin{verbatim}
13,algebra,surname,name,123456,77
\end{verbatim}


Нижче наведено код, який шукає усі назви предметів, що зустрічаються у файлі, та зберігає їх у змінній \id{codes}.



\begin{verbatim}
codes = set()
regex = re.compile(r'XXX')
with open('1.txt', encoding='utf8') as f:
    for s in f:
        res = re.fullmatch(regex, s)
        if not res: continue
        s = ''
        codes.add(r'YYY')
\end{verbatim}





Який рядок треба записати на місце \kw{XXX} ?\par
\vspace{5mm}


Який рядок треба записати на місце \kw{YYY} ?\par


\newpage

{\begin {center}\textbf{ 6}\end {center}}\par
Файл \code{'1.txt'} містить відомість з рядками, в яких через кому записа\-но назву предмету (знаків пунктуації не містить), прізвище студента, ім'я студента, номер залі\-ков\-ки, оцінку в балах за 100 бальною шкалою.

Білі символи навколо роздільників не допускаються.

Приклад рядка файлу



\begin{verbatim}
algebra,surname,name,123456,77
\end{verbatim}


Нижче наведено код, який має замінити у файлі усі номери заліковок на рядок \code{***}.



\begin{verbatim}
regex = re.compile(r'XXX')
with open('1.txt', encoding='utf8') as f:
    lst = f.readlines()
for i in range(len(lst)):
    lst[i] = re.sub(regex, r'YYY', lst[i])
with open('1.txt', 'w', encoding='utf8') as f:
    f.writelines(lst)
\end{verbatim}







Який рядок треба записати на місце \kw{XXX} ?\par
\vspace{5mm}


Який рядок треба записати на місце \kw{YYY} ?\par


\newpage

{\begin {center}\textbf{ 7}\end {center}}\par
Рядки журналу через двокрапку містять тип події (ERROR, WARNING, INFO, STATUS), ім'я користувача (складається з латиниці), час події (фор\-мат як у прикладі), код події (ціле).

Білі символи навколо роздільників не допускаються.

Приклад такого рядка присвоєно змінній \kw{s}.



\begin{verbatim}
s = 'ERROR:username:2022-05-05 13-26:404'
\end{verbatim}


Нижче наведено код, за виконання якого значенням змінної \kw{out} має стати ім'я користувача.

Для наведеного прикладу значенням змінної \kw{out} має стати рядок \code{username}



\begin{verbatim}
res = re.fullmatch(r'XXX', s)
s = ''
out = r'YYY'
\end{verbatim}


Код має працювати не тільки для конкретного рядка з прикладу, але й для кожного рядка з журналу.




Який рядок треба записати на місце \kw{XXX} ?\par
\vspace{5mm}


Який рядок треба записати на місце \kw{YYY} ?\par


\newpage

{\begin {center}\textbf{ 8}\end {center}}\par
Файл \code{'1.txt'} містить відомість з рядками, в яких через кому записа\-но поряд\-ковий номер (ціле), назву предмету (знаків пунктуації не містить), прізвище студента, ім'я студента, номер заліковки, оцінку в балах за 100 бальною шкалою.

Білі символи навколо роздільників не допускаються.

Приклад рядка файлу



\begin{verbatim}
13,algebra,surname,name,123456,77
\end{verbatim}


Нижче наведено код, який має залишити у файлі в кожному рядку тільки номер заліковки, назву предмету та оцінку в балах за 100-бальною шкалою (у заданій послі\-дов\-ності).

Рядок з прикладу має бути замінений на такий



\begin{verbatim}
123456,algebra,77
\end{verbatim}




Код



\begin{verbatim}
regex = re.compile(r'XXX')
with open('1.txt', encoding='utf8') as f:
    lst = f.readlines()
for i in range(len(lst)):
    lst[i] = re.sub(regex, r'YYY', lst[i])
with open('1.txt', 'w', encoding='utf8') as f:
    f.writelines(lst)
\end{verbatim}







Який рядок треба записати на місце \kw{XXX} ?\par
\vspace{5mm}


Який рядок треба записати на місце \kw{YYY} ?\par


\newpage

{\begin {center}\textbf{ 9}\end {center}}\par
Файл \code{'1.txt'} містить журнал з рядками, в яких через кому записано поряд\-ко\-вий номер запису (ціле), ім'я користувача (складається з латиниці), час події (фор\-мат як у прикладі), тип події (ERROR, WARNING, INFO, STATUS), код події (ціле).

Білі символи навколо роздільників не допускаються.

Приклад рядка журналу



\begin{verbatim}
12,username,2022-05-05 13:26,ERROR,404
\end{verbatim}


Нижче наведено код, який шукає усі події, що відбулись\\ для користувача \code{anonymous}, та зберігає їх типи в змінній \kw{codes}.



\begin{verbatim}
codes = set()
regex = re.compile(r'XXX')
with open('1.txt', encoding='utf8') as f:
    for s in f:
        res = re.fullmatch(regex, s)
        if not res: continue
        s = ''
        codes.add(r'YYY')
\end{verbatim}





Який рядок треба записати на місце \kw{XXX} ?\par
\vspace{5mm}


Який рядок треба записати на місце \kw{YYY} ?\par


\newpage

{\begin {center}\textbf{ 10}\end {center}}\par
Файл \code{'1.txt'} містить журнал з рядками, в яких через двокрапку записано поряд\-ко\-вий номер запису в журналі (ціле), тип події (ERROR, WARNING, INFO, STATUS), ім'я користувача (складається з латиниці), час події (фор\-мат як у прикладі), код події (ціле).

Білі символи навколо роздільників не допускаються.

Приклад рядка файлу



\begin{verbatim}
13:ERROR:username:2022-05-05 13-26:404
\end{verbatim}


Нижче наведено код, який має замінити у файлі імена користувачів на рядок \code{unknown}.



\begin{verbatim}
regex = re.compile(r'XXX')
with open('1.txt', encoding='utf8') as f:
    lst = f.readlines()
for i in range(len(lst)):
    lst[i] = re.sub(regex, r'YYY', lst[i])
with open('1.txt', 'w', encoding='utf8') as f:
    f.writelines(lst)
\end{verbatim}







Який рядок треба записати на місце \kw{XXX} ?\par
\vspace{5mm}


Який рядок треба записати на місце \kw{YYY} ?\par


\newpage

{\begin {center}\textbf{ 11}\end {center}}\par
Файл \code{'1.txt'} містить відомість з рядками, в яких через двокрапку записано поряд\-ковий номер (ціле), назву предмету (знаків пунктуації не містить), прізвище студента, ім'я студента, номер заліковки, оцінку в балах за 100 бальною шкалою.

Білі символи навколо роздільників не допускаються.

Приклад рядка файлу



\begin{verbatim}
13:algebra:surname:name:123456:77
\end{verbatim}


Нижче наведено код, який має в кожному рядку файлу переставити місцями назву предмету та номер заліковки.

Рядок з прикладу має бути замінений на такий



\begin{verbatim}
13:123456:surname:name:algebra:77
\end{verbatim}




Код:



\begin{verbatim}
regex = re.compile(r'XXX')
with open('1.txt', encoding='utf8') as f:
    lst = f.readlines()
for i in range(len(lst)):
    lst[i] = re.sub(regex, r'YYY', lst[i])
with open('1.txt', 'w', encoding='utf8') as f:
    f.writelines(lst)
\end{verbatim}







Який рядок треба записати на місце \kw{XXX} ?\par
\vspace{5mm}


Який рядок треба записати на місце \kw{YYY} ?\par


\newpage

{\begin {center}\textbf{ 12}\end {center}}\par
Файл \code{'1.txt'} містить журнал з рядками, в яких через кому записано поряд\-ко\-вий номер запису в журналі (ціле), тип події (ERROR, WARNING, INFO, STATUS), ім'я користувача (складається з латиниці), час події (фор\-мат як у прикладі), опис події.

Білі символи навколо роздільників не допускаються.

Приклад рядка файлу



\begin{verbatim}
13,ERROR,username,2022-05-05 13:26,something happens
\end{verbatim}


Нижче наведено код, який має залишити у файлі в кожному рядку тільки тип події та її опис.

Рядок з прикладу має бути замінений на такий



\begin{verbatim}
ERROR,something happens
\end{verbatim}




Код:



\begin{verbatim}
regex = re.compile(r'XXX')
with open('1.txt', encoding='utf8') as f:
    lst = f.readlines()
for i in range(len(lst)):
    lst[i] = re.sub(regex, r'YYY', lst[i])
with open('1.txt', 'w', encoding='utf8') as f:
    f.writelines(lst)
\end{verbatim}









Який рядок треба записати на місце \kw{XXX} ?\par
\vspace{5mm}


Який рядок треба записати на місце \kw{YYY} ?\par


\newpage

{\begin {center}\textbf{ 13}\end {center}}\par
Файл \code{'1.txt'} містить відомість з рядками, в яких через кому записано поряд\-ковий номер (ціле), назву предмету (знаків пунктуації не містить), прізвище студента, ім'я студента, номер заліковки, оцінку в балах за 100 бальною шкалою.

Білі символи навколо роздільників не допускаються.

Приклад рядка файлу



\begin{verbatim}
13,algebra,surname,name,123456,77
\end{verbatim}


Нижче наведено код, який має в кожному рядку файлу переставити місцями назву предмету та номер заліковки.

Рядок з прикладу має бути замінений на такий



\begin{verbatim}
13,123456,surname,name,algebra,77
\end{verbatim}




Код



\begin{verbatim}
regex = re.compile(r'XXX')
with open('1.txt', encoding='utf8') as f:
    lst = f.readlines()
for i in range(len(lst)):
    lst[i] = re.sub(regex, r'YYY', lst[i])
with open('1.txt', 'w', encoding='utf8') as f:
    f.writelines(lst)
\end{verbatim}







Який рядок треба записати на місце \kw{XXX} ?\par
\vspace{5mm}


Який рядок треба записати на місце \kw{YYY} ?\par


\newpage

{\begin {center}\textbf{ 14}\end {center}}\par
Рядки журналу через кому містять ім'я користувача (складається з латиниці), час події (фор\-мат як у прикладі), тип події (ERROR, WARNING, INFO, STATUS), код події (ціле).

Білі символи навколо роздільників не допускаються.

Приклад такого рядка присвоєно змінній \kw{s}.



\begin{verbatim}
s = 'username,2022-05-05 13:26,ERROR,404'
\end{verbatim}


Нижче наведено код, за виконання якого значенням змінної \kw{out} має стати код події.

Для наведеного прикладу значенням змінної \kw{out} має стати рядок \code{404}



\begin{verbatim}
res = re.fullmatch(r'XXX', s)
s = ''
out = r'YYY'
\end{verbatim}


Код має працювати не тільки для конкретного рядка з прикладу, але й для кожного рядка з журналу.




Який рядок треба записати на місце \kw{XXX} ?\par
\vspace{5mm}


Який рядок треба записати на місце \kw{YYY} ?\par


\newpage

{\begin {center}\textbf{ 15}\end {center}}\par
Файл \code{'1.txt'} містить журнал з рядками, в яких через двокрапку записано поряд\-ко\-вий номер запису (ціле), ім'я користувача (складається з латиниці), час події (фор\-мат як у прикладі), тип події (ERROR, WARNING, INFO, STATUS), код події (ціле).

Білі символи навколо роздільників не допускаються.

Приклад рядка журналу



\begin{verbatim}
12:username:2022-05-05 13-26:ERROR:404
\end{verbatim}


Нижче наведено код, який шукає усі попередження (подія WARNING), що зустрі\-ча\-ються в журналі і відбулися під час обідньої перерви деканату (з 13:00 по 13:59 включно), та зберігає їх коди в змінній \kw{codes}.



\begin{verbatim}
codes = set()
regex = re.compile(r'XXX')
with open('1.txt', encoding='utf8') as f:
    for s in f:
        res = re.fullmatch(regex, s)
        if not res: continue
        s = ''
        codes.add(r'YYY')
\end{verbatim}







Який рядок треба записати на місце \kw{XXX} ?\par
\vspace{5mm}


Який рядок треба записати на місце \kw{YYY} ?\par


\newpage

{\begin {center}\textbf{ 16}\end {center}}\par
Рядки відомості через кому містять порядковий номер (ціле), назву предмету (зна\-ків пунктуації не містить), прізвище студента, ім'я студента, оцінку в балах за 100 бальною шкалою.

Білі символи навколо роздільників не допускаються.

Приклад такого рядка присвоєно змінній \kw{s}.



\begin{verbatim}
s = '13,algebra,surname,name,77'
\end{verbatim}


Нижче наведено код, за виконання якого значенням змінної \kw{out} має стати прізвище студента.

Для наведеного прикладу значенням змінної \kw{out} має стати рядок \code{surname}



\begin{verbatim}
res = re.fullmatch(r'XXX', s)
s = ''
out = r'YYY'
\end{verbatim}


Код має працювати не тільки для конкретного рядка з прикладу, але й для кожного рядка з відомості.




Який рядок треба записати на місце \kw{XXX} ?\par
\vspace{5mm}


Який рядок треба записати на місце \kw{YYY} ?\par


\newpage

{\begin {center}\textbf{ 17}\end {center}}\par
Рядки відомості через кому містять порядковий номер (ціле), пріз\-ви\-ще студен\-та, ім'я студента, номер заліковки, оцінку в балах за 100 бальною шкалою.

Білі символи навколо роздільників не допускаються.

Приклад такого рядка присвоєно змінній \kw{s}.



\begin{verbatim}
s = '13,surname,name,123456,77'
\end{verbatim}


Нижче наведено код, за виконання якого для рядка \kw{s} значенням змінної \kw{out} має стати номер заліковки.

Для наведеного прикладу значенням змінної \kw{out} має стати рядок \code{123456} .



\begin{verbatim}
res = re.fullmatch(r'XXX', s)
s = ''
out = r'YYY'
\end{verbatim}


Код має працювати не тільки для конкретного рядка з прикладу, але й для кожного рядка з відомості.




Який рядок треба записати на місце \kw{XXX} ?\par
\vspace{5mm}


Який рядок треба записати на місце \kw{YYY} ?\par


\newpage

{\begin {center}\textbf{ 18}\end {center}}\par
Файл \code{'1.txt'} містить відомість з рядками, в яких через двокрапку записано поряд\-ковий номер (ціле), назву предмету (знаків пунктуації не містить), прізвище студента, ім'я студента, номер заліковки, оцінку в балах за 100 бальною шкалою.

Білі символи навколо роздільників не допускаються.

Приклад рядка файлу



\begin{verbatim}
13:algebra:surname:name:123456:77
\end{verbatim}


Нижче наведено код, який шукає усі предмети, який складав студент із заліковкою \code{777}, та зберігає їх типи в змінній \kw{codes}.



\begin{verbatim}
codes = set()
regex = re.compile(r'XXX')
with open('1.txt', encoding='utf8') as f:
    for s in f:
        res = re.fullmatch(regex, s)
        if not res: continue
        s = ''
        codes.add(r'YYY')
\end{verbatim}





Який рядок треба записати на місце \kw{XXX} ?\par
\vspace{5mm}


Який рядок треба записати на місце \kw{YYY} ?\par


\newpage

{\begin {center}\textbf{ 19}\end {center}}\par
Рядки відомості через двокрапку містять порядковий номер (ціле), назву предмету (знаків пунктуації не містить), прізвище студента, ім'я студента, оцінку в балах за 100 бальною шкалою.

Білі символи навколо роздільників не допускаються.

Приклад такого рядка присвоєно змінній \kw{s}.



\begin{verbatim}
s = '13:algebra:surname:name:77'
\end{verbatim}


Нижче наведено код, за виконання якого значенням змінної \kw{out} має стати прізвище студента.

Для наведеного прикладу значенням змінної \kw{out} має стати рядок \code{surname}



\begin{verbatim}
res = re.fullmatch(r'XXX', s)
s = ''
out = r'YYY'
\end{verbatim}


Код має працювати не тільки для конкретного рядка з прикладу, але й для кожного рядка з відомості.




Який рядок треба записати на місце \kw{XXX} ?\par
\vspace{5mm}


Який рядок треба записати на місце \kw{YYY} ?\par


\newpage

{\begin {center}\textbf{ 20}\end {center}}\par
Файл \code{'1.txt'} містить журнал з рядками, в яких через кому записано поряд\-ко\-вий номер запису в журналі (ціле), тип події (ERROR, WARNING, INFO, STATUS), ім'я користувача (складається з латиниці), час події (фор\-мат як у прикладі), код події (ціле).

Білі символи навколо роздільників не допускаються.

Приклад рядка файлу



\begin{verbatim}
13,ERROR,username,2022-05-05 13:26,404
\end{verbatim}


Нижче наведено код, який шукає усі коди помилок (подія ERROR), що зустріча\-ють\-ся в журналі, та зберігає їх у змінній \kw{codes}.



\begin{verbatim}
codes = set()
regex = re.compile(r'XXX')
with open('1.txt', encoding='utf8') as f:
    for s in f:
        res = re.fullmatch(regex, s)
        if not res: continue
        s = ''
        codes.add(r'YYY')
\end{verbatim}





Який рядок треба записати на місце \kw{XXX} ?\par
\vspace{5mm}


Який рядок треба записати на місце \kw{YYY} ?\par


\newpage

{\begin {center}\textbf{ 21}\end {center}}\par
Рядки журналу через двокрапку містять ім'я користувача (складається з латиниці), час події (фор\-мат як у прикладі), тип події (ERROR, WARNING, INFO, STATUS), код події (ціле).

Білі символи навколо роздільників не допускаються.

Приклад такого рядка присвоєно змінній \kw{s}.



\begin{verbatim}
s = 'username:2022-05-05 13-26:ERROR:404'
\end{verbatim}


Нижче наведено код, за виконання якого значенням змінної \kw{out} має стати код події.

Для наведеного прикладу значенням змінної \kw{out} має стати рядок \code{404}



\begin{verbatim}
res = re.fullmatch(r'XXX', s)
s = ''
out = r'YYY'
\end{verbatim}


Код має працювати не тільки для конкретного рядка з прикладу, але й для кожного рядка з журналу.




Який рядок треба записати на місце \kw{XXX} ?\par
\vspace{5mm}


Який рядок треба записати на місце \kw{YYY} ?\par


\newpage

{\begin {center}\textbf{ 22}\end {center}}\par
Файл \code{'1.txt'} містить журнал з рядками, в яких через двокрапку записано поряд\-ко\-вий номер запису в журналі (ціле), тип події (ERROR, WARNING, INFO, STATUS), ім'я користувача (складається з латиниці), час події (фор\-мат як у прикладі), код події (ціле).

Білі символи навколо роздільників не допускаються.

Приклад рядка файлу



\begin{verbatim}
13:ERROR:username:2022-05-05 13-26:404
\end{verbatim}


Нижче наведено код, який шукає усі коди помилок (подія ERROR), що зустріча\-ють\-ся в журналі, та зберігає їх у змінній \kw{codes}.



\begin{verbatim}
codes = set()
regex = re.compile(r'XXX')
with open('1.txt', encoding='utf8') as f:
    for s in f:
        res = re.fullmatch(regex, s)
        if not res: continue
        s = ''
        codes.add(r'YYY')
\end{verbatim}





Який рядок треба записати на місце \kw{XXX} ?\par
\vspace{5mm}


Який рядок треба записати на місце \kw{YYY} ?\par


\newpage

{\begin {center}\textbf{ 23}\end {center}}\par
Файл \code{'1.txt'} містить відомість з рядками, в яких через двокрапку записа\-но назву предмету (знаків пунктуації не містить), прізвище студента, ім'я студента, номер залі\-ков\-ки, оцінку в балах за 100 бальною шкалою.

Білі символи навколо роздільників не допускаються.

Приклад рядка файлу



\begin{verbatim}
algebra:surname:name:123456:77
\end{verbatim}


Нижче наведено код, який має замінити у файлі усі номери заліковок на рядок \code{***}.



\begin{verbatim}
regex = re.compile(r'XXX')
with open('1.txt', encoding='utf8') as f:
    lst = f.readlines()
for i in range(len(lst)):
    lst[i] = re.sub(regex, r'YYY', lst[i])
with open('1.txt', 'w', encoding='utf8') as f:
    f.writelines(lst)
\end{verbatim}







Який рядок треба записати на місце \kw{XXX} ?\par
\vspace{5mm}


Який рядок треба записати на місце \kw{YYY} ?\par


\newpage

{\begin {center}\textbf{ 24}\end {center}}\par
Файл \code{'1.txt'} містить журнал з рядками, в яких через кому записано поряд\-ко\-вий номер запису (ціле), ім'я користувача (складається з латиниці), час події (фор\-мат як у прикладі), тип події (ERROR, WARNING, INFO, STATUS), код події (ціле).

Білі символи навколо роздільників не допускаються.

Приклад рядка журналу



\begin{verbatim}
12,username,2022-05-05 13:26,ERROR,404
\end{verbatim}


Нижче наведено код, який шукає усі попередження (подія WARNING), що зустрі\-ча\-ються в журналі і відбулися під час обідньої перерви деканату (з 13,00 по 13,59 включно), та зберігає їх коди в змінній \kw{codes}.



\begin{verbatim}
codes = set()
regex = re.compile(r'XXX')
with open('1.txt', encoding='utf8') as f:
    for s in f:
        res = re.fullmatch(regex, s)
        if not res: continue
        s = ''
        codes.add(r'YYY')
\end{verbatim}







Який рядок треба записати на місце \kw{XXX} ?\par
\vspace{5mm}


Який рядок треба записати на місце \kw{YYY} ?\par


\newpage

{\begin {center}\textbf{ 25}\end {center}}\par
Файл \code{'1.txt'} містить відомість з рядками, в яких через кому записано поряд\-ковий номер (ціле), назву предмету (знаків пунктуації не містить), прізвище студента, ім'я студента, номер заліковки, оцінку в балах за 100 бальною шкалою.

Білі символи навколо роздільників не допускаються.

Приклад рядка файлу



\begin{verbatim}
13,algebra,surname,name,123456,77
\end{verbatim}


Нижче наведено код, який шукає усі предмети, який складав студент із заліковкою \code{777}, та зберігає їх типи в змінній \kw{codes}.



\begin{verbatim}
codes = set()
regex = re.compile(r'XXX')
with open('1.txt', encoding='utf8') as f:
    for s in f:
        res = re.fullmatch(regex, s)
        if not res: continue
        s = ''
        codes.add(r'YYY')
\end{verbatim}





Який рядок треба записати на місце \kw{XXX} ?\par
\vspace{5mm}


Який рядок треба записати на місце \kw{YYY} ?\par


\newpage

{\begin {center}\textbf{ 26}\end {center}}\par
Файл \code{'1.txt'} містить журнал з рядками, в яких через кому записано поряд\-ко\-вий номер запису в журналі (ціле), тип події (ERROR, WARNING, INFO, STATUS), ім'я користувача (складається з латиниці), час події (фор\-мат як у прикладі), код події (ціле).

Білі символи навколо роздільників не допускаються.

Приклад рядка файлу



\begin{verbatim}
13,ERROR,username,2022-05-05 13:26,404
\end{verbatim}


Нижче наведено код, який має замінити у файлі імена користувачів на рядок \code{unknown}.



\begin{verbatim}
regex = re.compile(r'XXX')
with open('1.txt', encoding='utf8') as f:
    lst = f.readlines()
for i in range(len(lst)):
    lst[i] = re.sub(regex, r'YYY', lst[i])
with open('1.txt', 'w', encoding='utf8') as f:
    f.writelines(lst)
\end{verbatim}







Який рядок треба записати на місце \kw{XXX} ?\par
\vspace{5mm}


Який рядок треба записати на місце \kw{YYY} ?\par


\newpage

{\begin {center}\textbf{ 27}\end {center}}\par
Файл \code{'1.txt'} містить журнал з рядками, в яких через двокрапку записано поряд\-ко\-вий номер запису в журналі (ціле), код події (ціле), тип події (ERROR, \\WARNING, INFO, STATUS), ім'я користувача (складається з латиниці), час події (фор\-мат як у прикладі), опис події.

Білі символи навколо роздільників не допускаються.

Приклад рядка файлу



\begin{verbatim}
13:404:ERROR:username:2022-05-05 13-26:something happens
\end{verbatim}


Нижче наведено код, який має залишити у файлі в кожному рядку тільки код події та ім'я користувача. Рядок з прикладу має бути замінений на такий



\begin{verbatim}
404:username
\end{verbatim}




Код:



\begin{verbatim}
regex = re.compile(r'XXX')
with open('1.txt', encoding='utf8') as f:
    lst = f.readlines()
for i in range(len(lst)):
    lst[i] = re.sub(regex, r'YYY', lst[i])
with open('1.txt', 'w', encoding='utf8') as f:
    f.writelines(lst)
\end{verbatim}







Який рядок треба записати на місце \kw{XXX} ?\par
\vspace{5mm}


Який рядок треба записати на місце \kw{YYY} ?\par


\newpage

{\begin {center}\textbf{ 28}\end {center}}\par
Рядки журналу через кому містять тип події (ERROR, WARNING, INFO, STATUS), ім'я користувача (складається з латиниці), час події (фор\-мат як у прикладі), код події (ціле).

Білі символи навколо роздільників не допускаються.

Приклад такого рядка присвоєно змінній \kw{s}.



\begin{verbatim}
s = 'ERROR,username,2022-05-05 13:26,404'
\end{verbatim}


Нижче наведено код, за виконання якого значенням змінної \kw{out} має стати ім'я користувача.

Для наведеного прикладу значенням змінної \kw{out} має стати рядок \code{username}



\begin{verbatim}
res = re.fullmatch(r'XXX', s)
s = ''
out = r'YYY'
\end{verbatim}


Код має працювати не тільки для конкретного рядка з прикладу, але й для кожного рядка з журналу.




Який рядок треба записати на місце \kw{XXX} ?\par
\vspace{5mm}


Який рядок треба записати на місце \kw{YYY} ?\par


\newpage

{\begin {center}\textbf{ 29}\end {center}}\par
Файл \code{'1.txt'} містить журнал з рядками, в яких через кому записано поряд\-ко\-вий номер запису в журналі (ціле), код події (ціле), тип події (ERROR,\\ WARNING, INFO, STATUS), ім'я користувача (складається з латиниці), час події (фор\-мат як у прикладі), опис події.

Білі символи навколо роздільників не допускаються.

Приклад рядка файлу



\begin{verbatim}
13,404,ERROR,username,2022-05-05 13:26,something happens
\end{verbatim}


Нижче наведено код, який має в кожному рядку файли переставити місцями код та тип події.

Рядок з прикладу має бути замінений на такий



\begin{verbatim}
13,ERROR,404,username,2022-05-05 13:26,something happens
\end{verbatim}




Код:



\begin{verbatim}
regex = re.compile(r'XXX')
with open('1.txt', encoding='utf8') as f:
    lst = f.readlines()
for i in range(len(lst)):
    lst[i] = re.sub(regex, r'YYY', lst[i])
with open('1.txt', 'w', encoding='utf8') as f:
    f.writelines(lst)
\end{verbatim}







Який рядок треба записати на місце \kw{XXX} ?\par
\vspace{5mm}


Який рядок треба записати на місце \kw{YYY} ?\par


\newpage

{\begin {center}\textbf{ 30}\end {center}}\par
Файл \code{'1.txt'} містить журнал з рядками, в яких через двокрапку записано поряд\-ко\-вий номер запису в журналі (ціле), тип події (ERROR, WARNING, INFO, STATUS), ім'я користувача (складається з латиниці), час події (фор\-мат як у прикладі), опис події.

Білі символи навколо роздільників не допускаються.

Приклад рядка файлу



\begin{verbatim}
13:ERROR:username:2022-05-05 13-26:something happens
\end{verbatim}


Нижче наведено код, який має залишити у файлі в кожному рядку тільки тип події та її опис.

Рядок з прикладу має бути замінений на такий



\begin{verbatim}
ERROR:something happens
\end{verbatim}




Код:



\begin{verbatim}
regex = re.compile(r'XXX')
with open('1.txt', encoding='utf8') as f:
    lst = f.readlines()
for i in range(len(lst)):
    lst[i] = re.sub(regex, r'YYY', lst[i])
with open('1.txt', 'w', encoding='utf8') as f:
    f.writelines(lst)
\end{verbatim}









Який рядок треба записати на місце \kw{XXX} ?\par
\vspace{5mm}


Який рядок треба записати на місце \kw{YYY} ?\par


\newpage

{\begin {center}\textbf{ 31}\end {center}}\par
Файл \code{'1.txt'} містить журнал з рядками, в яких через кому записано поряд\-ко\-вий номер запису в журналі (ціле), код події (ціле), тип події (ERROR, \\WARNING, INFO, STATUS), ім'я користувача (складається з латиниці), час події (фор\-мат як у прикладі), опис події.

Білі символи навколо роздільників не допускаються.

Приклад рядка файлу



\begin{verbatim}
13,404,ERROR,username,2022-05-05 13:26,something happens
\end{verbatim}


Нижче наведено код, який має залишити у файлі в кожному рядку тільки код події та ім'я користувача. Рядок з прикладу має бути замінений на такий



\begin{verbatim}
404,username
\end{verbatim}




Код:



\begin{verbatim}
regex = re.compile(r'XXX')
with open('1.txt', encoding='utf8') as f:
    lst = f.readlines()
for i in range(len(lst)):
    lst[i] = re.sub(regex, r'YYY', lst[i])
with open('1.txt', 'w', encoding='utf8') as f:
    f.writelines(lst)
\end{verbatim}







Який рядок треба записати на місце \kw{XXX} ?\par
\vspace{5mm}


Який рядок треба записати на місце \kw{YYY} ?\par


\newpage

{\begin {center}\textbf{ 32}\end {center}}\par
Файл \code{'1.txt'} містить відомість з рядками, в яких через двокрапку записано поряд\-ковий номер (ціле), назву предмету (знаків пунктуації не містить), прізвище студента, ім'я студента, номер заліковки, оцінку в балах за 100 бальною шкалою.

Білі символи навколо роздільників не допускаються.

Приклад рядка файлу



\begin{verbatim}
13:algebra:surname:name:123456:77
\end{verbatim}


Нижче наведено код, який має залишити у файлі в кожному рядку тільки номер заліковки та предмет.

Рядок з прикладу має бути замінений на такий



\begin{verbatim}
123456:algebra
\end{verbatim}




Код:



\begin{verbatim}
regex = re.compile(r'XXX')
with open('1.txt', encoding='utf8') as f:
    lst = f.readlines()
for i in range(len(lst)):
    lst[i] = re.sub(regex, r'YYY', lst[i])
with open('1.txt', 'w', encoding='utf8') as f:
    f.writelines(lst)
\end{verbatim}









Який рядок треба записати на місце \kw{XXX} ?\par
\vspace{5mm}


Який рядок треба записати на місце \kw{YYY} ?\par


\newpage

{\begin {center}\textbf{ 33}\end {center}}\par
Файл \code{'1.txt'} містить журнал з рядками, в яких через двокрапку записано поряд\-ко\-вий номер запису в журналі (ціле), код події (ціле), тип події (ERROR,\\ WARNING, INFO, STATUS), ім'я користувача (складається з латиниці), час події (фор\-мат як у прикладі), опис події.

Білі символи навколо роздільників не допускаються.

Приклад рядка файлу



\begin{verbatim}
13:404:ERROR:username:2022-05-05 13:26:something happens
\end{verbatim}


Нижче наведено код, який має в кожному рядку файли переставити місцями код та тип події.

Рядок з прикладу має бути замінений на такий



\begin{verbatim}
13:ERROR:404:username:2022-05-05 13:26:something happens
\end{verbatim}




Код:



\begin{verbatim}
regex = re.compile(r'XXX')
with open('1.txt', encoding='utf8') as f:
    lst = f.readlines()
for i in range(len(lst)):
    lst[i] = re.sub(regex, r'YYY', lst[i])
with open('1.txt', 'w', encoding='utf8') as f:
    f.writelines(lst)
\end{verbatim}







Який рядок треба записати на місце \kw{XXX} ?\par
\vspace{5mm}


Який рядок треба записати на місце \kw{YYY} ?\par


\newpage

{\begin {center}\textbf{ 34}\end {center}}\par
Файл \code{'1.txt'} містить відомість з рядками, в яких через двокрапку записано поряд\-ковий номер (ціле), назву предмету (знаків пунктуації не містить), прізвище студента, ім'я студента, номер заліковки, оцінку в балах за 100 бальною шкалою.

Білі символи навколо роздільників не допускаються.

Приклад рядка файлу



\begin{verbatim}
13:algebra:surname:name:123456:77
\end{verbatim}


Нижче наведено код, який шукає усі назви предметів, що зустрічаються у файлі, та зберігає їх у змінній \id{codes}.



\begin{verbatim}
codes = set()
regex = re.compile(r'XXX')
with open('1.txt', encoding='utf8') as f:
    for s in f:
        res = re.fullmatch(regex, s)
        if not res: continue
        s = ''
        codes.add(r'YYY')
\end{verbatim}





Який рядок треба записати на місце \kw{XXX} ?\par
\vspace{5mm}


Який рядок треба записати на місце \kw{YYY} ?\par


\newpage

{\begin {center}\textbf{ 35}\end {center}}\par
Файл \code{'1.txt'} містить відомість з рядками, в яких через двокрапку записа\-но назву предмету (знаків пунктуації не містить), прізвище студента, ім'я студента, номер залі\-ков\-ки, оцінку в балах за 100 бальною шкалою.

Білі символи навколо роздільників не допускаються.

Приклад рядка файлу



\begin{verbatim}
algebra:surname:name:123456:77
\end{verbatim}


Нижче наведено код, який має замінити у файлі усі номери заліковок на рядок \code{***}.



\begin{verbatim}
regex = re.compile(r'XXX')
with open('1.txt', encoding='utf8') as f:
    lst = f.readlines()
for i in range(len(lst)):
    lst[i] = re.sub(regex, r'YYY', lst[i])
with open('1.txt', 'w', encoding='utf8') as f:
    f.writelines(lst)
\end{verbatim}







Який рядок треба записати на місце \kw{XXX} ?\par
\vspace{5mm}


Який рядок треба записати на місце \kw{YYY} ?\par


\newpage

{\begin {center}\textbf{ 36}\end {center}}\par
Файл \code{'1.txt'} містить журнал з рядками, в яких через двокрапку записано поряд\-ко\-вий номер запису в журналі (ціле), тип події (ERROR, WARNING, INFO, STATUS), ім'я користувача (складається з латиниці), час події (фор\-мат як у прикладі), опис події.

Білі символи навколо роздільників не допускаються.

Приклад рядка файлу



\begin{verbatim}
13:ERROR:username:2022-05-05 13-26:something happens
\end{verbatim}


Нижче наведено код, який має залишити у файлі в кожному рядку тільки тип події та її опис.

Рядок з прикладу має бути замінений на такий



\begin{verbatim}
ERROR:something happens
\end{verbatim}




Код:



\begin{verbatim}
regex = re.compile(r'XXX')
with open('1.txt', encoding='utf8') as f:
    lst = f.readlines()
for i in range(len(lst)):
    lst[i] = re.sub(regex, r'YYY', lst[i])
with open('1.txt', 'w', encoding='utf8') as f:
    f.writelines(lst)
\end{verbatim}









Який рядок треба записати на місце \kw{XXX} ?\par
\vspace{5mm}


Який рядок треба записати на місце \kw{YYY} ?\par


\newpage

{\begin {center}\textbf{ 37}\end {center}}\par
Файл \code{'1.txt'} містить журнал з рядками, в яких через кому записано поряд\-ко\-вий номер запису в журналі (ціле), тип події (ERROR, WARNING, INFO, STATUS), ім'я користувача (складається з латиниці), час події (фор\-мат як у прикладі), опис події.

Білі символи навколо роздільників не допускаються.

Приклад рядка файлу



\begin{verbatim}
13,ERROR,username,2022-05-05 13:26,something happens
\end{verbatim}


Нижче наведено код, який має залишити у файлі в кожному рядку тільки тип події та її опис.

Рядок з прикладу має бути замінений на такий



\begin{verbatim}
ERROR,something happens
\end{verbatim}




Код:



\begin{verbatim}
regex = re.compile(r'XXX')
with open('1.txt', encoding='utf8') as f:
    lst = f.readlines()
for i in range(len(lst)):
    lst[i] = re.sub(regex, r'YYY', lst[i])
with open('1.txt', 'w', encoding='utf8') as f:
    f.writelines(lst)
\end{verbatim}









Який рядок треба записати на місце \kw{XXX} ?\par
\vspace{5mm}


Який рядок треба записати на місце \kw{YYY} ?\par


\newpage

{\begin {center}\textbf{ 38}\end {center}}\par
Файл \code{'1.txt'} містить журнал з рядками, в яких через кому записано поряд\-ко\-вий номер запису в журналі (ціле), код події (ціле), тип події (ERROR, \\WARNING, INFO, STATUS), ім'я користувача (складається з латиниці), час події (фор\-мат як у прикладі), опис події.

Білі символи навколо роздільників не допускаються.

Приклад рядка файлу



\begin{verbatim}
13,404,ERROR,username,2022-05-05 13:26,something happens
\end{verbatim}


Нижче наведено код, який має залишити у файлі в кожному рядку тільки код події та ім'я користувача. Рядок з прикладу має бути замінений на такий



\begin{verbatim}
404,username
\end{verbatim}




Код:



\begin{verbatim}
regex = re.compile(r'XXX')
with open('1.txt', encoding='utf8') as f:
    lst = f.readlines()
for i in range(len(lst)):
    lst[i] = re.sub(regex, r'YYY', lst[i])
with open('1.txt', 'w', encoding='utf8') as f:
    f.writelines(lst)
\end{verbatim}







Який рядок треба записати на місце \kw{XXX} ?\par
\vspace{5mm}


Який рядок треба записати на місце \kw{YYY} ?\par


\newpage

{\begin {center}\textbf{ 39}\end {center}}\par
Файл \code{'1.txt'} містить відомість з рядками, в яких через кому записано поряд\-ковий номер (ціле), назву предмету (знаків пунктуації не містить), прізвище студента, ім'я студента, номер заліковки, оцінку в балах за 100 бальною шкалою.

Білі символи навколо роздільників не допускаються.

Приклад рядка файлу



\begin{verbatim}
13,algebra,surname,name,123456,77
\end{verbatim}


Нижче наведено код, який має в кожному рядку файлу переставити місцями назву предмету та номер заліковки.

Рядок з прикладу має бути замінений на такий



\begin{verbatim}
13,123456,surname,name,algebra,77
\end{verbatim}




Код



\begin{verbatim}
regex = re.compile(r'XXX')
with open('1.txt', encoding='utf8') as f:
    lst = f.readlines()
for i in range(len(lst)):
    lst[i] = re.sub(regex, r'YYY', lst[i])
with open('1.txt', 'w', encoding='utf8') as f:
    f.writelines(lst)
\end{verbatim}







Який рядок треба записати на місце \kw{XXX} ?\par
\vspace{5mm}


Який рядок треба записати на місце \kw{YYY} ?\par


\newpage

{\begin {center}\textbf{ 40}\end {center}}\par
Файл \code{'1.txt'} містить журнал з рядками, в яких через кому записано поряд\-ко\-вий номер запису (ціле), ім'я користувача (складається з латиниці), час події (фор\-мат як у прикладі), тип події (ERROR, WARNING, INFO, STATUS), код події (ціле).

Білі символи навколо роздільників не допускаються.

Приклад рядка журналу



\begin{verbatim}
12,username,2022-05-05 13:26,ERROR,404
\end{verbatim}


Нижче наведено код, який шукає усі події, що відбулись\\ для користувача \code{anonymous}, та зберігає їх типи в змінній \kw{codes}.



\begin{verbatim}
codes = set()
regex = re.compile(r'XXX')
with open('1.txt', encoding='utf8') as f:
    for s in f:
        res = re.fullmatch(regex, s)
        if not res: continue
        s = ''
        codes.add(r'YYY')
\end{verbatim}





Який рядок треба записати на місце \kw{XXX} ?\par
\vspace{5mm}


Який рядок треба записати на місце \kw{YYY} ?\par


\newpage

{\begin {center}\textbf{ 41}\end {center}}\par
Рядки відомості через двокрапку містять порядковий номер (ціле), пріз\-ви\-ще студен\-та, ім'я студента, номер заліковки, оцінку в балах за 100 бальною шкалою.

Білі символи навколо роздільників не допускаються.

Приклад такого рядка присвоєно змінній \kw{s}.



\begin{verbatim}
s = '13:surname:name:123456:77'
\end{verbatim}


Нижче наведено код, за виконання якого для рядка \kw{s} значенням змінної \kw{out} має стати номер заліковки.

Для наведеного прикладу значенням змінної \kw{out} має стати рядок \code{123456} .



\begin{verbatim}
res = re.fullmatch(r'XXX', s)
s = ''
out = r'YYY'
\end{verbatim}


Код має працювати не тільки для конкретного рядка з прикладу, але й для кожного рядка з відомості.




Який рядок треба записати на місце \kw{XXX} ?\par
\vspace{5mm}


Який рядок треба записати на місце \kw{YYY} ?\par


\newpage

{\begin {center}\textbf{ 42}\end {center}}\par
Файл \code{'1.txt'} містить журнал з рядками, в яких через двокрапку записано поряд\-ко\-вий номер запису в журналі (ціле), тип події (ERROR, WARNING, INFO, STATUS), ім'я користувача (складається з латиниці), час події (фор\-мат як у прикладі), код події (ціле).

Білі символи навколо роздільників не допускаються.

Приклад рядка файлу



\begin{verbatim}
13:ERROR:username:2022-05-05 13-26:404
\end{verbatim}


Нижче наведено код, який шукає усі коди помилок (подія ERROR), що зустріча\-ють\-ся в журналі, та зберігає їх у змінній \kw{codes}.



\begin{verbatim}
codes = set()
regex = re.compile(r'XXX')
with open('1.txt', encoding='utf8') as f:
    for s in f:
        res = re.fullmatch(regex, s)
        if not res: continue
        s = ''
        codes.add(r'YYY')
\end{verbatim}





Який рядок треба записати на місце \kw{XXX} ?\par
\vspace{5mm}


Який рядок треба записати на місце \kw{YYY} ?\par


\newpage

{\begin {center}\textbf{ 43}\end {center}}\par
Рядки журналу через двокрапку містять тип події (ERROR, WARNING, INFO, STATUS), ім'я користувача (складається з латиниці), час події (фор\-мат як у прикладі), код події (ціле).

Білі символи навколо роздільників не допускаються.

Приклад такого рядка присвоєно змінній \kw{s}.



\begin{verbatim}
s = 'ERROR:username:2022-05-05 13-26:404'
\end{verbatim}


Нижче наведено код, за виконання якого значенням змінної \kw{out} має стати ім'я користувача.

Для наведеного прикладу значенням змінної \kw{out} має стати рядок \code{username}



\begin{verbatim}
res = re.fullmatch(r'XXX', s)
s = ''
out = r'YYY'
\end{verbatim}


Код має працювати не тільки для конкретного рядка з прикладу, але й для кожного рядка з журналу.




Який рядок треба записати на місце \kw{XXX} ?\par
\vspace{5mm}


Який рядок треба записати на місце \kw{YYY} ?\par


\newpage

{\begin {center}\textbf{ 44}\end {center}}\par
Файл \code{'1.txt'} містить відомість з рядками, в яких через двокрапку записано поряд\-ковий номер (ціле), назву предмету (знаків пунктуації не містить), прізвище студента, ім'я студента, номер заліковки, оцінку в балах за 100 бальною шкалою.

Білі символи навколо роздільників не допускаються.

Приклад рядка файлу



\begin{verbatim}
13:algebra:surname:name:123456:77
\end{verbatim}


Нижче наведено код, який має залишити у файлі в кожному рядку тільки номер заліковки та предмет.

Рядок з прикладу має бути замінений на такий



\begin{verbatim}
123456:algebra
\end{verbatim}




Код:



\begin{verbatim}
regex = re.compile(r'XXX')
with open('1.txt', encoding='utf8') as f:
    lst = f.readlines()
for i in range(len(lst)):
    lst[i] = re.sub(regex, r'YYY', lst[i])
with open('1.txt', 'w', encoding='utf8') as f:
    f.writelines(lst)
\end{verbatim}









Який рядок треба записати на місце \kw{XXX} ?\par
\vspace{5mm}


Який рядок треба записати на місце \kw{YYY} ?\par


\newpage

{\begin {center}\textbf{ 45}\end {center}}\par
Файл \code{'1.txt'} містить відомість з рядками, в яких через кому записано поряд\-ковий номер (ціле), назву предмету (знаків пунктуації не містить), прізвище студента, ім'я студента, номер заліковки, оцінку в балах за 100 бальною шкалою.

Білі символи навколо роздільників не допускаються.

Приклад рядка файлу



\begin{verbatim}
13,algebra,surname,name,123456,77
\end{verbatim}


Нижче наведено код, який має залишити у файлі в кожному рядку тільки номер заліковки та предмет.

Рядок з прикладу має бути замінений на такий



\begin{verbatim}
123456,algebra
\end{verbatim}




Код



\begin{verbatim}
regex = re.compile(r'XXX')
with open('1.txt', encoding='utf8') as f:
    lst = f.readlines()
for i in range(len(lst)):
    lst[i] = re.sub(regex, r'YYY', lst[i])
with open('1.txt', 'w', encoding='utf8') as f:
    f.writelines(lst)
\end{verbatim}









Який рядок треба записати на місце \kw{XXX} ?\par
\vspace{5mm}


Який рядок треба записати на місце \kw{YYY} ?\par


\newpage

{\begin {center}\textbf{ 46}\end {center}}\par
Файл \code{'1.txt'} містить відомість з рядками, в яких через кому записано поряд\-ковий номер (ціле), назву предмету (знаків пунктуації не містить), прізвище студента, ім'я студента, номер заліковки, оцінку в балах за 100 бальною шкалою.

Білі символи навколо роздільників не допускаються.

Приклад рядка файлу



\begin{verbatim}
13,algebra,surname,name,123456,77
\end{verbatim}


Нижче наведено код, який шукає усі назви предметів, що зустрічаються у файлі, та зберігає їх у змінній \id{codes}.



\begin{verbatim}
codes = set()
regex = re.compile(r'XXX')
with open('1.txt', encoding='utf8') as f:
    for s in f:
        res = re.fullmatch(regex, s)
        if not res: continue
        s = ''
        codes.add(r'YYY')
\end{verbatim}





Який рядок треба записати на місце \kw{XXX} ?\par
\vspace{5mm}


Який рядок треба записати на місце \kw{YYY} ?\par


\newpage

{\begin {center}\textbf{ 47}\end {center}}\par
Файл \code{'1.txt'} містить журнал з рядками, в яких через двокрапку записано поряд\-ко\-вий номер запису (ціле), ім'я користувача (складається з латиниці), час події (фор\-мат як у прикладі), тип події (ERROR, WARNING, INFO, STATUS), код події (ціле).

Білі символи навколо роздільників не допускаються.

Приклад рядка журналу



\begin{verbatim}
12:username:2022-05-05 13-26:ERROR:404
\end{verbatim}


Нижче наведено код, який шукає усі попередження (подія WARNING), що зустрі\-ча\-ються в журналі і відбулися під час обідньої перерви деканату (з 13:00 по 13:59 включно), та зберігає їх коди в змінній \kw{codes}.



\begin{verbatim}
codes = set()
regex = re.compile(r'XXX')
with open('1.txt', encoding='utf8') as f:
    for s in f:
        res = re.fullmatch(regex, s)
        if not res: continue
        s = ''
        codes.add(r'YYY')
\end{verbatim}







Який рядок треба записати на місце \kw{XXX} ?\par
\vspace{5mm}


Який рядок треба записати на місце \kw{YYY} ?\par


\newpage

{\begin {center}\textbf{ 48}\end {center}}\par
Файл \code{'1.txt'} містить відомість з рядками, в яких через двокрапку записано поряд\-ковий номер (ціле), назву предмету (знаків пунктуації не містить), прізвище студента, ім'я студента, номер заліковки, оцінку в балах за 100 бальною шкалою.

Білі символи навколо роздільників не допускаються.

Приклад рядка файлу



\begin{verbatim}
13:algebra:surname:name:123456:77
\end{verbatim}


Нижче наведено код, який шукає усі назви предметів, що зустрічаються у файлі, та зберігає їх у змінній \id{codes}.



\begin{verbatim}
codes = set()
regex = re.compile(r'XXX')
with open('1.txt', encoding='utf8') as f:
    for s in f:
        res = re.fullmatch(regex, s)
        if not res: continue
        s = ''
        codes.add(r'YYY')
\end{verbatim}





Який рядок треба записати на місце \kw{XXX} ?\par
\vspace{5mm}


Який рядок треба записати на місце \kw{YYY} ?\par


\newpage

{\begin {center}\textbf{ 49}\end {center}}\par
Файл \code{'1.txt'} містить відомість з рядками, в яких через кому записа\-но назву предмету (знаків пунктуації не містить), прізвище студента, ім'я студента, номер залі\-ков\-ки, оцінку в балах за 100 бальною шкалою.

Білі символи навколо роздільників не допускаються.

Приклад рядка файлу



\begin{verbatim}
algebra,surname,name,123456,77
\end{verbatim}


Нижче наведено код, який має замінити у файлі усі номери заліковок на рядок \code{***}.



\begin{verbatim}
regex = re.compile(r'XXX')
with open('1.txt', encoding='utf8') as f:
    lst = f.readlines()
for i in range(len(lst)):
    lst[i] = re.sub(regex, r'YYY', lst[i])
with open('1.txt', 'w', encoding='utf8') as f:
    f.writelines(lst)
\end{verbatim}







Який рядок треба записати на місце \kw{XXX} ?\par
\vspace{5mm}


Який рядок треба записати на місце \kw{YYY} ?\par


\newpage

{\begin {center}\textbf{ 50}\end {center}}\par
Рядки журналу через кому містять тип події (ERROR, WARNING, INFO, STATUS), ім'я користувача (складається з латиниці), час події (фор\-мат як у прикладі), код події (ціле).

Білі символи навколо роздільників не допускаються.

Приклад такого рядка присвоєно змінній \kw{s}.



\begin{verbatim}
s = 'ERROR,username,2022-05-05 13:26,404'
\end{verbatim}


Нижче наведено код, за виконання якого значенням змінної \kw{out} має стати ім'я користувача.

Для наведеного прикладу значенням змінної \kw{out} має стати рядок \code{username}



\begin{verbatim}
res = re.fullmatch(r'XXX', s)
s = ''
out = r'YYY'
\end{verbatim}


Код має працювати не тільки для конкретного рядка з прикладу, але й для кожного рядка з журналу.




Який рядок треба записати на місце \kw{XXX} ?\par
\vspace{5mm}


Який рядок треба записати на місце \kw{YYY} ?\par


\newpage

{\begin {center}\textbf{ 51}\end {center}}\par
Файл \code{'1.txt'} містить журнал з рядками, в яких через двокрапку записано поряд\-ко\-вий номер запису в журналі (ціле), тип події (ERROR, WARNING, INFO, STATUS), ім'я користувача (складається з латиниці), час події (фор\-мат як у прикладі), код події (ціле).

Білі символи навколо роздільників не допускаються.

Приклад рядка файлу



\begin{verbatim}
13:ERROR:username:2022-05-05 13-26:404
\end{verbatim}


Нижче наведено код, який має замінити у файлі імена користувачів на рядок \code{unknown}.



\begin{verbatim}
regex = re.compile(r'XXX')
with open('1.txt', encoding='utf8') as f:
    lst = f.readlines()
for i in range(len(lst)):
    lst[i] = re.sub(regex, r'YYY', lst[i])
with open('1.txt', 'w', encoding='utf8') as f:
    f.writelines(lst)
\end{verbatim}







Який рядок треба записати на місце \kw{XXX} ?\par
\vspace{5mm}


Який рядок треба записати на місце \kw{YYY} ?\par


\newpage

{\begin {center}\textbf{ 52}\end {center}}\par
Рядки відомості через кому містять порядковий номер (ціле), назву предмету (зна\-ків пунктуації не містить), прізвище студента, ім'я студента, оцінку в балах за 100 бальною шкалою.

Білі символи навколо роздільників не допускаються.

Приклад такого рядка присвоєно змінній \kw{s}.



\begin{verbatim}
s = '13,algebra,surname,name,77'
\end{verbatim}


Нижче наведено код, за виконання якого значенням змінної \kw{out} має стати прізвище студента.

Для наведеного прикладу значенням змінної \kw{out} має стати рядок \code{surname}



\begin{verbatim}
res = re.fullmatch(r'XXX', s)
s = ''
out = r'YYY'
\end{verbatim}


Код має працювати не тільки для конкретного рядка з прикладу, але й для кожного рядка з відомості.




Який рядок треба записати на місце \kw{XXX} ?\par
\vspace{5mm}


Який рядок треба записати на місце \kw{YYY} ?\par


\newpage

{\begin {center}\textbf{ 53}\end {center}}\par
Файл \code{'1.txt'} містить відомість з рядками, в яких через кому записа\-но поряд\-ковий номер (ціле), назву предмету (знаків пунктуації не містить), прізвище студента, ім'я студента, номер заліковки, оцінку в балах за 100 бальною шкалою.

Білі символи навколо роздільників не допускаються.

Приклад рядка файлу



\begin{verbatim}
13,algebra,surname,name,123456,77
\end{verbatim}


Нижче наведено код, який має залишити у файлі в кожному рядку тільки номер заліковки, назву предмету та оцінку в балах за 100-бальною шкалою (у заданій послі\-дов\-ності).

Рядок з прикладу має бути замінений на такий



\begin{verbatim}
123456,algebra,77
\end{verbatim}




Код



\begin{verbatim}
regex = re.compile(r'XXX')
with open('1.txt', encoding='utf8') as f:
    lst = f.readlines()
for i in range(len(lst)):
    lst[i] = re.sub(regex, r'YYY', lst[i])
with open('1.txt', 'w', encoding='utf8') as f:
    f.writelines(lst)
\end{verbatim}







Який рядок треба записати на місце \kw{XXX} ?\par
\vspace{5mm}


Який рядок треба записати на місце \kw{YYY} ?\par


\newpage

{\begin {center}\textbf{ 54}\end {center}}\par
Файл \code{'1.txt'} містить відомість з рядками, в яких через двокрапку записа\-но поряд\-ковий номер (ціле), назву предмету (знаків пунктуації не містить), прізвище студента, ім'я студента, номер заліковки, оцінку в балах за 100 бальною шкалою.

Білі символи навколо роздільників не допускаються.

Приклад рядка файлу



\begin{verbatim}
13:algebra:surname:name:123456:77
\end{verbatim}


Нижче наведено код, який має залишити у файлі в кожному рядку тільки номер заліковки, назву предмету та оцінку в балах за 100-бальною шкалою (у заданій послі\-дов\-ності).

Рядок з прикладу має бути замінений на такий



\begin{verbatim}
123456:algebra:77
\end{verbatim}




Код:



\begin{verbatim}
regex = re.compile(r'XXX')
with open('1.txt', encoding='utf8') as f:
    lst = f.readlines()
for i in range(len(lst)):
    lst[i] = re.sub(regex, r'YYY', lst[i])
with open('1.txt', 'w', encoding='utf8') as f:
    f.writelines(lst)
\end{verbatim}







Який рядок треба записати на місце \kw{XXX} ?\par
\vspace{5mm}


Який рядок треба записати на місце \kw{YYY} ?\par


\newpage

{\begin {center}\textbf{ 55}\end {center}}\par
Файл \code{'1.txt'} містить журнал з рядками, в яких через кому записано поряд\-ко\-вий номер запису в журналі (ціле), код події (ціле), тип події (ERROR,\\ WARNING, INFO, STATUS), ім'я користувача (складається з латиниці), час події (фор\-мат як у прикладі), опис події.

Білі символи навколо роздільників не допускаються.

Приклад рядка файлу



\begin{verbatim}
13,404,ERROR,username,2022-05-05 13:26,something happens
\end{verbatim}


Нижче наведено код, який має в кожному рядку файли переставити місцями код та тип події.

Рядок з прикладу має бути замінений на такий



\begin{verbatim}
13,ERROR,404,username,2022-05-05 13:26,something happens
\end{verbatim}




Код:



\begin{verbatim}
regex = re.compile(r'XXX')
with open('1.txt', encoding='utf8') as f:
    lst = f.readlines()
for i in range(len(lst)):
    lst[i] = re.sub(regex, r'YYY', lst[i])
with open('1.txt', 'w', encoding='utf8') as f:
    f.writelines(lst)
\end{verbatim}







Який рядок треба записати на місце \kw{XXX} ?\par
\vspace{5mm}


Який рядок треба записати на місце \kw{YYY} ?\par


\newpage

{\begin {center}\textbf{ 56}\end {center}}\par
Файл \code{'1.txt'} містить журнал з рядками, в яких через двокрапку записано поряд\-ко\-вий номер запису в журналі (ціле), код події (ціле), тип події (ERROR,\\ WARNING, INFO, STATUS), ім'я користувача (складається з латиниці), час події (фор\-мат як у прикладі), опис події.

Білі символи навколо роздільників не допускаються.

Приклад рядка файлу



\begin{verbatim}
13:404:ERROR:username:2022-05-05 13:26:something happens
\end{verbatim}


Нижче наведено код, який має в кожному рядку файли переставити місцями код та тип події.

Рядок з прикладу має бути замінений на такий



\begin{verbatim}
13:ERROR:404:username:2022-05-05 13:26:something happens
\end{verbatim}




Код:



\begin{verbatim}
regex = re.compile(r'XXX')
with open('1.txt', encoding='utf8') as f:
    lst = f.readlines()
for i in range(len(lst)):
    lst[i] = re.sub(regex, r'YYY', lst[i])
with open('1.txt', 'w', encoding='utf8') as f:
    f.writelines(lst)
\end{verbatim}







Який рядок треба записати на місце \kw{XXX} ?\par
\vspace{5mm}


Який рядок треба записати на місце \kw{YYY} ?\par


\newpage

{\begin {center}\textbf{ 57}\end {center}}\par
Файл \code{'1.txt'} містить відомість з рядками, в яких через кому записано поряд\-ковий номер (ціле), назву предмету (знаків пунктуації не містить), прізвище студента, ім'я студента, номер заліковки, оцінку в балах за 100 бальною шкалою.

Білі символи навколо роздільників не допускаються.

Приклад рядка файлу



\begin{verbatim}
13,algebra,surname,name,123456,77
\end{verbatim}


Нижче наведено код, який шукає усі предмети, який складав студент із заліковкою \code{777}, та зберігає їх типи в змінній \kw{codes}.



\begin{verbatim}
codes = set()
regex = re.compile(r'XXX')
with open('1.txt', encoding='utf8') as f:
    for s in f:
        res = re.fullmatch(regex, s)
        if not res: continue
        s = ''
        codes.add(r'YYY')
\end{verbatim}





Який рядок треба записати на місце \kw{XXX} ?\par
\vspace{5mm}


Який рядок треба записати на місце \kw{YYY} ?\par


\newpage

{\begin {center}\textbf{ 58}\end {center}}\par
Файл \code{'1.txt'} містить відомість з рядками, в яких через двокрапку записано поряд\-ковий номер (ціле), назву предмету (знаків пунктуації не містить), прізвище студента, ім'я студента, номер заліковки, оцінку в балах за 100 бальною шкалою.

Білі символи навколо роздільників не допускаються.

Приклад рядка файлу



\begin{verbatim}
13:algebra:surname:name:123456:77
\end{verbatim}


Нижче наведено код, який шукає усі предмети, який складав студент із заліковкою \code{777}, та зберігає їх типи в змінній \kw{codes}.



\begin{verbatim}
codes = set()
regex = re.compile(r'XXX')
with open('1.txt', encoding='utf8') as f:
    for s in f:
        res = re.fullmatch(regex, s)
        if not res: continue
        s = ''
        codes.add(r'YYY')
\end{verbatim}





Який рядок треба записати на місце \kw{XXX} ?\par
\vspace{5mm}


Який рядок треба записати на місце \kw{YYY} ?\par


\newpage

{\begin {center}\textbf{ 59}\end {center}}\par
Рядки журналу через кому містять ім'я користувача (складається з латиниці), час події (фор\-мат як у прикладі), тип події (ERROR, WARNING, INFO, STATUS), код події (ціле).

Білі символи навколо роздільників не допускаються.

Приклад такого рядка присвоєно змінній \kw{s}.



\begin{verbatim}
s = 'username,2022-05-05 13:26,ERROR,404'
\end{verbatim}


Нижче наведено код, за виконання якого значенням змінної \kw{out} має стати код події.

Для наведеного прикладу значенням змінної \kw{out} має стати рядок \code{404}



\begin{verbatim}
res = re.fullmatch(r'XXX', s)
s = ''
out = r'YYY'
\end{verbatim}


Код має працювати не тільки для конкретного рядка з прикладу, але й для кожного рядка з журналу.




Який рядок треба записати на місце \kw{XXX} ?\par
\vspace{5mm}


Який рядок треба записати на місце \kw{YYY} ?\par


\newpage

{\begin {center}\textbf{ 60}\end {center}}\par
Файл \code{'1.txt'} містить журнал з рядками, в яких через двокрапку записано поряд\-ко\-вий номер запису (ціле), ім'я користувача (складається з латиниці), час події (фор\-мат як у прикладі), тип події (ERROR, WARNING, INFO, STATUS), код події (ціле).

Білі символи навколо роздільників не допускаються.

Приклад рядка журналу



\begin{verbatim}
12:username:2022-05-05 13-26:ERROR:404
\end{verbatim}


Нижче наведено код, який шукає усі події, що відбулись\\ для користувача \code{anonymous}, та зберігає їх типи в змінній \kw{codes}.



\begin{verbatim}
codes = set()
regex = re.compile(r'XXX')
with open('1.txt', encoding='utf8') as f:
    for s in f:
        res = re.fullmatch(regex, s)
        if not res: continue
        s = ''
        codes.add(r'YYY')
\end{verbatim}





Який рядок треба записати на місце \kw{XXX} ?\par
\vspace{5mm}


Який рядок треба записати на місце \kw{YYY} ?\par


\newpage

{\begin {center}\textbf{ 61}\end {center}}\par
Файл \code{'1.txt'} містить журнал з рядками, в яких через кому записано поряд\-ко\-вий номер запису в журналі (ціле), тип події (ERROR, WARNING, INFO, STATUS), ім'я користувача (складається з латиниці), час події (фор\-мат як у прикладі), код події (ціле).

Білі символи навколо роздільників не допускаються.

Приклад рядка файлу



\begin{verbatim}
13,ERROR,username,2022-05-05 13:26,404
\end{verbatim}


Нижче наведено код, який має замінити у файлі імена користувачів на рядок \code{unknown}.



\begin{verbatim}
regex = re.compile(r'XXX')
with open('1.txt', encoding='utf8') as f:
    lst = f.readlines()
for i in range(len(lst)):
    lst[i] = re.sub(regex, r'YYY', lst[i])
with open('1.txt', 'w', encoding='utf8') as f:
    f.writelines(lst)
\end{verbatim}







Який рядок треба записати на місце \kw{XXX} ?\par
\vspace{5mm}


Який рядок треба записати на місце \kw{YYY} ?\par


\newpage

{\begin {center}\textbf{ 62}\end {center}}\par
Рядки журналу через двокрапку містять ім'я користувача (складається з латиниці), час події (фор\-мат як у прикладі), тип події (ERROR, WARNING, INFO, STATUS), код події (ціле).

Білі символи навколо роздільників не допускаються.

Приклад такого рядка присвоєно змінній \kw{s}.



\begin{verbatim}
s = 'username:2022-05-05 13-26:ERROR:404'
\end{verbatim}


Нижче наведено код, за виконання якого значенням змінної \kw{out} має стати код події.

Для наведеного прикладу значенням змінної \kw{out} має стати рядок \code{404}



\begin{verbatim}
res = re.fullmatch(r'XXX', s)
s = ''
out = r'YYY'
\end{verbatim}


Код має працювати не тільки для конкретного рядка з прикладу, але й для кожного рядка з журналу.




Який рядок треба записати на місце \kw{XXX} ?\par
\vspace{5mm}


Який рядок треба записати на місце \kw{YYY} ?\par


\newpage

{\begin {center}\textbf{ 63}\end {center}}\par
Рядки відомості через кому містять порядковий номер (ціле), пріз\-ви\-ще студен\-та, ім'я студента, номер заліковки, оцінку в балах за 100 бальною шкалою.

Білі символи навколо роздільників не допускаються.

Приклад такого рядка присвоєно змінній \kw{s}.



\begin{verbatim}
s = '13,surname,name,123456,77'
\end{verbatim}


Нижче наведено код, за виконання якого для рядка \kw{s} значенням змінної \kw{out} має стати номер заліковки.

Для наведеного прикладу значенням змінної \kw{out} має стати рядок \code{123456} .



\begin{verbatim}
res = re.fullmatch(r'XXX', s)
s = ''
out = r'YYY'
\end{verbatim}


Код має працювати не тільки для конкретного рядка з прикладу, але й для кожного рядка з відомості.




Який рядок треба записати на місце \kw{XXX} ?\par
\vspace{5mm}


Який рядок треба записати на місце \kw{YYY} ?\par


\newpage

{\begin {center}\textbf{ 64}\end {center}}\par
Файл \code{'1.txt'} містить відомість з рядками, в яких через двокрапку записано поряд\-ковий номер (ціле), назву предмету (знаків пунктуації не містить), прізвище студента, ім'я студента, номер заліковки, оцінку в балах за 100 бальною шкалою.

Білі символи навколо роздільників не допускаються.

Приклад рядка файлу



\begin{verbatim}
13:algebra:surname:name:123456:77
\end{verbatim}


Нижче наведено код, який має в кожному рядку файлу переставити місцями назву предмету та номер заліковки.

Рядок з прикладу має бути замінений на такий



\begin{verbatim}
13:123456:surname:name:algebra:77
\end{verbatim}




Код:



\begin{verbatim}
regex = re.compile(r'XXX')
with open('1.txt', encoding='utf8') as f:
    lst = f.readlines()
for i in range(len(lst)):
    lst[i] = re.sub(regex, r'YYY', lst[i])
with open('1.txt', 'w', encoding='utf8') as f:
    f.writelines(lst)
\end{verbatim}







Який рядок треба записати на місце \kw{XXX} ?\par
\vspace{5mm}


Який рядок треба записати на місце \kw{YYY} ?\par


\newpage

{\begin {center}\textbf{ 65}\end {center}}\par
Файл \code{'1.txt'} містить журнал з рядками, в яких через кому записано поряд\-ко\-вий номер запису (ціле), ім'я користувача (складається з латиниці), час події (фор\-мат як у прикладі), тип події (ERROR, WARNING, INFO, STATUS), код події (ціле).

Білі символи навколо роздільників не допускаються.

Приклад рядка журналу



\begin{verbatim}
12,username,2022-05-05 13:26,ERROR,404
\end{verbatim}


Нижче наведено код, який шукає усі попередження (подія WARNING), що зустрі\-ча\-ються в журналі і відбулися під час обідньої перерви деканату (з 13,00 по 13,59 включно), та зберігає їх коди в змінній \kw{codes}.



\begin{verbatim}
codes = set()
regex = re.compile(r'XXX')
with open('1.txt', encoding='utf8') as f:
    for s in f:
        res = re.fullmatch(regex, s)
        if not res: continue
        s = ''
        codes.add(r'YYY')
\end{verbatim}







Який рядок треба записати на місце \kw{XXX} ?\par
\vspace{5mm}


Який рядок треба записати на місце \kw{YYY} ?\par


\newpage

{\begin {center}\textbf{ 66}\end {center}}\par
Файл \code{'1.txt'} містить журнал з рядками, в яких через двокрапку записано поряд\-ко\-вий номер запису в журналі (ціле), код події (ціле), тип події (ERROR, \\WARNING, INFO, STATUS), ім'я користувача (складається з латиниці), час події (фор\-мат як у прикладі), опис події.

Білі символи навколо роздільників не допускаються.

Приклад рядка файлу



\begin{verbatim}
13:404:ERROR:username:2022-05-05 13-26:something happens
\end{verbatim}


Нижче наведено код, який має залишити у файлі в кожному рядку тільки код події та ім'я користувача. Рядок з прикладу має бути замінений на такий



\begin{verbatim}
404:username
\end{verbatim}




Код:



\begin{verbatim}
regex = re.compile(r'XXX')
with open('1.txt', encoding='utf8') as f:
    lst = f.readlines()
for i in range(len(lst)):
    lst[i] = re.sub(regex, r'YYY', lst[i])
with open('1.txt', 'w', encoding='utf8') as f:
    f.writelines(lst)
\end{verbatim}







Який рядок треба записати на місце \kw{XXX} ?\par
\vspace{5mm}


Який рядок треба записати на місце \kw{YYY} ?\par


\newpage

{\begin {center}\textbf{ 67}\end {center}}\par
Файл \code{'1.txt'} містить журнал з рядками, в яких через кому записано поряд\-ко\-вий номер запису в журналі (ціле), тип події (ERROR, WARNING, INFO, STATUS), ім'я користувача (складається з латиниці), час події (фор\-мат як у прикладі), код події (ціле).

Білі символи навколо роздільників не допускаються.

Приклад рядка файлу



\begin{verbatim}
13,ERROR,username,2022-05-05 13:26,404
\end{verbatim}


Нижче наведено код, який шукає усі коди помилок (подія ERROR), що зустріча\-ють\-ся в журналі, та зберігає їх у змінній \kw{codes}.



\begin{verbatim}
codes = set()
regex = re.compile(r'XXX')
with open('1.txt', encoding='utf8') as f:
    for s in f:
        res = re.fullmatch(regex, s)
        if not res: continue
        s = ''
        codes.add(r'YYY')
\end{verbatim}





Який рядок треба записати на місце \kw{XXX} ?\par
\vspace{5mm}


Який рядок треба записати на місце \kw{YYY} ?\par


\newpage

{\begin {center}\textbf{ 68}\end {center}}\par
Рядки відомості через двокрапку містять порядковий номер (ціле), назву предмету (знаків пунктуації не містить), прізвище студента, ім'я студента, оцінку в балах за 100 бальною шкалою.

Білі символи навколо роздільників не допускаються.

Приклад такого рядка присвоєно змінній \kw{s}.



\begin{verbatim}
s = '13:algebra:surname:name:77'
\end{verbatim}


Нижче наведено код, за виконання якого значенням змінної \kw{out} має стати прізвище студента.

Для наведеного прикладу значенням змінної \kw{out} має стати рядок \code{surname}



\begin{verbatim}
res = re.fullmatch(r'XXX', s)
s = ''
out = r'YYY'
\end{verbatim}


Код має працювати не тільки для конкретного рядка з прикладу, але й для кожного рядка з відомості.




Який рядок треба записати на місце \kw{XXX} ?\par
\vspace{5mm}


Який рядок треба записати на місце \kw{YYY} ?\par


\newpage

{\begin {center}\textbf{ 69}\end {center}}\par
Рядки журналу через кому містять тип події (ERROR, WARNING, INFO, STATUS), ім'я користувача (складається з латиниці), час події (фор\-мат як у прикладі), код події (ціле).

Білі символи навколо роздільників не допускаються.

Приклад такого рядка присвоєно змінній \kw{s}.



\begin{verbatim}
s = 'ERROR,username,2022-05-05 13:26,404'
\end{verbatim}


Нижче наведено код, за виконання якого значенням змінної \kw{out} має стати ім'я користувача.

Для наведеного прикладу значенням змінної \kw{out} має стати рядок \code{username}



\begin{verbatim}
res = re.fullmatch(r'XXX', s)
s = ''
out = r'YYY'
\end{verbatim}


Код має працювати не тільки для конкретного рядка з прикладу, але й для кожного рядка з журналу.




Який рядок треба записати на місце \kw{XXX} ?\par
\vspace{5mm}


Який рядок треба записати на місце \kw{YYY} ?\par


\newpage

{\begin {center}\textbf{ 70}\end {center}}\par
Файл \code{'1.txt'} містить відомість з рядками, в яких через кому записано поряд\-ковий номер (ціле), назву предмету (знаків пунктуації не містить), прізвище студента, ім'я студента, номер заліковки, оцінку в балах за 100 бальною шкалою.

Білі символи навколо роздільників не допускаються.

Приклад рядка файлу



\begin{verbatim}
13,algebra,surname,name,123456,77
\end{verbatim}


Нижче наведено код, який шукає усі назви предметів, що зустрічаються у файлі, та зберігає їх у змінній \id{codes}.



\begin{verbatim}
codes = set()
regex = re.compile(r'XXX')
with open('1.txt', encoding='utf8') as f:
    for s in f:
        res = re.fullmatch(regex, s)
        if not res: continue
        s = ''
        codes.add(r'YYY')
\end{verbatim}





Який рядок треба записати на місце \kw{XXX} ?\par
\vspace{5mm}


Який рядок треба записати на місце \kw{YYY} ?\par


\newpage

{\begin {center}\textbf{ 71}\end {center}}\par
Файл \code{'1.txt'} містить відомість з рядками, в яких через двокрапку записано поряд\-ковий номер (ціле), назву предмету (знаків пунктуації не містить), прізвище студента, ім'я студента, номер заліковки, оцінку в балах за 100 бальною шкалою.

Білі символи навколо роздільників не допускаються.

Приклад рядка файлу



\begin{verbatim}
13:algebra:surname:name:123456:77
\end{verbatim}


Нижче наведено код, який має залишити у файлі в кожному рядку тільки номер заліковки та предмет.

Рядок з прикладу має бути замінений на такий



\begin{verbatim}
123456:algebra
\end{verbatim}




Код:



\begin{verbatim}
regex = re.compile(r'XXX')
with open('1.txt', encoding='utf8') as f:
    lst = f.readlines()
for i in range(len(lst)):
    lst[i] = re.sub(regex, r'YYY', lst[i])
with open('1.txt', 'w', encoding='utf8') as f:
    f.writelines(lst)
\end{verbatim}









Який рядок треба записати на місце \kw{XXX} ?\par
\vspace{5mm}


Який рядок треба записати на місце \kw{YYY} ?\par


\newpage

{\begin {center}\textbf{ 72}\end {center}}\par
Файл \code{'1.txt'} містить відомість з рядками, в яких через двокрапку записа\-но назву предмету (знаків пунктуації не містить), прізвище студента, ім'я студента, номер залі\-ков\-ки, оцінку в балах за 100 бальною шкалою.

Білі символи навколо роздільників не допускаються.

Приклад рядка файлу



\begin{verbatim}
algebra:surname:name:123456:77
\end{verbatim}


Нижче наведено код, який має замінити у файлі усі номери заліковок на рядок \code{***}.



\begin{verbatim}
regex = re.compile(r'XXX')
with open('1.txt', encoding='utf8') as f:
    lst = f.readlines()
for i in range(len(lst)):
    lst[i] = re.sub(regex, r'YYY', lst[i])
with open('1.txt', 'w', encoding='utf8') as f:
    f.writelines(lst)
\end{verbatim}







Який рядок треба записати на місце \kw{XXX} ?\par
\vspace{5mm}


Який рядок треба записати на місце \kw{YYY} ?\par


\newpage

{\begin {center}\textbf{ 73}\end {center}}\par
Файл \code{'1.txt'} містить журнал з рядками, в яких через кому записано поряд\-ко\-вий номер запису в журналі (ціле), тип події (ERROR, WARNING, INFO, STATUS), ім'я користувача (складається з латиниці), час події (фор\-мат як у прикладі), код події (ціле).

Білі символи навколо роздільників не допускаються.

Приклад рядка файлу



\begin{verbatim}
13,ERROR,username,2022-05-05 13:26,404
\end{verbatim}


Нижче наведено код, який шукає усі коди помилок (подія ERROR), що зустріча\-ють\-ся в журналі, та зберігає їх у змінній \kw{codes}.



\begin{verbatim}
codes = set()
regex = re.compile(r'XXX')
with open('1.txt', encoding='utf8') as f:
    for s in f:
        res = re.fullmatch(regex, s)
        if not res: continue
        s = ''
        codes.add(r'YYY')
\end{verbatim}





Який рядок треба записати на місце \kw{XXX} ?\par
\vspace{5mm}


Який рядок треба записати на місце \kw{YYY} ?\par


\newpage

{\begin {center}\textbf{ 74}\end {center}}\par
Файл \code{'1.txt'} містить відомість з рядками, в яких через двокрапку записа\-но поряд\-ковий номер (ціле), назву предмету (знаків пунктуації не містить), прізвище студента, ім'я студента, номер заліковки, оцінку в балах за 100 бальною шкалою.

Білі символи навколо роздільників не допускаються.

Приклад рядка файлу



\begin{verbatim}
13:algebra:surname:name:123456:77
\end{verbatim}


Нижче наведено код, який має залишити у файлі в кожному рядку тільки номер заліковки, назву предмету та оцінку в балах за 100-бальною шкалою (у заданій послі\-дов\-ності).

Рядок з прикладу має бути замінений на такий



\begin{verbatim}
123456:algebra:77
\end{verbatim}




Код:



\begin{verbatim}
regex = re.compile(r'XXX')
with open('1.txt', encoding='utf8') as f:
    lst = f.readlines()
for i in range(len(lst)):
    lst[i] = re.sub(regex, r'YYY', lst[i])
with open('1.txt', 'w', encoding='utf8') as f:
    f.writelines(lst)
\end{verbatim}







Який рядок треба записати на місце \kw{XXX} ?\par
\vspace{5mm}


Який рядок треба записати на місце \kw{YYY} ?\par


\newpage

{\begin {center}\textbf{ 75}\end {center}}\par
Файл \code{'1.txt'} містить журнал з рядками, в яких через кому записано поряд\-ко\-вий номер запису в журналі (ціле), тип події (ERROR, WARNING, INFO, STATUS), ім'я користувача (складається з латиниці), час події (фор\-мат як у прикладі), код події (ціле).

Білі символи навколо роздільників не допускаються.

Приклад рядка файлу



\begin{verbatim}
13,ERROR,username,2022-05-05 13:26,404
\end{verbatim}


Нижче наведено код, який має замінити у файлі імена користувачів на рядок \code{unknown}.



\begin{verbatim}
regex = re.compile(r'XXX')
with open('1.txt', encoding='utf8') as f:
    lst = f.readlines()
for i in range(len(lst)):
    lst[i] = re.sub(regex, r'YYY', lst[i])
with open('1.txt', 'w', encoding='utf8') as f:
    f.writelines(lst)
\end{verbatim}







Який рядок треба записати на місце \kw{XXX} ?\par
\vspace{5mm}


Який рядок треба записати на місце \kw{YYY} ?\par


\newpage

{\begin {center}\textbf{ 76}\end {center}}\par
Файл \code{'1.txt'} містить журнал з рядками, в яких через кому записано поряд\-ко\-вий номер запису в журналі (ціле), код події (ціле), тип події (ERROR, \\WARNING, INFO, STATUS), ім'я користувача (складається з латиниці), час події (фор\-мат як у прикладі), опис події.

Білі символи навколо роздільників не допускаються.

Приклад рядка файлу



\begin{verbatim}
13,404,ERROR,username,2022-05-05 13:26,something happens
\end{verbatim}


Нижче наведено код, який має залишити у файлі в кожному рядку тільки код події та ім'я користувача. Рядок з прикладу має бути замінений на такий



\begin{verbatim}
404,username
\end{verbatim}




Код:



\begin{verbatim}
regex = re.compile(r'XXX')
with open('1.txt', encoding='utf8') as f:
    lst = f.readlines()
for i in range(len(lst)):
    lst[i] = re.sub(regex, r'YYY', lst[i])
with open('1.txt', 'w', encoding='utf8') as f:
    f.writelines(lst)
\end{verbatim}







Який рядок треба записати на місце \kw{XXX} ?\par
\vspace{5mm}


Який рядок треба записати на місце \kw{YYY} ?\par


\newpage

{\begin {center}\textbf{ 77}\end {center}}\par
Файл \code{'1.txt'} містить журнал з рядками, в яких через кому записано поряд\-ко\-вий номер запису (ціле), ім'я користувача (складається з латиниці), час події (фор\-мат як у прикладі), тип події (ERROR, WARNING, INFO, STATUS), код події (ціле).

Білі символи навколо роздільників не допускаються.

Приклад рядка журналу



\begin{verbatim}
12,username,2022-05-05 13:26,ERROR,404
\end{verbatim}


Нижче наведено код, який шукає усі події, що відбулись\\ для користувача \code{anonymous}, та зберігає їх типи в змінній \kw{codes}.



\begin{verbatim}
codes = set()
regex = re.compile(r'XXX')
with open('1.txt', encoding='utf8') as f:
    for s in f:
        res = re.fullmatch(regex, s)
        if not res: continue
        s = ''
        codes.add(r'YYY')
\end{verbatim}





Який рядок треба записати на місце \kw{XXX} ?\par
\vspace{5mm}


Який рядок треба записати на місце \kw{YYY} ?\par


\newpage

{\begin {center}\textbf{ 78}\end {center}}\par
Файл \code{'1.txt'} містить відомість з рядками, в яких через кому записа\-но назву предмету (знаків пунктуації не містить), прізвище студента, ім'я студента, номер залі\-ков\-ки, оцінку в балах за 100 бальною шкалою.

Білі символи навколо роздільників не допускаються.

Приклад рядка файлу



\begin{verbatim}
algebra,surname,name,123456,77
\end{verbatim}


Нижче наведено код, який має замінити у файлі усі номери заліковок на рядок \code{***}.



\begin{verbatim}
regex = re.compile(r'XXX')
with open('1.txt', encoding='utf8') as f:
    lst = f.readlines()
for i in range(len(lst)):
    lst[i] = re.sub(regex, r'YYY', lst[i])
with open('1.txt', 'w', encoding='utf8') as f:
    f.writelines(lst)
\end{verbatim}







Який рядок треба записати на місце \kw{XXX} ?\par
\vspace{5mm}


Який рядок треба записати на місце \kw{YYY} ?\par


\newpage

{\begin {center}\textbf{ 79}\end {center}}\par
Файл \code{'1.txt'} містить журнал з рядками, в яких через двокрапку записано поряд\-ко\-вий номер запису (ціле), ім'я користувача (складається з латиниці), час події (фор\-мат як у прикладі), тип події (ERROR, WARNING, INFO, STATUS), код події (ціле).

Білі символи навколо роздільників не допускаються.

Приклад рядка журналу



\begin{verbatim}
12:username:2022-05-05 13-26:ERROR:404
\end{verbatim}


Нижче наведено код, який шукає усі події, що відбулись\\ для користувача \code{anonymous}, та зберігає їх типи в змінній \kw{codes}.



\begin{verbatim}
codes = set()
regex = re.compile(r'XXX')
with open('1.txt', encoding='utf8') as f:
    for s in f:
        res = re.fullmatch(regex, s)
        if not res: continue
        s = ''
        codes.add(r'YYY')
\end{verbatim}





Який рядок треба записати на місце \kw{XXX} ?\par
\vspace{5mm}


Який рядок треба записати на місце \kw{YYY} ?\par


\newpage

{\begin {center}\textbf{ 80}\end {center}}\par
Файл \code{'1.txt'} містить журнал з рядками, в яких через двокрапку записано поряд\-ко\-вий номер запису (ціле), ім'я користувача (складається з латиниці), час події (фор\-мат як у прикладі), тип події (ERROR, WARNING, INFO, STATUS), код події (ціле).

Білі символи навколо роздільників не допускаються.

Приклад рядка журналу



\begin{verbatim}
12:username:2022-05-05 13-26:ERROR:404
\end{verbatim}


Нижче наведено код, який шукає усі попередження (подія WARNING), що зустрі\-ча\-ються в журналі і відбулися під час обідньої перерви деканату (з 13:00 по 13:59 включно), та зберігає їх коди в змінній \kw{codes}.



\begin{verbatim}
codes = set()
regex = re.compile(r'XXX')
with open('1.txt', encoding='utf8') as f:
    for s in f:
        res = re.fullmatch(regex, s)
        if not res: continue
        s = ''
        codes.add(r'YYY')
\end{verbatim}







Який рядок треба записати на місце \kw{XXX} ?\par
\vspace{5mm}


Який рядок треба записати на місце \kw{YYY} ?\par


\newpage

{\begin {center}\textbf{ 81}\end {center}}\par
Файл \code{'1.txt'} містить відомість з рядками, в яких через двокрапку записано поряд\-ковий номер (ціле), назву предмету (знаків пунктуації не містить), прізвище студента, ім'я студента, номер заліковки, оцінку в балах за 100 бальною шкалою.

Білі символи навколо роздільників не допускаються.

Приклад рядка файлу



\begin{verbatim}
13:algebra:surname:name:123456:77
\end{verbatim}


Нижче наведено код, який має в кожному рядку файлу переставити місцями назву предмету та номер заліковки.

Рядок з прикладу має бути замінений на такий



\begin{verbatim}
13:123456:surname:name:algebra:77
\end{verbatim}




Код:



\begin{verbatim}
regex = re.compile(r'XXX')
with open('1.txt', encoding='utf8') as f:
    lst = f.readlines()
for i in range(len(lst)):
    lst[i] = re.sub(regex, r'YYY', lst[i])
with open('1.txt', 'w', encoding='utf8') as f:
    f.writelines(lst)
\end{verbatim}







Який рядок треба записати на місце \kw{XXX} ?\par
\vspace{5mm}


Який рядок треба записати на місце \kw{YYY} ?\par


\newpage

{\begin {center}\textbf{ 82}\end {center}}\par
Файл \code{'1.txt'} містить журнал з рядками, в яких через двокрапку записано поряд\-ко\-вий номер запису в журналі (ціле), тип події (ERROR, WARNING, INFO, STATUS), ім'я користувача (складається з латиниці), час події (фор\-мат як у прикладі), опис події.

Білі символи навколо роздільників не допускаються.

Приклад рядка файлу



\begin{verbatim}
13:ERROR:username:2022-05-05 13-26:something happens
\end{verbatim}


Нижче наведено код, який має залишити у файлі в кожному рядку тільки тип події та її опис.

Рядок з прикладу має бути замінений на такий



\begin{verbatim}
ERROR:something happens
\end{verbatim}




Код:



\begin{verbatim}
regex = re.compile(r'XXX')
with open('1.txt', encoding='utf8') as f:
    lst = f.readlines()
for i in range(len(lst)):
    lst[i] = re.sub(regex, r'YYY', lst[i])
with open('1.txt', 'w', encoding='utf8') as f:
    f.writelines(lst)
\end{verbatim}









Який рядок треба записати на місце \kw{XXX} ?\par
\vspace{5mm}


Який рядок треба записати на місце \kw{YYY} ?\par


\newpage

{\begin {center}\textbf{ 83}\end {center}}\par
Рядки відомості через кому містять порядковий номер (ціле), пріз\-ви\-ще студен\-та, ім'я студента, номер заліковки, оцінку в балах за 100 бальною шкалою.

Білі символи навколо роздільників не допускаються.

Приклад такого рядка присвоєно змінній \kw{s}.



\begin{verbatim}
s = '13,surname,name,123456,77'
\end{verbatim}


Нижче наведено код, за виконання якого для рядка \kw{s} значенням змінної \kw{out} має стати номер заліковки.

Для наведеного прикладу значенням змінної \kw{out} має стати рядок \code{123456} .



\begin{verbatim}
res = re.fullmatch(r'XXX', s)
s = ''
out = r'YYY'
\end{verbatim}


Код має працювати не тільки для конкретного рядка з прикладу, але й для кожного рядка з відомості.




Який рядок треба записати на місце \kw{XXX} ?\par
\vspace{5mm}


Який рядок треба записати на місце \kw{YYY} ?\par


\newpage

{\begin {center}\textbf{ 84}\end {center}}\par
Файл \code{'1.txt'} містить журнал з рядками, в яких через двокрапку записано поряд\-ко\-вий номер запису в журналі (ціле), код події (ціле), тип події (ERROR,\\ WARNING, INFO, STATUS), ім'я користувача (складається з латиниці), час події (фор\-мат як у прикладі), опис події.

Білі символи навколо роздільників не допускаються.

Приклад рядка файлу



\begin{verbatim}
13:404:ERROR:username:2022-05-05 13:26:something happens
\end{verbatim}


Нижче наведено код, який має в кожному рядку файли переставити місцями код та тип події.

Рядок з прикладу має бути замінений на такий



\begin{verbatim}
13:ERROR:404:username:2022-05-05 13:26:something happens
\end{verbatim}




Код:



\begin{verbatim}
regex = re.compile(r'XXX')
with open('1.txt', encoding='utf8') as f:
    lst = f.readlines()
for i in range(len(lst)):
    lst[i] = re.sub(regex, r'YYY', lst[i])
with open('1.txt', 'w', encoding='utf8') as f:
    f.writelines(lst)
\end{verbatim}







Який рядок треба записати на місце \kw{XXX} ?\par
\vspace{5mm}


Який рядок треба записати на місце \kw{YYY} ?\par


\newpage

{\begin {center}\textbf{ 85}\end {center}}\par
Файл \code{'1.txt'} містить журнал з рядками, в яких через кому записано поряд\-ко\-вий номер запису в журналі (ціле), код події (ціле), тип події (ERROR,\\ WARNING, INFO, STATUS), ім'я користувача (складається з латиниці), час події (фор\-мат як у прикладі), опис події.

Білі символи навколо роздільників не допускаються.

Приклад рядка файлу



\begin{verbatim}
13,404,ERROR,username,2022-05-05 13:26,something happens
\end{verbatim}


Нижче наведено код, який має в кожному рядку файли переставити місцями код та тип події.

Рядок з прикладу має бути замінений на такий



\begin{verbatim}
13,ERROR,404,username,2022-05-05 13:26,something happens
\end{verbatim}




Код:



\begin{verbatim}
regex = re.compile(r'XXX')
with open('1.txt', encoding='utf8') as f:
    lst = f.readlines()
for i in range(len(lst)):
    lst[i] = re.sub(regex, r'YYY', lst[i])
with open('1.txt', 'w', encoding='utf8') as f:
    f.writelines(lst)
\end{verbatim}







Який рядок треба записати на місце \kw{XXX} ?\par
\vspace{5mm}


Який рядок треба записати на місце \kw{YYY} ?\par


\newpage

{\begin {center}\textbf{ 86}\end {center}}\par
Рядки відомості через кому містять порядковий номер (ціле), назву предмету (зна\-ків пунктуації не містить), прізвище студента, ім'я студента, оцінку в балах за 100 бальною шкалою.

Білі символи навколо роздільників не допускаються.

Приклад такого рядка присвоєно змінній \kw{s}.



\begin{verbatim}
s = '13,algebra,surname,name,77'
\end{verbatim}


Нижче наведено код, за виконання якого значенням змінної \kw{out} має стати прізвище студента.

Для наведеного прикладу значенням змінної \kw{out} має стати рядок \code{surname}



\begin{verbatim}
res = re.fullmatch(r'XXX', s)
s = ''
out = r'YYY'
\end{verbatim}


Код має працювати не тільки для конкретного рядка з прикладу, але й для кожного рядка з відомості.




Який рядок треба записати на місце \kw{XXX} ?\par
\vspace{5mm}


Який рядок треба записати на місце \kw{YYY} ?\par


\newpage

{\begin {center}\textbf{ 87}\end {center}}\par
Файл \code{'1.txt'} містить відомість з рядками, в яких через двокрапку записано поряд\-ковий номер (ціле), назву предмету (знаків пунктуації не містить), прізвище студента, ім'я студента, номер заліковки, оцінку в балах за 100 бальною шкалою.

Білі символи навколо роздільників не допускаються.

Приклад рядка файлу



\begin{verbatim}
13:algebra:surname:name:123456:77
\end{verbatim}


Нижче наведено код, який шукає усі назви предметів, що зустрічаються у файлі, та зберігає їх у змінній \id{codes}.



\begin{verbatim}
codes = set()
regex = re.compile(r'XXX')
with open('1.txt', encoding='utf8') as f:
    for s in f:
        res = re.fullmatch(regex, s)
        if not res: continue
        s = ''
        codes.add(r'YYY')
\end{verbatim}





Який рядок треба записати на місце \kw{XXX} ?\par
\vspace{5mm}


Який рядок треба записати на місце \kw{YYY} ?\par


\newpage

{\begin {center}\textbf{ 88}\end {center}}\par
Файл \code{'1.txt'} містить відомість з рядками, в яких через кому записано поряд\-ковий номер (ціле), назву предмету (знаків пунктуації не містить), прізвище студента, ім'я студента, номер заліковки, оцінку в балах за 100 бальною шкалою.

Білі символи навколо роздільників не допускаються.

Приклад рядка файлу



\begin{verbatim}
13,algebra,surname,name,123456,77
\end{verbatim}


Нижче наведено код, який шукає усі предмети, який складав студент із заліковкою \code{777}, та зберігає їх типи в змінній \kw{codes}.



\begin{verbatim}
codes = set()
regex = re.compile(r'XXX')
with open('1.txt', encoding='utf8') as f:
    for s in f:
        res = re.fullmatch(regex, s)
        if not res: continue
        s = ''
        codes.add(r'YYY')
\end{verbatim}





Який рядок треба записати на місце \kw{XXX} ?\par
\vspace{5mm}


Який рядок треба записати на місце \kw{YYY} ?\par


\newpage

{\begin {center}\textbf{ 89}\end {center}}\par
Файл \code{'1.txt'} містить відомість з рядками, в яких через кому записа\-но поряд\-ковий номер (ціле), назву предмету (знаків пунктуації не містить), прізвище студента, ім'я студента, номер заліковки, оцінку в балах за 100 бальною шкалою.

Білі символи навколо роздільників не допускаються.

Приклад рядка файлу



\begin{verbatim}
13,algebra,surname,name,123456,77
\end{verbatim}


Нижче наведено код, який має залишити у файлі в кожному рядку тільки номер заліковки, назву предмету та оцінку в балах за 100-бальною шкалою (у заданій послі\-дов\-ності).

Рядок з прикладу має бути замінений на такий



\begin{verbatim}
123456,algebra,77
\end{verbatim}




Код



\begin{verbatim}
regex = re.compile(r'XXX')
with open('1.txt', encoding='utf8') as f:
    lst = f.readlines()
for i in range(len(lst)):
    lst[i] = re.sub(regex, r'YYY', lst[i])
with open('1.txt', 'w', encoding='utf8') as f:
    f.writelines(lst)
\end{verbatim}







Який рядок треба записати на місце \kw{XXX} ?\par
\vspace{5mm}


Який рядок треба записати на місце \kw{YYY} ?\par


\newpage

{\begin {center}\textbf{ 90}\end {center}}\par
Рядки журналу через двокрапку містять ім'я користувача (складається з латиниці), час події (фор\-мат як у прикладі), тип події (ERROR, WARNING, INFO, STATUS), код події (ціле).

Білі символи навколо роздільників не допускаються.

Приклад такого рядка присвоєно змінній \kw{s}.



\begin{verbatim}
s = 'username:2022-05-05 13-26:ERROR:404'
\end{verbatim}


Нижче наведено код, за виконання якого значенням змінної \kw{out} має стати код події.

Для наведеного прикладу значенням змінної \kw{out} має стати рядок \code{404}



\begin{verbatim}
res = re.fullmatch(r'XXX', s)
s = ''
out = r'YYY'
\end{verbatim}


Код має працювати не тільки для конкретного рядка з прикладу, але й для кожного рядка з журналу.




Який рядок треба записати на місце \kw{XXX} ?\par
\vspace{5mm}


Який рядок треба записати на місце \kw{YYY} ?\par


\newpage

{\begin {center}\textbf{ 91}\end {center}}\par
Рядки відомості через двокрапку містять порядковий номер (ціле), назву предмету (знаків пунктуації не містить), прізвище студента, ім'я студента, оцінку в балах за 100 бальною шкалою.

Білі символи навколо роздільників не допускаються.

Приклад такого рядка присвоєно змінній \kw{s}.



\begin{verbatim}
s = '13:algebra:surname:name:77'
\end{verbatim}


Нижче наведено код, за виконання якого значенням змінної \kw{out} має стати прізвище студента.

Для наведеного прикладу значенням змінної \kw{out} має стати рядок \code{surname}



\begin{verbatim}
res = re.fullmatch(r'XXX', s)
s = ''
out = r'YYY'
\end{verbatim}


Код має працювати не тільки для конкретного рядка з прикладу, але й для кожного рядка з відомості.




Який рядок треба записати на місце \kw{XXX} ?\par
\vspace{5mm}


Який рядок треба записати на місце \kw{YYY} ?\par


\newpage

{\begin {center}\textbf{ 92}\end {center}}\par
Файл \code{'1.txt'} містить журнал з рядками, в яких через двокрапку записано поряд\-ко\-вий номер запису в журналі (ціле), код події (ціле), тип події (ERROR, \\WARNING, INFO, STATUS), ім'я користувача (складається з латиниці), час події (фор\-мат як у прикладі), опис події.

Білі символи навколо роздільників не допускаються.

Приклад рядка файлу



\begin{verbatim}
13:404:ERROR:username:2022-05-05 13-26:something happens
\end{verbatim}


Нижче наведено код, який має залишити у файлі в кожному рядку тільки код події та ім'я користувача. Рядок з прикладу має бути замінений на такий



\begin{verbatim}
404:username
\end{verbatim}




Код:



\begin{verbatim}
regex = re.compile(r'XXX')
with open('1.txt', encoding='utf8') as f:
    lst = f.readlines()
for i in range(len(lst)):
    lst[i] = re.sub(regex, r'YYY', lst[i])
with open('1.txt', 'w', encoding='utf8') as f:
    f.writelines(lst)
\end{verbatim}







Який рядок треба записати на місце \kw{XXX} ?\par
\vspace{5mm}


Який рядок треба записати на місце \kw{YYY} ?\par


\newpage

{\begin {center}\textbf{ 93}\end {center}}\par
Рядки журналу через кому містять ім'я користувача (складається з латиниці), час події (фор\-мат як у прикладі), тип події (ERROR, WARNING, INFO, STATUS), код події (ціле).

Білі символи навколо роздільників не допускаються.

Приклад такого рядка присвоєно змінній \kw{s}.



\begin{verbatim}
s = 'username,2022-05-05 13:26,ERROR,404'
\end{verbatim}


Нижче наведено код, за виконання якого значенням змінної \kw{out} має стати код події.

Для наведеного прикладу значенням змінної \kw{out} має стати рядок \code{404}



\begin{verbatim}
res = re.fullmatch(r'XXX', s)
s = ''
out = r'YYY'
\end{verbatim}


Код має працювати не тільки для конкретного рядка з прикладу, але й для кожного рядка з журналу.




Який рядок треба записати на місце \kw{XXX} ?\par
\vspace{5mm}


Який рядок треба записати на місце \kw{YYY} ?\par


\newpage

{\begin {center}\textbf{ 94}\end {center}}\par
Файл \code{'1.txt'} містить журнал з рядками, в яких через двокрапку записано поряд\-ко\-вий номер запису в журналі (ціле), тип події (ERROR, WARNING, INFO, STATUS), ім'я користувача (складається з латиниці), час події (фор\-мат як у прикладі), код події (ціле).

Білі символи навколо роздільників не допускаються.

Приклад рядка файлу



\begin{verbatim}
13:ERROR:username:2022-05-05 13-26:404
\end{verbatim}


Нижче наведено код, який має замінити у файлі імена користувачів на рядок \code{unknown}.



\begin{verbatim}
regex = re.compile(r'XXX')
with open('1.txt', encoding='utf8') as f:
    lst = f.readlines()
for i in range(len(lst)):
    lst[i] = re.sub(regex, r'YYY', lst[i])
with open('1.txt', 'w', encoding='utf8') as f:
    f.writelines(lst)
\end{verbatim}







Який рядок треба записати на місце \kw{XXX} ?\par
\vspace{5mm}


Який рядок треба записати на місце \kw{YYY} ?\par


\newpage

{\begin {center}\textbf{ 95}\end {center}}\par
Рядки журналу через двокрапку містять тип події (ERROR, WARNING, INFO, STATUS), ім'я користувача (складається з латиниці), час події (фор\-мат як у прикладі), код події (ціле).

Білі символи навколо роздільників не допускаються.

Приклад такого рядка присвоєно змінній \kw{s}.



\begin{verbatim}
s = 'ERROR:username:2022-05-05 13-26:404'
\end{verbatim}


Нижче наведено код, за виконання якого значенням змінної \kw{out} має стати ім'я користувача.

Для наведеного прикладу значенням змінної \kw{out} має стати рядок \code{username}



\begin{verbatim}
res = re.fullmatch(r'XXX', s)
s = ''
out = r'YYY'
\end{verbatim}


Код має працювати не тільки для конкретного рядка з прикладу, але й для кожного рядка з журналу.




Який рядок треба записати на місце \kw{XXX} ?\par
\vspace{5mm}


Який рядок треба записати на місце \kw{YYY} ?\par


\newpage

{\begin {center}\textbf{ 96}\end {center}}\par
Файл \code{'1.txt'} містить журнал з рядками, в яких через кому записано поряд\-ко\-вий номер запису в журналі (ціле), тип події (ERROR, WARNING, INFO, STATUS), ім'я користувача (складається з латиниці), час події (фор\-мат як у прикладі), опис події.

Білі символи навколо роздільників не допускаються.

Приклад рядка файлу



\begin{verbatim}
13,ERROR,username,2022-05-05 13:26,something happens
\end{verbatim}


Нижче наведено код, який має залишити у файлі в кожному рядку тільки тип події та її опис.

Рядок з прикладу має бути замінений на такий



\begin{verbatim}
ERROR,something happens
\end{verbatim}




Код:



\begin{verbatim}
regex = re.compile(r'XXX')
with open('1.txt', encoding='utf8') as f:
    lst = f.readlines()
for i in range(len(lst)):
    lst[i] = re.sub(regex, r'YYY', lst[i])
with open('1.txt', 'w', encoding='utf8') as f:
    f.writelines(lst)
\end{verbatim}









Який рядок треба записати на місце \kw{XXX} ?\par
\vspace{5mm}


Який рядок треба записати на місце \kw{YYY} ?\par


\newpage

{\begin {center}\textbf{ 97}\end {center}}\par
Файл \code{'1.txt'} містить відомість з рядками, в яких через двокрапку записано поряд\-ковий номер (ціле), назву предмету (знаків пунктуації не містить), прізвище студента, ім'я студента, номер заліковки, оцінку в балах за 100 бальною шкалою.

Білі символи навколо роздільників не допускаються.

Приклад рядка файлу



\begin{verbatim}
13:algebra:surname:name:123456:77
\end{verbatim}


Нижче наведено код, який шукає усі предмети, який складав студент із заліковкою \code{777}, та зберігає їх типи в змінній \kw{codes}.



\begin{verbatim}
codes = set()
regex = re.compile(r'XXX')
with open('1.txt', encoding='utf8') as f:
    for s in f:
        res = re.fullmatch(regex, s)
        if not res: continue
        s = ''
        codes.add(r'YYY')
\end{verbatim}





Який рядок треба записати на місце \kw{XXX} ?\par
\vspace{5mm}


Який рядок треба записати на місце \kw{YYY} ?\par


\newpage

{\begin {center}\textbf{ 98}\end {center}}\par
Файл \code{'1.txt'} містить відомість з рядками, в яких через кому записано поряд\-ковий номер (ціле), назву предмету (знаків пунктуації не містить), прізвище студента, ім'я студента, номер заліковки, оцінку в балах за 100 бальною шкалою.

Білі символи навколо роздільників не допускаються.

Приклад рядка файлу



\begin{verbatim}
13,algebra,surname,name,123456,77
\end{verbatim}


Нижче наведено код, який має залишити у файлі в кожному рядку тільки номер заліковки та предмет.

Рядок з прикладу має бути замінений на такий



\begin{verbatim}
123456,algebra
\end{verbatim}




Код



\begin{verbatim}
regex = re.compile(r'XXX')
with open('1.txt', encoding='utf8') as f:
    lst = f.readlines()
for i in range(len(lst)):
    lst[i] = re.sub(regex, r'YYY', lst[i])
with open('1.txt', 'w', encoding='utf8') as f:
    f.writelines(lst)
\end{verbatim}









Який рядок треба записати на місце \kw{XXX} ?\par
\vspace{5mm}


Який рядок треба записати на місце \kw{YYY} ?\par


\newpage

{\begin {center}\textbf{ 99}\end {center}}\par
Файл \code{'1.txt'} містить відомість з рядками, в яких через кому записано поряд\-ковий номер (ціле), назву предмету (знаків пунктуації не містить), прізвище студента, ім'я студента, номер заліковки, оцінку в балах за 100 бальною шкалою.

Білі символи навколо роздільників не допускаються.

Приклад рядка файлу



\begin{verbatim}
13,algebra,surname,name,123456,77
\end{verbatim}


Нижче наведено код, який має в кожному рядку файлу переставити місцями назву предмету та номер заліковки.

Рядок з прикладу має бути замінений на такий



\begin{verbatim}
13,123456,surname,name,algebra,77
\end{verbatim}




Код



\begin{verbatim}
regex = re.compile(r'XXX')
with open('1.txt', encoding='utf8') as f:
    lst = f.readlines()
for i in range(len(lst)):
    lst[i] = re.sub(regex, r'YYY', lst[i])
with open('1.txt', 'w', encoding='utf8') as f:
    f.writelines(lst)
\end{verbatim}







Який рядок треба записати на місце \kw{XXX} ?\par
\vspace{5mm}


Який рядок треба записати на місце \kw{YYY} ?\par


\newpage

{\begin {center}\textbf{ 100}\end {center}}\par
Рядки відомості через двокрапку містять порядковий номер (ціле), пріз\-ви\-ще студен\-та, ім'я студента, номер заліковки, оцінку в балах за 100 бальною шкалою.

Білі символи навколо роздільників не допускаються.

Приклад такого рядка присвоєно змінній \kw{s}.



\begin{verbatim}
s = '13:surname:name:123456:77'
\end{verbatim}


Нижче наведено код, за виконання якого для рядка \kw{s} значенням змінної \kw{out} має стати номер заліковки.

Для наведеного прикладу значенням змінної \kw{out} має стати рядок \code{123456} .



\begin{verbatim}
res = re.fullmatch(r'XXX', s)
s = ''
out = r'YYY'
\end{verbatim}


Код має працювати не тільки для конкретного рядка з прикладу, але й для кожного рядка з відомості.




Який рядок треба записати на місце \kw{XXX} ?\par
\vspace{5mm}


Який рядок треба записати на місце \kw{YYY} ?\par


\newpage

{\begin {center}\textbf{ 101}\end {center}}\par
Файл \code{'1.txt'} містить журнал з рядками, в яких через кому записано поряд\-ко\-вий номер запису (ціле), ім'я користувача (складається з латиниці), час події (фор\-мат як у прикладі), тип події (ERROR, WARNING, INFO, STATUS), код події (ціле).

Білі символи навколо роздільників не допускаються.

Приклад рядка журналу



\begin{verbatim}
12,username,2022-05-05 13:26,ERROR,404
\end{verbatim}


Нижче наведено код, який шукає усі попередження (подія WARNING), що зустрі\-ча\-ються в журналі і відбулися під час обідньої перерви деканату (з 13,00 по 13,59 включно), та зберігає їх коди в змінній \kw{codes}.



\begin{verbatim}
codes = set()
regex = re.compile(r'XXX')
with open('1.txt', encoding='utf8') as f:
    for s in f:
        res = re.fullmatch(regex, s)
        if not res: continue
        s = ''
        codes.add(r'YYY')
\end{verbatim}







Який рядок треба записати на місце \kw{XXX} ?\par
\vspace{5mm}


Який рядок треба записати на місце \kw{YYY} ?\par


\newpage

{\begin {center}\textbf{ 102}\end {center}}\par
Файл \code{'1.txt'} містить журнал з рядками, в яких через двокрапку записано поряд\-ко\-вий номер запису в журналі (ціле), тип події (ERROR, WARNING, INFO, STATUS), ім'я користувача (складається з латиниці), час події (фор\-мат як у прикладі), код події (ціле).

Білі символи навколо роздільників не допускаються.

Приклад рядка файлу



\begin{verbatim}
13:ERROR:username:2022-05-05 13-26:404
\end{verbatim}


Нижче наведено код, який шукає усі коди помилок (подія ERROR), що зустріча\-ють\-ся в журналі, та зберігає їх у змінній \kw{codes}.



\begin{verbatim}
codes = set()
regex = re.compile(r'XXX')
with open('1.txt', encoding='utf8') as f:
    for s in f:
        res = re.fullmatch(regex, s)
        if not res: continue
        s = ''
        codes.add(r'YYY')
\end{verbatim}





Який рядок треба записати на місце \kw{XXX} ?\par
\vspace{5mm}


Який рядок треба записати на місце \kw{YYY} ?\par


\newpage

{\begin {center}\textbf{ 103}\end {center}}\par
Файл \code{'1.txt'} містить журнал з рядками, в яких через кому записано поряд\-ко\-вий номер запису в журналі (ціле), тип події (ERROR, WARNING, INFO, STATUS), ім'я користувача (складається з латиниці), час події (фор\-мат як у прикладі), опис події.

Білі символи навколо роздільників не допускаються.

Приклад рядка файлу



\begin{verbatim}
13,ERROR,username,2022-05-05 13:26,something happens
\end{verbatim}


Нижче наведено код, який має залишити у файлі в кожному рядку тільки тип події та її опис.

Рядок з прикладу має бути замінений на такий



\begin{verbatim}
ERROR,something happens
\end{verbatim}




Код:



\begin{verbatim}
regex = re.compile(r'XXX')
with open('1.txt', encoding='utf8') as f:
    lst = f.readlines()
for i in range(len(lst)):
    lst[i] = re.sub(regex, r'YYY', lst[i])
with open('1.txt', 'w', encoding='utf8') as f:
    f.writelines(lst)
\end{verbatim}









Який рядок треба записати на місце \kw{XXX} ?\par
\vspace{5mm}


Який рядок треба записати на місце \kw{YYY} ?\par


\newpage

{\begin {center}\textbf{ 104}\end {center}}\par
Файл \code{'1.txt'} містить журнал з рядками, в яких через кому записано поряд\-ко\-вий номер запису в журналі (ціле), тип події (ERROR, WARNING, INFO, STATUS), ім'я користувача (складається з латиниці), час події (фор\-мат як у прикладі), код події (ціле).

Білі символи навколо роздільників не допускаються.

Приклад рядка файлу



\begin{verbatim}
13,ERROR,username,2022-05-05 13:26,404
\end{verbatim}


Нижче наведено код, який шукає усі коди помилок (подія ERROR), що зустріча\-ють\-ся в журналі, та зберігає їх у змінній \kw{codes}.



\begin{verbatim}
codes = set()
regex = re.compile(r'XXX')
with open('1.txt', encoding='utf8') as f:
    for s in f:
        res = re.fullmatch(regex, s)
        if not res: continue
        s = ''
        codes.add(r'YYY')
\end{verbatim}





Який рядок треба записати на місце \kw{XXX} ?\par
\vspace{5mm}


Який рядок треба записати на місце \kw{YYY} ?\par


\newpage

{\begin {center}\textbf{ 105}\end {center}}\par
Файл \code{'1.txt'} містить відомість з рядками, в яких через двокрапку записано поряд\-ковий номер (ціле), назву предмету (знаків пунктуації не містить), прізвище студента, ім'я студента, номер заліковки, оцінку в балах за 100 бальною шкалою.

Білі символи навколо роздільників не допускаються.

Приклад рядка файлу



\begin{verbatim}
13:algebra:surname:name:123456:77
\end{verbatim}


Нижче наведено код, який шукає усі назви предметів, що зустрічаються у файлі, та зберігає їх у змінній \id{codes}.



\begin{verbatim}
codes = set()
regex = re.compile(r'XXX')
with open('1.txt', encoding='utf8') as f:
    for s in f:
        res = re.fullmatch(regex, s)
        if not res: continue
        s = ''
        codes.add(r'YYY')
\end{verbatim}





Який рядок треба записати на місце \kw{XXX} ?\par
\vspace{5mm}


Який рядок треба записати на місце \kw{YYY} ?\par


\newpage

{\begin {center}\textbf{ 106}\end {center}}\par
Файл \code{'1.txt'} містить відомість з рядками, в яких через двокрапку записано поряд\-ковий номер (ціле), назву предмету (знаків пунктуації не містить), прізвище студента, ім'я студента, номер заліковки, оцінку в балах за 100 бальною шкалою.

Білі символи навколо роздільників не допускаються.

Приклад рядка файлу



\begin{verbatim}
13:algebra:surname:name:123456:77
\end{verbatim}


Нижче наведено код, який має в кожному рядку файлу переставити місцями назву предмету та номер заліковки.

Рядок з прикладу має бути замінений на такий



\begin{verbatim}
13:123456:surname:name:algebra:77
\end{verbatim}




Код:



\begin{verbatim}
regex = re.compile(r'XXX')
with open('1.txt', encoding='utf8') as f:
    lst = f.readlines()
for i in range(len(lst)):
    lst[i] = re.sub(regex, r'YYY', lst[i])
with open('1.txt', 'w', encoding='utf8') as f:
    f.writelines(lst)
\end{verbatim}







Який рядок треба записати на місце \kw{XXX} ?\par
\vspace{5mm}


Який рядок треба записати на місце \kw{YYY} ?\par


\newpage

{\begin {center}\textbf{ 107}\end {center}}\par
Файл \code{'1.txt'} містить журнал з рядками, в яких через двокрапку записано поряд\-ко\-вий номер запису в журналі (ціле), код події (ціле), тип події (ERROR,\\ WARNING, INFO, STATUS), ім'я користувача (складається з латиниці), час події (фор\-мат як у прикладі), опис події.

Білі символи навколо роздільників не допускаються.

Приклад рядка файлу



\begin{verbatim}
13:404:ERROR:username:2022-05-05 13:26:something happens
\end{verbatim}


Нижче наведено код, який має в кожному рядку файли переставити місцями код та тип події.

Рядок з прикладу має бути замінений на такий



\begin{verbatim}
13:ERROR:404:username:2022-05-05 13:26:something happens
\end{verbatim}




Код:



\begin{verbatim}
regex = re.compile(r'XXX')
with open('1.txt', encoding='utf8') as f:
    lst = f.readlines()
for i in range(len(lst)):
    lst[i] = re.sub(regex, r'YYY', lst[i])
with open('1.txt', 'w', encoding='utf8') as f:
    f.writelines(lst)
\end{verbatim}







Який рядок треба записати на місце \kw{XXX} ?\par
\vspace{5mm}


Який рядок треба записати на місце \kw{YYY} ?\par


\newpage

{\begin {center}\textbf{ 108}\end {center}}\par
Файл \code{'1.txt'} містить відомість з рядками, в яких через кому записа\-но поряд\-ковий номер (ціле), назву предмету (знаків пунктуації не містить), прізвище студента, ім'я студента, номер заліковки, оцінку в балах за 100 бальною шкалою.

Білі символи навколо роздільників не допускаються.

Приклад рядка файлу



\begin{verbatim}
13,algebra,surname,name,123456,77
\end{verbatim}


Нижче наведено код, який має залишити у файлі в кожному рядку тільки номер заліковки, назву предмету та оцінку в балах за 100-бальною шкалою (у заданій послі\-дов\-ності).

Рядок з прикладу має бути замінений на такий



\begin{verbatim}
123456,algebra,77
\end{verbatim}




Код



\begin{verbatim}
regex = re.compile(r'XXX')
with open('1.txt', encoding='utf8') as f:
    lst = f.readlines()
for i in range(len(lst)):
    lst[i] = re.sub(regex, r'YYY', lst[i])
with open('1.txt', 'w', encoding='utf8') as f:
    f.writelines(lst)
\end{verbatim}







Який рядок треба записати на місце \kw{XXX} ?\par
\vspace{5mm}


Який рядок треба записати на місце \kw{YYY} ?\par


\newpage

{\begin {center}\textbf{ 109}\end {center}}\par
Файл \code{'1.txt'} містить журнал з рядками, в яких через кому записано поряд\-ко\-вий номер запису в журналі (ціле), код події (ціле), тип події (ERROR, \\WARNING, INFO, STATUS), ім'я користувача (складається з латиниці), час події (фор\-мат як у прикладі), опис події.

Білі символи навколо роздільників не допускаються.

Приклад рядка файлу



\begin{verbatim}
13,404,ERROR,username,2022-05-05 13:26,something happens
\end{verbatim}


Нижче наведено код, який має залишити у файлі в кожному рядку тільки код події та ім'я користувача. Рядок з прикладу має бути замінений на такий



\begin{verbatim}
404,username
\end{verbatim}




Код:



\begin{verbatim}
regex = re.compile(r'XXX')
with open('1.txt', encoding='utf8') as f:
    lst = f.readlines()
for i in range(len(lst)):
    lst[i] = re.sub(regex, r'YYY', lst[i])
with open('1.txt', 'w', encoding='utf8') as f:
    f.writelines(lst)
\end{verbatim}







Який рядок треба записати на місце \kw{XXX} ?\par
\vspace{5mm}


Який рядок треба записати на місце \kw{YYY} ?\par


\newpage

{\begin {center}\textbf{ 110}\end {center}}\par
Файл \code{'1.txt'} містить відомість з рядками, в яких через кому записано поряд\-ковий номер (ціле), назву предмету (знаків пунктуації не містить), прізвище студента, ім'я студента, номер заліковки, оцінку в балах за 100 бальною шкалою.

Білі символи навколо роздільників не допускаються.

Приклад рядка файлу



\begin{verbatim}
13,algebra,surname,name,123456,77
\end{verbatim}


Нижче наведено код, який шукає усі назви предметів, що зустрічаються у файлі, та зберігає їх у змінній \id{codes}.



\begin{verbatim}
codes = set()
regex = re.compile(r'XXX')
with open('1.txt', encoding='utf8') as f:
    for s in f:
        res = re.fullmatch(regex, s)
        if not res: continue
        s = ''
        codes.add(r'YYY')
\end{verbatim}





Який рядок треба записати на місце \kw{XXX} ?\par
\vspace{5mm}


Який рядок треба записати на місце \kw{YYY} ?\par


\newpage

{\begin {center}\textbf{ 111}\end {center}}\par
Файл \code{'1.txt'} містить журнал з рядками, в яких через кому записано поряд\-ко\-вий номер запису (ціле), ім'я користувача (складається з латиниці), час події (фор\-мат як у прикладі), тип події (ERROR, WARNING, INFO, STATUS), код події (ціле).

Білі символи навколо роздільників не допускаються.

Приклад рядка журналу



\begin{verbatim}
12,username,2022-05-05 13:26,ERROR,404
\end{verbatim}


Нижче наведено код, який шукає усі події, що відбулись\\ для користувача \code{anonymous}, та зберігає їх типи в змінній \kw{codes}.



\begin{verbatim}
codes = set()
regex = re.compile(r'XXX')
with open('1.txt', encoding='utf8') as f:
    for s in f:
        res = re.fullmatch(regex, s)
        if not res: continue
        s = ''
        codes.add(r'YYY')
\end{verbatim}





Який рядок треба записати на місце \kw{XXX} ?\par
\vspace{5mm}


Який рядок треба записати на місце \kw{YYY} ?\par


\newpage

{\begin {center}\textbf{ 112}\end {center}}\par
Файл \code{'1.txt'} містить відомість з рядками, в яких через кому записано поряд\-ковий номер (ціле), назву предмету (знаків пунктуації не містить), прізвище студента, ім'я студента, номер заліковки, оцінку в балах за 100 бальною шкалою.

Білі символи навколо роздільників не допускаються.

Приклад рядка файлу



\begin{verbatim}
13,algebra,surname,name,123456,77
\end{verbatim}


Нижче наведено код, який шукає усі предмети, який складав студент із заліковкою \code{777}, та зберігає їх типи в змінній \kw{codes}.



\begin{verbatim}
codes = set()
regex = re.compile(r'XXX')
with open('1.txt', encoding='utf8') as f:
    for s in f:
        res = re.fullmatch(regex, s)
        if not res: continue
        s = ''
        codes.add(r'YYY')
\end{verbatim}





Який рядок треба записати на місце \kw{XXX} ?\par
\vspace{5mm}


Який рядок треба записати на місце \kw{YYY} ?\par


\newpage

{\begin {center}\textbf{ 113}\end {center}}\par
Файл \code{'1.txt'} містить журнал з рядками, в яких через двокрапку записано поряд\-ко\-вий номер запису в журналі (ціле), тип події (ERROR, WARNING, INFO, STATUS), ім'я користувача (складається з латиниці), час події (фор\-мат як у прикладі), код події (ціле).

Білі символи навколо роздільників не допускаються.

Приклад рядка файлу



\begin{verbatim}
13:ERROR:username:2022-05-05 13-26:404
\end{verbatim}


Нижче наведено код, який має замінити у файлі імена користувачів на рядок \code{unknown}.



\begin{verbatim}
regex = re.compile(r'XXX')
with open('1.txt', encoding='utf8') as f:
    lst = f.readlines()
for i in range(len(lst)):
    lst[i] = re.sub(regex, r'YYY', lst[i])
with open('1.txt', 'w', encoding='utf8') as f:
    f.writelines(lst)
\end{verbatim}







Який рядок треба записати на місце \kw{XXX} ?\par
\vspace{5mm}


Який рядок треба записати на місце \kw{YYY} ?\par


\newpage

{\begin {center}\textbf{ 114}\end {center}}\par
Рядки журналу через кому містять ім'я користувача (складається з латиниці), час події (фор\-мат як у прикладі), тип події (ERROR, WARNING, INFO, STATUS), код події (ціле).

Білі символи навколо роздільників не допускаються.

Приклад такого рядка присвоєно змінній \kw{s}.



\begin{verbatim}
s = 'username,2022-05-05 13:26,ERROR,404'
\end{verbatim}


Нижче наведено код, за виконання якого значенням змінної \kw{out} має стати код події.

Для наведеного прикладу значенням змінної \kw{out} має стати рядок \code{404}



\begin{verbatim}
res = re.fullmatch(r'XXX', s)
s = ''
out = r'YYY'
\end{verbatim}


Код має працювати не тільки для конкретного рядка з прикладу, але й для кожного рядка з журналу.




Який рядок треба записати на місце \kw{XXX} ?\par
\vspace{5mm}


Який рядок треба записати на місце \kw{YYY} ?\par


\newpage

{\begin {center}\textbf{ 115}\end {center}}\par
Файл \code{'1.txt'} містить журнал з рядками, в яких через кому записано поряд\-ко\-вий номер запису (ціле), ім'я користувача (складається з латиниці), час події (фор\-мат як у прикладі), тип події (ERROR, WARNING, INFO, STATUS), код події (ціле).

Білі символи навколо роздільників не допускаються.

Приклад рядка журналу



\begin{verbatim}
12,username,2022-05-05 13:26,ERROR,404
\end{verbatim}


Нижче наведено код, який шукає усі попередження (подія WARNING), що зустрі\-ча\-ються в журналі і відбулися під час обідньої перерви деканату (з 13,00 по 13,59 включно), та зберігає їх коди в змінній \kw{codes}.



\begin{verbatim}
codes = set()
regex = re.compile(r'XXX')
with open('1.txt', encoding='utf8') as f:
    for s in f:
        res = re.fullmatch(regex, s)
        if not res: continue
        s = ''
        codes.add(r'YYY')
\end{verbatim}







Який рядок треба записати на місце \kw{XXX} ?\par
\vspace{5mm}


Який рядок треба записати на місце \kw{YYY} ?\par


\newpage

{\begin {center}\textbf{ 116}\end {center}}\par
Файл \code{'1.txt'} містить журнал з рядками, в яких через двокрапку записано поряд\-ко\-вий номер запису в журналі (ціле), тип події (ERROR, WARNING, INFO, STATUS), ім'я користувача (складається з латиниці), час події (фор\-мат як у прикладі), опис події.

Білі символи навколо роздільників не допускаються.

Приклад рядка файлу



\begin{verbatim}
13:ERROR:username:2022-05-05 13-26:something happens
\end{verbatim}


Нижче наведено код, який має залишити у файлі в кожному рядку тільки тип події та її опис.

Рядок з прикладу має бути замінений на такий



\begin{verbatim}
ERROR:something happens
\end{verbatim}




Код:



\begin{verbatim}
regex = re.compile(r'XXX')
with open('1.txt', encoding='utf8') as f:
    lst = f.readlines()
for i in range(len(lst)):
    lst[i] = re.sub(regex, r'YYY', lst[i])
with open('1.txt', 'w', encoding='utf8') as f:
    f.writelines(lst)
\end{verbatim}









Який рядок треба записати на місце \kw{XXX} ?\par
\vspace{5mm}


Який рядок треба записати на місце \kw{YYY} ?\par


\newpage

{\begin {center}\textbf{ 117}\end {center}}\par
Рядки відомості через двокрапку містять порядковий номер (ціле), пріз\-ви\-ще студен\-та, ім'я студента, номер заліковки, оцінку в балах за 100 бальною шкалою.

Білі символи навколо роздільників не допускаються.

Приклад такого рядка присвоєно змінній \kw{s}.



\begin{verbatim}
s = '13:surname:name:123456:77'
\end{verbatim}


Нижче наведено код, за виконання якого для рядка \kw{s} значенням змінної \kw{out} має стати номер заліковки.

Для наведеного прикладу значенням змінної \kw{out} має стати рядок \code{123456} .



\begin{verbatim}
res = re.fullmatch(r'XXX', s)
s = ''
out = r'YYY'
\end{verbatim}


Код має працювати не тільки для конкретного рядка з прикладу, але й для кожного рядка з відомості.




Який рядок треба записати на місце \kw{XXX} ?\par
\vspace{5mm}


Який рядок треба записати на місце \kw{YYY} ?\par


\newpage

{\begin {center}\textbf{ 118}\end {center}}\par
Рядки відомості через двокрапку містять порядковий номер (ціле), назву предмету (знаків пунктуації не містить), прізвище студента, ім'я студента, оцінку в балах за 100 бальною шкалою.

Білі символи навколо роздільників не допускаються.

Приклад такого рядка присвоєно змінній \kw{s}.



\begin{verbatim}
s = '13:algebra:surname:name:77'
\end{verbatim}


Нижче наведено код, за виконання якого значенням змінної \kw{out} має стати прізвище студента.

Для наведеного прикладу значенням змінної \kw{out} має стати рядок \code{surname}



\begin{verbatim}
res = re.fullmatch(r'XXX', s)
s = ''
out = r'YYY'
\end{verbatim}


Код має працювати не тільки для конкретного рядка з прикладу, але й для кожного рядка з відомості.




Який рядок треба записати на місце \kw{XXX} ?\par
\vspace{5mm}


Який рядок треба записати на місце \kw{YYY} ?\par


\newpage

{\begin {center}\textbf{ 119}\end {center}}\par
Файл \code{'1.txt'} містить відомість з рядками, в яких через двокрапку записано поряд\-ковий номер (ціле), назву предмету (знаків пунктуації не містить), прізвище студента, ім'я студента, номер заліковки, оцінку в балах за 100 бальною шкалою.

Білі символи навколо роздільників не допускаються.

Приклад рядка файлу



\begin{verbatim}
13:algebra:surname:name:123456:77
\end{verbatim}


Нижче наведено код, який шукає усі предмети, який складав студент із заліковкою \code{777}, та зберігає їх типи в змінній \kw{codes}.



\begin{verbatim}
codes = set()
regex = re.compile(r'XXX')
with open('1.txt', encoding='utf8') as f:
    for s in f:
        res = re.fullmatch(regex, s)
        if not res: continue
        s = ''
        codes.add(r'YYY')
\end{verbatim}





Який рядок треба записати на місце \kw{XXX} ?\par
\vspace{5mm}


Який рядок треба записати на місце \kw{YYY} ?\par


\newpage

{\begin {center}\textbf{ 120}\end {center}}\par
Файл \code{'1.txt'} містить журнал з рядками, в яких через двокрапку записано поряд\-ко\-вий номер запису в журналі (ціле), код події (ціле), тип події (ERROR, \\WARNING, INFO, STATUS), ім'я користувача (складається з латиниці), час події (фор\-мат як у прикладі), опис події.

Білі символи навколо роздільників не допускаються.

Приклад рядка файлу



\begin{verbatim}
13:404:ERROR:username:2022-05-05 13-26:something happens
\end{verbatim}


Нижче наведено код, який має залишити у файлі в кожному рядку тільки код події та ім'я користувача. Рядок з прикладу має бути замінений на такий



\begin{verbatim}
404:username
\end{verbatim}




Код:



\begin{verbatim}
regex = re.compile(r'XXX')
with open('1.txt', encoding='utf8') as f:
    lst = f.readlines()
for i in range(len(lst)):
    lst[i] = re.sub(regex, r'YYY', lst[i])
with open('1.txt', 'w', encoding='utf8') as f:
    f.writelines(lst)
\end{verbatim}







Який рядок треба записати на місце \kw{XXX} ?\par
\vspace{5mm}


Який рядок треба записати на місце \kw{YYY} ?\par


\newpage

{\begin {center}\textbf{ 121}\end {center}}\par
Рядки журналу через двокрапку містять ім'я користувача (складається з латиниці), час події (фор\-мат як у прикладі), тип події (ERROR, WARNING, INFO, STATUS), код події (ціле).

Білі символи навколо роздільників не допускаються.

Приклад такого рядка присвоєно змінній \kw{s}.



\begin{verbatim}
s = 'username:2022-05-05 13-26:ERROR:404'
\end{verbatim}


Нижче наведено код, за виконання якого значенням змінної \kw{out} має стати код події.

Для наведеного прикладу значенням змінної \kw{out} має стати рядок \code{404}



\begin{verbatim}
res = re.fullmatch(r'XXX', s)
s = ''
out = r'YYY'
\end{verbatim}


Код має працювати не тільки для конкретного рядка з прикладу, але й для кожного рядка з журналу.




Який рядок треба записати на місце \kw{XXX} ?\par
\vspace{5mm}


Який рядок треба записати на місце \kw{YYY} ?\par


\newpage

{\begin {center}\textbf{ 122}\end {center}}\par
Файл \code{'1.txt'} містить відомість з рядками, в яких через двокрапку записа\-но назву предмету (знаків пунктуації не містить), прізвище студента, ім'я студента, номер залі\-ков\-ки, оцінку в балах за 100 бальною шкалою.

Білі символи навколо роздільників не допускаються.

Приклад рядка файлу



\begin{verbatim}
algebra:surname:name:123456:77
\end{verbatim}


Нижче наведено код, який має замінити у файлі усі номери заліковок на рядок \code{***}.



\begin{verbatim}
regex = re.compile(r'XXX')
with open('1.txt', encoding='utf8') as f:
    lst = f.readlines()
for i in range(len(lst)):
    lst[i] = re.sub(regex, r'YYY', lst[i])
with open('1.txt', 'w', encoding='utf8') as f:
    f.writelines(lst)
\end{verbatim}







Який рядок треба записати на місце \kw{XXX} ?\par
\vspace{5mm}


Який рядок треба записати на місце \kw{YYY} ?\par


\newpage

{\begin {center}\textbf{ 123}\end {center}}\par
Файл \code{'1.txt'} містить відомість з рядками, в яких через кому записа\-но назву предмету (знаків пунктуації не містить), прізвище студента, ім'я студента, номер залі\-ков\-ки, оцінку в балах за 100 бальною шкалою.

Білі символи навколо роздільників не допускаються.

Приклад рядка файлу



\begin{verbatim}
algebra,surname,name,123456,77
\end{verbatim}


Нижче наведено код, який має замінити у файлі усі номери заліковок на рядок \code{***}.



\begin{verbatim}
regex = re.compile(r'XXX')
with open('1.txt', encoding='utf8') as f:
    lst = f.readlines()
for i in range(len(lst)):
    lst[i] = re.sub(regex, r'YYY', lst[i])
with open('1.txt', 'w', encoding='utf8') as f:
    f.writelines(lst)
\end{verbatim}







Який рядок треба записати на місце \kw{XXX} ?\par
\vspace{5mm}


Який рядок треба записати на місце \kw{YYY} ?\par


\newpage

{\begin {center}\textbf{ 124}\end {center}}\par
Файл \code{'1.txt'} містить відомість з рядками, в яких через кому записано поряд\-ковий номер (ціле), назву предмету (знаків пунктуації не містить), прізвище студента, ім'я студента, номер заліковки, оцінку в балах за 100 бальною шкалою.

Білі символи навколо роздільників не допускаються.

Приклад рядка файлу



\begin{verbatim}
13,algebra,surname,name,123456,77
\end{verbatim}


Нижче наведено код, який має залишити у файлі в кожному рядку тільки номер заліковки та предмет.

Рядок з прикладу має бути замінений на такий



\begin{verbatim}
123456,algebra
\end{verbatim}




Код



\begin{verbatim}
regex = re.compile(r'XXX')
with open('1.txt', encoding='utf8') as f:
    lst = f.readlines()
for i in range(len(lst)):
    lst[i] = re.sub(regex, r'YYY', lst[i])
with open('1.txt', 'w', encoding='utf8') as f:
    f.writelines(lst)
\end{verbatim}









Який рядок треба записати на місце \kw{XXX} ?\par
\vspace{5mm}


Який рядок треба записати на місце \kw{YYY} ?\par


\newpage

{\begin {center}\textbf{ 125}\end {center}}\par
Файл \code{'1.txt'} містить журнал з рядками, в яких через кому записано поряд\-ко\-вий номер запису в журналі (ціле), тип події (ERROR, WARNING, INFO, STATUS), ім'я користувача (складається з латиниці), час події (фор\-мат як у прикладі), код події (ціле).

Білі символи навколо роздільників не допускаються.

Приклад рядка файлу



\begin{verbatim}
13,ERROR,username,2022-05-05 13:26,404
\end{verbatim}


Нижче наведено код, який має замінити у файлі імена користувачів на рядок \code{unknown}.



\begin{verbatim}
regex = re.compile(r'XXX')
with open('1.txt', encoding='utf8') as f:
    lst = f.readlines()
for i in range(len(lst)):
    lst[i] = re.sub(regex, r'YYY', lst[i])
with open('1.txt', 'w', encoding='utf8') as f:
    f.writelines(lst)
\end{verbatim}







Який рядок треба записати на місце \kw{XXX} ?\par
\vspace{5mm}


Який рядок треба записати на місце \kw{YYY} ?\par


\newpage

{\begin {center}\textbf{ 126}\end {center}}\par
Рядки відомості через кому містять порядковий номер (ціле), пріз\-ви\-ще студен\-та, ім'я студента, номер заліковки, оцінку в балах за 100 бальною шкалою.

Білі символи навколо роздільників не допускаються.

Приклад такого рядка присвоєно змінній \kw{s}.



\begin{verbatim}
s = '13,surname,name,123456,77'
\end{verbatim}


Нижче наведено код, за виконання якого для рядка \kw{s} значенням змінної \kw{out} має стати номер заліковки.

Для наведеного прикладу значенням змінної \kw{out} має стати рядок \code{123456} .



\begin{verbatim}
res = re.fullmatch(r'XXX', s)
s = ''
out = r'YYY'
\end{verbatim}


Код має працювати не тільки для конкретного рядка з прикладу, але й для кожного рядка з відомості.




Який рядок треба записати на місце \kw{XXX} ?\par
\vspace{5mm}


Який рядок треба записати на місце \kw{YYY} ?\par


\newpage

{\begin {center}\textbf{ 127}\end {center}}\par
Рядки журналу через кому містять тип події (ERROR, WARNING, INFO, STATUS), ім'я користувача (складається з латиниці), час події (фор\-мат як у прикладі), код події (ціле).

Білі символи навколо роздільників не допускаються.

Приклад такого рядка присвоєно змінній \kw{s}.



\begin{verbatim}
s = 'ERROR,username,2022-05-05 13:26,404'
\end{verbatim}


Нижче наведено код, за виконання якого значенням змінної \kw{out} має стати ім'я користувача.

Для наведеного прикладу значенням змінної \kw{out} має стати рядок \code{username}



\begin{verbatim}
res = re.fullmatch(r'XXX', s)
s = ''
out = r'YYY'
\end{verbatim}


Код має працювати не тільки для конкретного рядка з прикладу, але й для кожного рядка з журналу.




Який рядок треба записати на місце \kw{XXX} ?\par
\vspace{5mm}


Який рядок треба записати на місце \kw{YYY} ?\par


\newpage

{\begin {center}\textbf{ 128}\end {center}}\par
Файл \code{'1.txt'} містить журнал з рядками, в яких через двокрапку записано поряд\-ко\-вий номер запису в журналі (ціле), тип події (ERROR, WARNING, INFO, STATUS), ім'я користувача (складається з латиниці), час події (фор\-мат як у прикладі), код події (ціле).

Білі символи навколо роздільників не допускаються.

Приклад рядка файлу



\begin{verbatim}
13:ERROR:username:2022-05-05 13-26:404
\end{verbatim}


Нижче наведено код, який шукає усі коди помилок (подія ERROR), що зустріча\-ють\-ся в журналі, та зберігає їх у змінній \kw{codes}.



\begin{verbatim}
codes = set()
regex = re.compile(r'XXX')
with open('1.txt', encoding='utf8') as f:
    for s in f:
        res = re.fullmatch(regex, s)
        if not res: continue
        s = ''
        codes.add(r'YYY')
\end{verbatim}





Який рядок треба записати на місце \kw{XXX} ?\par
\vspace{5mm}


Який рядок треба записати на місце \kw{YYY} ?\par


\newpage

{\begin {center}\textbf{ 129}\end {center}}\par
Файл \code{'1.txt'} містить журнал з рядками, в яких через кому записано поряд\-ко\-вий номер запису в журналі (ціле), код події (ціле), тип події (ERROR,\\ WARNING, INFO, STATUS), ім'я користувача (складається з латиниці), час події (фор\-мат як у прикладі), опис події.

Білі символи навколо роздільників не допускаються.

Приклад рядка файлу



\begin{verbatim}
13,404,ERROR,username,2022-05-05 13:26,something happens
\end{verbatim}


Нижче наведено код, який має в кожному рядку файли переставити місцями код та тип події.

Рядок з прикладу має бути замінений на такий



\begin{verbatim}
13,ERROR,404,username,2022-05-05 13:26,something happens
\end{verbatim}




Код:



\begin{verbatim}
regex = re.compile(r'XXX')
with open('1.txt', encoding='utf8') as f:
    lst = f.readlines()
for i in range(len(lst)):
    lst[i] = re.sub(regex, r'YYY', lst[i])
with open('1.txt', 'w', encoding='utf8') as f:
    f.writelines(lst)
\end{verbatim}







Який рядок треба записати на місце \kw{XXX} ?\par
\vspace{5mm}


Який рядок треба записати на місце \kw{YYY} ?\par


\newpage

{\begin {center}\textbf{ 130}\end {center}}\par
Файл \code{'1.txt'} містить журнал з рядками, в яких через двокрапку записано поряд\-ко\-вий номер запису (ціле), ім'я користувача (складається з латиниці), час події (фор\-мат як у прикладі), тип події (ERROR, WARNING, INFO, STATUS), код події (ціле).

Білі символи навколо роздільників не допускаються.

Приклад рядка журналу



\begin{verbatim}
12:username:2022-05-05 13-26:ERROR:404
\end{verbatim}


Нижче наведено код, який шукає усі події, що відбулись\\ для користувача \code{anonymous}, та зберігає їх типи в змінній \kw{codes}.



\begin{verbatim}
codes = set()
regex = re.compile(r'XXX')
with open('1.txt', encoding='utf8') as f:
    for s in f:
        res = re.fullmatch(regex, s)
        if not res: continue
        s = ''
        codes.add(r'YYY')
\end{verbatim}





Який рядок треба записати на місце \kw{XXX} ?\par
\vspace{5mm}


Який рядок треба записати на місце \kw{YYY} ?\par


\newpage

{\begin {center}\textbf{ 131}\end {center}}\par
Файл \code{'1.txt'} містить журнал з рядками, в яких через двокрапку записано поряд\-ко\-вий номер запису (ціле), ім'я користувача (складається з латиниці), час події (фор\-мат як у прикладі), тип події (ERROR, WARNING, INFO, STATUS), код події (ціле).

Білі символи навколо роздільників не допускаються.

Приклад рядка журналу



\begin{verbatim}
12:username:2022-05-05 13-26:ERROR:404
\end{verbatim}


Нижче наведено код, який шукає усі попередження (подія WARNING), що зустрі\-ча\-ються в журналі і відбулися під час обідньої перерви деканату (з 13:00 по 13:59 включно), та зберігає їх коди в змінній \kw{codes}.



\begin{verbatim}
codes = set()
regex = re.compile(r'XXX')
with open('1.txt', encoding='utf8') as f:
    for s in f:
        res = re.fullmatch(regex, s)
        if not res: continue
        s = ''
        codes.add(r'YYY')
\end{verbatim}







Який рядок треба записати на місце \kw{XXX} ?\par
\vspace{5mm}


Який рядок треба записати на місце \kw{YYY} ?\par


\newpage

{\begin {center}\textbf{ 132}\end {center}}\par
Файл \code{'1.txt'} містить відомість з рядками, в яких через кому записано поряд\-ковий номер (ціле), назву предмету (знаків пунктуації не містить), прізвище студента, ім'я студента, номер заліковки, оцінку в балах за 100 бальною шкалою.

Білі символи навколо роздільників не допускаються.

Приклад рядка файлу



\begin{verbatim}
13,algebra,surname,name,123456,77
\end{verbatim}


Нижче наведено код, який має в кожному рядку файлу переставити місцями назву предмету та номер заліковки.

Рядок з прикладу має бути замінений на такий



\begin{verbatim}
13,123456,surname,name,algebra,77
\end{verbatim}




Код



\begin{verbatim}
regex = re.compile(r'XXX')
with open('1.txt', encoding='utf8') as f:
    lst = f.readlines()
for i in range(len(lst)):
    lst[i] = re.sub(regex, r'YYY', lst[i])
with open('1.txt', 'w', encoding='utf8') as f:
    f.writelines(lst)
\end{verbatim}







Який рядок треба записати на місце \kw{XXX} ?\par
\vspace{5mm}


Який рядок треба записати на місце \kw{YYY} ?\par


\newpage

{\begin {center}\textbf{ 133}\end {center}}\par
Рядки відомості через кому містять порядковий номер (ціле), назву предмету (зна\-ків пунктуації не містить), прізвище студента, ім'я студента, оцінку в балах за 100 бальною шкалою.

Білі символи навколо роздільників не допускаються.

Приклад такого рядка присвоєно змінній \kw{s}.



\begin{verbatim}
s = '13,algebra,surname,name,77'
\end{verbatim}


Нижче наведено код, за виконання якого значенням змінної \kw{out} має стати прізвище студента.

Для наведеного прикладу значенням змінної \kw{out} має стати рядок \code{surname}



\begin{verbatim}
res = re.fullmatch(r'XXX', s)
s = ''
out = r'YYY'
\end{verbatim}


Код має працювати не тільки для конкретного рядка з прикладу, але й для кожного рядка з відомості.




Який рядок треба записати на місце \kw{XXX} ?\par
\vspace{5mm}


Який рядок треба записати на місце \kw{YYY} ?\par


\newpage

{\begin {center}\textbf{ 134}\end {center}}\par
Рядки журналу через двокрапку містять тип події (ERROR, WARNING, INFO, STATUS), ім'я користувача (складається з латиниці), час події (фор\-мат як у прикладі), код події (ціле).

Білі символи навколо роздільників не допускаються.

Приклад такого рядка присвоєно змінній \kw{s}.



\begin{verbatim}
s = 'ERROR:username:2022-05-05 13-26:404'
\end{verbatim}


Нижче наведено код, за виконання якого значенням змінної \kw{out} має стати ім'я користувача.

Для наведеного прикладу значенням змінної \kw{out} має стати рядок \code{username}



\begin{verbatim}
res = re.fullmatch(r'XXX', s)
s = ''
out = r'YYY'
\end{verbatim}


Код має працювати не тільки для конкретного рядка з прикладу, але й для кожного рядка з журналу.




Який рядок треба записати на місце \kw{XXX} ?\par
\vspace{5mm}


Який рядок треба записати на місце \kw{YYY} ?\par


\newpage

{\begin {center}\textbf{ 135}\end {center}}\par
Файл \code{'1.txt'} містить відомість з рядками, в яких через двокрапку записа\-но поряд\-ковий номер (ціле), назву предмету (знаків пунктуації не містить), прізвище студента, ім'я студента, номер заліковки, оцінку в балах за 100 бальною шкалою.

Білі символи навколо роздільників не допускаються.

Приклад рядка файлу



\begin{verbatim}
13:algebra:surname:name:123456:77
\end{verbatim}


Нижче наведено код, який має залишити у файлі в кожному рядку тільки номер заліковки, назву предмету та оцінку в балах за 100-бальною шкалою (у заданій послі\-дов\-ності).

Рядок з прикладу має бути замінений на такий



\begin{verbatim}
123456:algebra:77
\end{verbatim}




Код:



\begin{verbatim}
regex = re.compile(r'XXX')
with open('1.txt', encoding='utf8') as f:
    lst = f.readlines()
for i in range(len(lst)):
    lst[i] = re.sub(regex, r'YYY', lst[i])
with open('1.txt', 'w', encoding='utf8') as f:
    f.writelines(lst)
\end{verbatim}







Який рядок треба записати на місце \kw{XXX} ?\par
\vspace{5mm}


Який рядок треба записати на місце \kw{YYY} ?\par


\newpage

{\begin {center}\textbf{ 136}\end {center}}\par
Файл \code{'1.txt'} містить відомість з рядками, в яких через двокрапку записано поряд\-ковий номер (ціле), назву предмету (знаків пунктуації не містить), прізвище студента, ім'я студента, номер заліковки, оцінку в балах за 100 бальною шкалою.

Білі символи навколо роздільників не допускаються.

Приклад рядка файлу



\begin{verbatim}
13:algebra:surname:name:123456:77
\end{verbatim}


Нижче наведено код, який має залишити у файлі в кожному рядку тільки номер заліковки та предмет.

Рядок з прикладу має бути замінений на такий



\begin{verbatim}
123456:algebra
\end{verbatim}




Код:



\begin{verbatim}
regex = re.compile(r'XXX')
with open('1.txt', encoding='utf8') as f:
    lst = f.readlines()
for i in range(len(lst)):
    lst[i] = re.sub(regex, r'YYY', lst[i])
with open('1.txt', 'w', encoding='utf8') as f:
    f.writelines(lst)
\end{verbatim}









Який рядок треба записати на місце \kw{XXX} ?\par
\vspace{5mm}


Який рядок треба записати на місце \kw{YYY} ?\par


\newpage

{\begin {center}\textbf{ 137}\end {center}}\par
Файл \code{'1.txt'} містить журнал з рядками, в яких через кому записано поряд\-ко\-вий номер запису в журналі (ціле), тип події (ERROR, WARNING, INFO, STATUS), ім'я користувача (складається з латиниці), час події (фор\-мат як у прикладі), код події (ціле).

Білі символи навколо роздільників не допускаються.

Приклад рядка файлу



\begin{verbatim}
13,ERROR,username,2022-05-05 13:26,404
\end{verbatim}


Нижче наведено код, який має замінити у файлі імена користувачів на рядок \code{unknown}.



\begin{verbatim}
regex = re.compile(r'XXX')
with open('1.txt', encoding='utf8') as f:
    lst = f.readlines()
for i in range(len(lst)):
    lst[i] = re.sub(regex, r'YYY', lst[i])
with open('1.txt', 'w', encoding='utf8') as f:
    f.writelines(lst)
\end{verbatim}







Який рядок треба записати на місце \kw{XXX} ?\par
\vspace{5mm}


Який рядок треба записати на місце \kw{YYY} ?\par


\newpage

{\begin {center}\textbf{ 138}\end {center}}\par
Рядки журналу через кому містять ім'я користувача (складається з латиниці), час події (фор\-мат як у прикладі), тип події (ERROR, WARNING, INFO, STATUS), код події (ціле).

Білі символи навколо роздільників не допускаються.

Приклад такого рядка присвоєно змінній \kw{s}.



\begin{verbatim}
s = 'username,2022-05-05 13:26,ERROR,404'
\end{verbatim}


Нижче наведено код, за виконання якого значенням змінної \kw{out} має стати код події.

Для наведеного прикладу значенням змінної \kw{out} має стати рядок \code{404}



\begin{verbatim}
res = re.fullmatch(r'XXX', s)
s = ''
out = r'YYY'
\end{verbatim}


Код має працювати не тільки для конкретного рядка з прикладу, але й для кожного рядка з журналу.




Який рядок треба записати на місце \kw{XXX} ?\par
\vspace{5mm}


Який рядок треба записати на місце \kw{YYY} ?\par


\newpage

{\begin {center}\textbf{ 139}\end {center}}\par
Рядки відомості через двокрапку містять порядковий номер (ціле), назву предмету (знаків пунктуації не містить), прізвище студента, ім'я студента, оцінку в балах за 100 бальною шкалою.

Білі символи навколо роздільників не допускаються.

Приклад такого рядка присвоєно змінній \kw{s}.



\begin{verbatim}
s = '13:algebra:surname:name:77'
\end{verbatim}


Нижче наведено код, за виконання якого значенням змінної \kw{out} має стати прізвище студента.

Для наведеного прикладу значенням змінної \kw{out} має стати рядок \code{surname}



\begin{verbatim}
res = re.fullmatch(r'XXX', s)
s = ''
out = r'YYY'
\end{verbatim}


Код має працювати не тільки для конкретного рядка з прикладу, але й для кожного рядка з відомості.




Який рядок треба записати на місце \kw{XXX} ?\par
\vspace{5mm}


Який рядок треба записати на місце \kw{YYY} ?\par


\newpage

{\begin {center}\textbf{ 140}\end {center}}\par
Файл \code{'1.txt'} містить відомість з рядками, в яких через кому записано поряд\-ковий номер (ціле), назву предмету (знаків пунктуації не містить), прізвище студента, ім'я студента, номер заліковки, оцінку в балах за 100 бальною шкалою.

Білі символи навколо роздільників не допускаються.

Приклад рядка файлу



\begin{verbatim}
13,algebra,surname,name,123456,77
\end{verbatim}


Нижче наведено код, який шукає усі предмети, який складав студент із заліковкою \code{777}, та зберігає їх типи в змінній \kw{codes}.



\begin{verbatim}
codes = set()
regex = re.compile(r'XXX')
with open('1.txt', encoding='utf8') as f:
    for s in f:
        res = re.fullmatch(regex, s)
        if not res: continue
        s = ''
        codes.add(r'YYY')
\end{verbatim}





Який рядок треба записати на місце \kw{XXX} ?\par
\vspace{5mm}


Який рядок треба записати на місце \kw{YYY} ?\par


\newpage

{\begin {center}\textbf{ 141}\end {center}}\par
Файл \code{'1.txt'} містить журнал з рядками, в яких через кому записано поряд\-ко\-вий номер запису (ціле), ім'я користувача (складається з латиниці), час події (фор\-мат як у прикладі), тип події (ERROR, WARNING, INFO, STATUS), код події (ціле).

Білі символи навколо роздільників не допускаються.

Приклад рядка журналу



\begin{verbatim}
12,username,2022-05-05 13:26,ERROR,404
\end{verbatim}


Нижче наведено код, який шукає усі попередження (подія WARNING), що зустрі\-ча\-ються в журналі і відбулися під час обідньої перерви деканату (з 13,00 по 13,59 включно), та зберігає їх коди в змінній \kw{codes}.



\begin{verbatim}
codes = set()
regex = re.compile(r'XXX')
with open('1.txt', encoding='utf8') as f:
    for s in f:
        res = re.fullmatch(regex, s)
        if not res: continue
        s = ''
        codes.add(r'YYY')
\end{verbatim}







Який рядок треба записати на місце \kw{XXX} ?\par
\vspace{5mm}


Який рядок треба записати на місце \kw{YYY} ?\par


\newpage

{\begin {center}\textbf{ 142}\end {center}}\par
Рядки журналу через двокрапку містять ім'я користувача (складається з латиниці), час події (фор\-мат як у прикладі), тип події (ERROR, WARNING, INFO, STATUS), код події (ціле).

Білі символи навколо роздільників не допускаються.

Приклад такого рядка присвоєно змінній \kw{s}.



\begin{verbatim}
s = 'username:2022-05-05 13-26:ERROR:404'
\end{verbatim}


Нижче наведено код, за виконання якого значенням змінної \kw{out} має стати код події.

Для наведеного прикладу значенням змінної \kw{out} має стати рядок \code{404}



\begin{verbatim}
res = re.fullmatch(r'XXX', s)
s = ''
out = r'YYY'
\end{verbatim}


Код має працювати не тільки для конкретного рядка з прикладу, але й для кожного рядка з журналу.




Який рядок треба записати на місце \kw{XXX} ?\par
\vspace{5mm}


Який рядок треба записати на місце \kw{YYY} ?\par


\newpage

{\begin {center}\textbf{ 143}\end {center}}\par
Рядки відомості через кому містять порядковий номер (ціле), назву предмету (зна\-ків пунктуації не містить), прізвище студента, ім'я студента, оцінку в балах за 100 бальною шкалою.

Білі символи навколо роздільників не допускаються.

Приклад такого рядка присвоєно змінній \kw{s}.



\begin{verbatim}
s = '13,algebra,surname,name,77'
\end{verbatim}


Нижче наведено код, за виконання якого значенням змінної \kw{out} має стати прізвище студента.

Для наведеного прикладу значенням змінної \kw{out} має стати рядок \code{surname}



\begin{verbatim}
res = re.fullmatch(r'XXX', s)
s = ''
out = r'YYY'
\end{verbatim}


Код має працювати не тільки для конкретного рядка з прикладу, але й для кожного рядка з відомості.




Який рядок треба записати на місце \kw{XXX} ?\par
\vspace{5mm}


Який рядок треба записати на місце \kw{YYY} ?\par


\newpage

{\begin {center}\textbf{ 144}\end {center}}\par
Рядки відомості через двокрапку містять порядковий номер (ціле), пріз\-ви\-ще студен\-та, ім'я студента, номер заліковки, оцінку в балах за 100 бальною шкалою.

Білі символи навколо роздільників не допускаються.

Приклад такого рядка присвоєно змінній \kw{s}.



\begin{verbatim}
s = '13:surname:name:123456:77'
\end{verbatim}


Нижче наведено код, за виконання якого для рядка \kw{s} значенням змінної \kw{out} має стати номер заліковки.

Для наведеного прикладу значенням змінної \kw{out} має стати рядок \code{123456} .



\begin{verbatim}
res = re.fullmatch(r'XXX', s)
s = ''
out = r'YYY'
\end{verbatim}


Код має працювати не тільки для конкретного рядка з прикладу, але й для кожного рядка з відомості.




Який рядок треба записати на місце \kw{XXX} ?\par
\vspace{5mm}


Який рядок треба записати на місце \kw{YYY} ?\par


\newpage

{\begin {center}\textbf{ 145}\end {center}}\par
Файл \code{'1.txt'} містить журнал з рядками, в яких через двокрапку записано поряд\-ко\-вий номер запису (ціле), ім'я користувача (складається з латиниці), час події (фор\-мат як у прикладі), тип події (ERROR, WARNING, INFO, STATUS), код події (ціле).

Білі символи навколо роздільників не допускаються.

Приклад рядка журналу



\begin{verbatim}
12:username:2022-05-05 13-26:ERROR:404
\end{verbatim}


Нижче наведено код, який шукає усі події, що відбулись\\ для користувача \code{anonymous}, та зберігає їх типи в змінній \kw{codes}.



\begin{verbatim}
codes = set()
regex = re.compile(r'XXX')
with open('1.txt', encoding='utf8') as f:
    for s in f:
        res = re.fullmatch(regex, s)
        if not res: continue
        s = ''
        codes.add(r'YYY')
\end{verbatim}





Який рядок треба записати на місце \kw{XXX} ?\par
\vspace{5mm}


Який рядок треба записати на місце \kw{YYY} ?\par


\newpage

{\begin {center}\textbf{ 146}\end {center}}\par
Файл \code{'1.txt'} містить відомість з рядками, в яких через кому записано поряд\-ковий номер (ціле), назву предмету (знаків пунктуації не містить), прізвище студента, ім'я студента, номер заліковки, оцінку в балах за 100 бальною шкалою.

Білі символи навколо роздільників не допускаються.

Приклад рядка файлу



\begin{verbatim}
13,algebra,surname,name,123456,77
\end{verbatim}


Нижче наведено код, який шукає усі назви предметів, що зустрічаються у файлі, та зберігає їх у змінній \id{codes}.



\begin{verbatim}
codes = set()
regex = re.compile(r'XXX')
with open('1.txt', encoding='utf8') as f:
    for s in f:
        res = re.fullmatch(regex, s)
        if not res: continue
        s = ''
        codes.add(r'YYY')
\end{verbatim}





Який рядок треба записати на місце \kw{XXX} ?\par
\vspace{5mm}


Який рядок треба записати на місце \kw{YYY} ?\par


\newpage

{\begin {center}\textbf{ 147}\end {center}}\par
Рядки журналу через кому містять тип події (ERROR, WARNING, INFO, STATUS), ім'я користувача (складається з латиниці), час події (фор\-мат як у прикладі), код події (ціле).

Білі символи навколо роздільників не допускаються.

Приклад такого рядка присвоєно змінній \kw{s}.



\begin{verbatim}
s = 'ERROR,username,2022-05-05 13:26,404'
\end{verbatim}


Нижче наведено код, за виконання якого значенням змінної \kw{out} має стати ім'я користувача.

Для наведеного прикладу значенням змінної \kw{out} має стати рядок \code{username}



\begin{verbatim}
res = re.fullmatch(r'XXX', s)
s = ''
out = r'YYY'
\end{verbatim}


Код має працювати не тільки для конкретного рядка з прикладу, але й для кожного рядка з журналу.




Який рядок треба записати на місце \kw{XXX} ?\par
\vspace{5mm}


Який рядок треба записати на місце \kw{YYY} ?\par


\newpage

{\begin {center}\textbf{ 148}\end {center}}\par
Файл \code{'1.txt'} містить журнал з рядками, в яких через кому записано поряд\-ко\-вий номер запису в журналі (ціле), код події (ціле), тип події (ERROR, \\WARNING, INFO, STATUS), ім'я користувача (складається з латиниці), час події (фор\-мат як у прикладі), опис події.

Білі символи навколо роздільників не допускаються.

Приклад рядка файлу



\begin{verbatim}
13,404,ERROR,username,2022-05-05 13:26,something happens
\end{verbatim}


Нижче наведено код, який має залишити у файлі в кожному рядку тільки код події та ім'я користувача. Рядок з прикладу має бути замінений на такий



\begin{verbatim}
404,username
\end{verbatim}




Код:



\begin{verbatim}
regex = re.compile(r'XXX')
with open('1.txt', encoding='utf8') as f:
    lst = f.readlines()
for i in range(len(lst)):
    lst[i] = re.sub(regex, r'YYY', lst[i])
with open('1.txt', 'w', encoding='utf8') as f:
    f.writelines(lst)
\end{verbatim}







Який рядок треба записати на місце \kw{XXX} ?\par
\vspace{5mm}


Який рядок треба записати на місце \kw{YYY} ?\par


\newpage

{\begin {center}\textbf{ 149}\end {center}}\par
Рядки журналу через двокрапку містять тип події (ERROR, WARNING, INFO, STATUS), ім'я користувача (складається з латиниці), час події (фор\-мат як у прикладі), код події (ціле).

Білі символи навколо роздільників не допускаються.

Приклад такого рядка присвоєно змінній \kw{s}.



\begin{verbatim}
s = 'ERROR:username:2022-05-05 13-26:404'
\end{verbatim}


Нижче наведено код, за виконання якого значенням змінної \kw{out} має стати ім'я користувача.

Для наведеного прикладу значенням змінної \kw{out} має стати рядок \code{username}



\begin{verbatim}
res = re.fullmatch(r'XXX', s)
s = ''
out = r'YYY'
\end{verbatim}


Код має працювати не тільки для конкретного рядка з прикладу, але й для кожного рядка з журналу.




Який рядок треба записати на місце \kw{XXX} ?\par
\vspace{5mm}


Який рядок треба записати на місце \kw{YYY} ?\par


\newpage

{\begin {center}\textbf{ 150}\end {center}}\par
Файл \code{'1.txt'} містить журнал з рядками, в яких через кому записано поряд\-ко\-вий номер запису в журналі (ціле), код події (ціле), тип події (ERROR,\\ WARNING, INFO, STATUS), ім'я користувача (складається з латиниці), час події (фор\-мат як у прикладі), опис події.

Білі символи навколо роздільників не допускаються.

Приклад рядка файлу



\begin{verbatim}
13,404,ERROR,username,2022-05-05 13:26,something happens
\end{verbatim}


Нижче наведено код, який має в кожному рядку файли переставити місцями код та тип події.

Рядок з прикладу має бути замінений на такий



\begin{verbatim}
13,ERROR,404,username,2022-05-05 13:26,something happens
\end{verbatim}




Код:



\begin{verbatim}
regex = re.compile(r'XXX')
with open('1.txt', encoding='utf8') as f:
    lst = f.readlines()
for i in range(len(lst)):
    lst[i] = re.sub(regex, r'YYY', lst[i])
with open('1.txt', 'w', encoding='utf8') as f:
    f.writelines(lst)
\end{verbatim}







Який рядок треба записати на місце \kw{XXX} ?\par
\vspace{5mm}


Який рядок треба записати на місце \kw{YYY} ?\par


\newpage

{\begin {center}\textbf{ 151}\end {center}}\par
Файл \code{'1.txt'} містить журнал з рядками, в яких через кому записано поряд\-ко\-вий номер запису в журналі (ціле), тип події (ERROR, WARNING, INFO, STATUS), ім'я користувача (складається з латиниці), час події (фор\-мат як у прикладі), код події (ціле).

Білі символи навколо роздільників не допускаються.

Приклад рядка файлу



\begin{verbatim}
13,ERROR,username,2022-05-05 13:26,404
\end{verbatim}


Нижче наведено код, який шукає усі коди помилок (подія ERROR), що зустріча\-ють\-ся в журналі, та зберігає їх у змінній \kw{codes}.



\begin{verbatim}
codes = set()
regex = re.compile(r'XXX')
with open('1.txt', encoding='utf8') as f:
    for s in f:
        res = re.fullmatch(regex, s)
        if not res: continue
        s = ''
        codes.add(r'YYY')
\end{verbatim}





Який рядок треба записати на місце \kw{XXX} ?\par
\vspace{5mm}


Який рядок треба записати на місце \kw{YYY} ?\par


\newpage

{\begin {center}\textbf{ 152}\end {center}}\par
Файл \code{'1.txt'} містить відомість з рядками, в яких через кому записа\-но поряд\-ковий номер (ціле), назву предмету (знаків пунктуації не містить), прізвище студента, ім'я студента, номер заліковки, оцінку в балах за 100 бальною шкалою.

Білі символи навколо роздільників не допускаються.

Приклад рядка файлу



\begin{verbatim}
13,algebra,surname,name,123456,77
\end{verbatim}


Нижче наведено код, який має залишити у файлі в кожному рядку тільки номер заліковки, назву предмету та оцінку в балах за 100-бальною шкалою (у заданій послі\-дов\-ності).

Рядок з прикладу має бути замінений на такий



\begin{verbatim}
123456,algebra,77
\end{verbatim}




Код



\begin{verbatim}
regex = re.compile(r'XXX')
with open('1.txt', encoding='utf8') as f:
    lst = f.readlines()
for i in range(len(lst)):
    lst[i] = re.sub(regex, r'YYY', lst[i])
with open('1.txt', 'w', encoding='utf8') as f:
    f.writelines(lst)
\end{verbatim}







Який рядок треба записати на місце \kw{XXX} ?\par
\vspace{5mm}


Який рядок треба записати на місце \kw{YYY} ?\par


\newpage

{\begin {center}\textbf{ 153}\end {center}}\par
Файл \code{'1.txt'} містить відомість з рядками, в яких через двокрапку записа\-но поряд\-ковий номер (ціле), назву предмету (знаків пунктуації не містить), прізвище студента, ім'я студента, номер заліковки, оцінку в балах за 100 бальною шкалою.

Білі символи навколо роздільників не допускаються.

Приклад рядка файлу



\begin{verbatim}
13:algebra:surname:name:123456:77
\end{verbatim}


Нижче наведено код, який має залишити у файлі в кожному рядку тільки номер заліковки, назву предмету та оцінку в балах за 100-бальною шкалою (у заданій послі\-дов\-ності).

Рядок з прикладу має бути замінений на такий



\begin{verbatim}
123456:algebra:77
\end{verbatim}




Код:



\begin{verbatim}
regex = re.compile(r'XXX')
with open('1.txt', encoding='utf8') as f:
    lst = f.readlines()
for i in range(len(lst)):
    lst[i] = re.sub(regex, r'YYY', lst[i])
with open('1.txt', 'w', encoding='utf8') as f:
    f.writelines(lst)
\end{verbatim}







Який рядок треба записати на місце \kw{XXX} ?\par
\vspace{5mm}


Який рядок треба записати на місце \kw{YYY} ?\par


\newpage

{\begin {center}\textbf{ 154}\end {center}}\par
Файл \code{'1.txt'} містить відомість з рядками, в яких через двокрапку записа\-но назву предмету (знаків пунктуації не містить), прізвище студента, ім'я студента, номер залі\-ков\-ки, оцінку в балах за 100 бальною шкалою.

Білі символи навколо роздільників не допускаються.

Приклад рядка файлу



\begin{verbatim}
algebra:surname:name:123456:77
\end{verbatim}


Нижче наведено код, який має замінити у файлі усі номери заліковок на рядок \code{***}.



\begin{verbatim}
regex = re.compile(r'XXX')
with open('1.txt', encoding='utf8') as f:
    lst = f.readlines()
for i in range(len(lst)):
    lst[i] = re.sub(regex, r'YYY', lst[i])
with open('1.txt', 'w', encoding='utf8') as f:
    f.writelines(lst)
\end{verbatim}







Який рядок треба записати на місце \kw{XXX} ?\par
\vspace{5mm}


Який рядок треба записати на місце \kw{YYY} ?\par


\newpage

{\begin {center}\textbf{ 155}\end {center}}\par
Файл \code{'1.txt'} містить журнал з рядками, в яких через двокрапку записано поряд\-ко\-вий номер запису в журналі (ціле), код події (ціле), тип події (ERROR, \\WARNING, INFO, STATUS), ім'я користувача (складається з латиниці), час події (фор\-мат як у прикладі), опис події.

Білі символи навколо роздільників не допускаються.

Приклад рядка файлу



\begin{verbatim}
13:404:ERROR:username:2022-05-05 13-26:something happens
\end{verbatim}


Нижче наведено код, який має залишити у файлі в кожному рядку тільки код події та ім'я користувача. Рядок з прикладу має бути замінений на такий



\begin{verbatim}
404:username
\end{verbatim}




Код:



\begin{verbatim}
regex = re.compile(r'XXX')
with open('1.txt', encoding='utf8') as f:
    lst = f.readlines()
for i in range(len(lst)):
    lst[i] = re.sub(regex, r'YYY', lst[i])
with open('1.txt', 'w', encoding='utf8') as f:
    f.writelines(lst)
\end{verbatim}







Який рядок треба записати на місце \kw{XXX} ?\par
\vspace{5mm}


Який рядок треба записати на місце \kw{YYY} ?\par


\newpage

{\begin {center}\textbf{ 156}\end {center}}\par
Файл \code{'1.txt'} містить відомість з рядками, в яких через двокрапку записано поряд\-ковий номер (ціле), назву предмету (знаків пунктуації не містить), прізвище студента, ім'я студента, номер заліковки, оцінку в балах за 100 бальною шкалою.

Білі символи навколо роздільників не допускаються.

Приклад рядка файлу



\begin{verbatim}
13:algebra:surname:name:123456:77
\end{verbatim}


Нижче наведено код, який шукає усі предмети, який складав студент із заліковкою \code{777}, та зберігає їх типи в змінній \kw{codes}.



\begin{verbatim}
codes = set()
regex = re.compile(r'XXX')
with open('1.txt', encoding='utf8') as f:
    for s in f:
        res = re.fullmatch(regex, s)
        if not res: continue
        s = ''
        codes.add(r'YYY')
\end{verbatim}





Який рядок треба записати на місце \kw{XXX} ?\par
\vspace{5mm}


Який рядок треба записати на місце \kw{YYY} ?\par


\newpage

{\begin {center}\textbf{ 157}\end {center}}\par
Файл \code{'1.txt'} містить журнал з рядками, в яких через двокрапку записано поряд\-ко\-вий номер запису в журналі (ціле), тип події (ERROR, WARNING, INFO, STATUS), ім'я користувача (складається з латиниці), час події (фор\-мат як у прикладі), опис події.

Білі символи навколо роздільників не допускаються.

Приклад рядка файлу



\begin{verbatim}
13:ERROR:username:2022-05-05 13-26:something happens
\end{verbatim}


Нижче наведено код, який має залишити у файлі в кожному рядку тільки тип події та її опис.

Рядок з прикладу має бути замінений на такий



\begin{verbatim}
ERROR:something happens
\end{verbatim}




Код:



\begin{verbatim}
regex = re.compile(r'XXX')
with open('1.txt', encoding='utf8') as f:
    lst = f.readlines()
for i in range(len(lst)):
    lst[i] = re.sub(regex, r'YYY', lst[i])
with open('1.txt', 'w', encoding='utf8') as f:
    f.writelines(lst)
\end{verbatim}









Який рядок треба записати на місце \kw{XXX} ?\par
\vspace{5mm}


Який рядок треба записати на місце \kw{YYY} ?\par


\newpage

{\begin {center}\textbf{ 158}\end {center}}\par
Файл \code{'1.txt'} містить журнал з рядками, в яких через кому записано поряд\-ко\-вий номер запису в журналі (ціле), тип події (ERROR, WARNING, INFO, STATUS), ім'я користувача (складається з латиниці), час події (фор\-мат як у прикладі), опис події.

Білі символи навколо роздільників не допускаються.

Приклад рядка файлу



\begin{verbatim}
13,ERROR,username,2022-05-05 13:26,something happens
\end{verbatim}


Нижче наведено код, який має залишити у файлі в кожному рядку тільки тип події та її опис.

Рядок з прикладу має бути замінений на такий



\begin{verbatim}
ERROR,something happens
\end{verbatim}




Код:



\begin{verbatim}
regex = re.compile(r'XXX')
with open('1.txt', encoding='utf8') as f:
    lst = f.readlines()
for i in range(len(lst)):
    lst[i] = re.sub(regex, r'YYY', lst[i])
with open('1.txt', 'w', encoding='utf8') as f:
    f.writelines(lst)
\end{verbatim}









Який рядок треба записати на місце \kw{XXX} ?\par
\vspace{5mm}


Який рядок треба записати на місце \kw{YYY} ?\par


\newpage

{\begin {center}\textbf{ 159}\end {center}}\par
Файл \code{'1.txt'} містить журнал з рядками, в яких через кому записано поряд\-ко\-вий номер запису (ціле), ім'я користувача (складається з латиниці), час події (фор\-мат як у прикладі), тип події (ERROR, WARNING, INFO, STATUS), код події (ціле).

Білі символи навколо роздільників не допускаються.

Приклад рядка журналу



\begin{verbatim}
12,username,2022-05-05 13:26,ERROR,404
\end{verbatim}


Нижче наведено код, який шукає усі події, що відбулись\\ для користувача \code{anonymous}, та зберігає їх типи в змінній \kw{codes}.



\begin{verbatim}
codes = set()
regex = re.compile(r'XXX')
with open('1.txt', encoding='utf8') as f:
    for s in f:
        res = re.fullmatch(regex, s)
        if not res: continue
        s = ''
        codes.add(r'YYY')
\end{verbatim}





Який рядок треба записати на місце \kw{XXX} ?\par
\vspace{5mm}


Який рядок треба записати на місце \kw{YYY} ?\par


\newpage

{\begin {center}\textbf{ 160}\end {center}}\par
Файл \code{'1.txt'} містить журнал з рядками, в яких через двокрапку записано поряд\-ко\-вий номер запису в журналі (ціле), тип події (ERROR, WARNING, INFO, STATUS), ім'я користувача (складається з латиниці), час події (фор\-мат як у прикладі), код події (ціле).

Білі символи навколо роздільників не допускаються.

Приклад рядка файлу



\begin{verbatim}
13:ERROR:username:2022-05-05 13-26:404
\end{verbatim}


Нижче наведено код, який шукає усі коди помилок (подія ERROR), що зустріча\-ють\-ся в журналі, та зберігає їх у змінній \kw{codes}.



\begin{verbatim}
codes = set()
regex = re.compile(r'XXX')
with open('1.txt', encoding='utf8') as f:
    for s in f:
        res = re.fullmatch(regex, s)
        if not res: continue
        s = ''
        codes.add(r'YYY')
\end{verbatim}





Який рядок треба записати на місце \kw{XXX} ?\par
\vspace{5mm}


Який рядок треба записати на місце \kw{YYY} ?\par


\newpage

{\begin {center}\textbf{ 161}\end {center}}\par
Файл \code{'1.txt'} містить відомість з рядками, в яких через кому записа\-но назву предмету (знаків пунктуації не містить), прізвище студента, ім'я студента, номер залі\-ков\-ки, оцінку в балах за 100 бальною шкалою.

Білі символи навколо роздільників не допускаються.

Приклад рядка файлу



\begin{verbatim}
algebra,surname,name,123456,77
\end{verbatim}


Нижче наведено код, який має замінити у файлі усі номери заліковок на рядок \code{***}.



\begin{verbatim}
regex = re.compile(r'XXX')
with open('1.txt', encoding='utf8') as f:
    lst = f.readlines()
for i in range(len(lst)):
    lst[i] = re.sub(regex, r'YYY', lst[i])
with open('1.txt', 'w', encoding='utf8') as f:
    f.writelines(lst)
\end{verbatim}







Який рядок треба записати на місце \kw{XXX} ?\par
\vspace{5mm}


Який рядок треба записати на місце \kw{YYY} ?\par


\newpage

{\begin {center}\textbf{ 162}\end {center}}\par
Файл \code{'1.txt'} містить відомість з рядками, в яких через двокрапку записано поряд\-ковий номер (ціле), назву предмету (знаків пунктуації не містить), прізвище студента, ім'я студента, номер заліковки, оцінку в балах за 100 бальною шкалою.

Білі символи навколо роздільників не допускаються.

Приклад рядка файлу



\begin{verbatim}
13:algebra:surname:name:123456:77
\end{verbatim}


Нижче наведено код, який має в кожному рядку файлу переставити місцями назву предмету та номер заліковки.

Рядок з прикладу має бути замінений на такий



\begin{verbatim}
13:123456:surname:name:algebra:77
\end{verbatim}




Код:



\begin{verbatim}
regex = re.compile(r'XXX')
with open('1.txt', encoding='utf8') as f:
    lst = f.readlines()
for i in range(len(lst)):
    lst[i] = re.sub(regex, r'YYY', lst[i])
with open('1.txt', 'w', encoding='utf8') as f:
    f.writelines(lst)
\end{verbatim}







Який рядок треба записати на місце \kw{XXX} ?\par
\vspace{5mm}


Який рядок треба записати на місце \kw{YYY} ?\par


\newpage

{\begin {center}\textbf{ 163}\end {center}}\par
Файл \code{'1.txt'} містить журнал з рядками, в яких через двокрапку записано поряд\-ко\-вий номер запису (ціле), ім'я користувача (складається з латиниці), час події (фор\-мат як у прикладі), тип події (ERROR, WARNING, INFO, STATUS), код події (ціле).

Білі символи навколо роздільників не допускаються.

Приклад рядка журналу



\begin{verbatim}
12:username:2022-05-05 13-26:ERROR:404
\end{verbatim}


Нижче наведено код, який шукає усі попередження (подія WARNING), що зустрі\-ча\-ються в журналі і відбулися під час обідньої перерви деканату (з 13:00 по 13:59 включно), та зберігає їх коди в змінній \kw{codes}.



\begin{verbatim}
codes = set()
regex = re.compile(r'XXX')
with open('1.txt', encoding='utf8') as f:
    for s in f:
        res = re.fullmatch(regex, s)
        if not res: continue
        s = ''
        codes.add(r'YYY')
\end{verbatim}







Який рядок треба записати на місце \kw{XXX} ?\par
\vspace{5mm}


Який рядок треба записати на місце \kw{YYY} ?\par


\newpage

{\begin {center}\textbf{ 164}\end {center}}\par
Файл \code{'1.txt'} містить журнал з рядками, в яких через двокрапку записано поряд\-ко\-вий номер запису в журналі (ціле), код події (ціле), тип події (ERROR,\\ WARNING, INFO, STATUS), ім'я користувача (складається з латиниці), час події (фор\-мат як у прикладі), опис події.

Білі символи навколо роздільників не допускаються.

Приклад рядка файлу



\begin{verbatim}
13:404:ERROR:username:2022-05-05 13:26:something happens
\end{verbatim}


Нижче наведено код, який має в кожному рядку файли переставити місцями код та тип події.

Рядок з прикладу має бути замінений на такий



\begin{verbatim}
13:ERROR:404:username:2022-05-05 13:26:something happens
\end{verbatim}




Код:



\begin{verbatim}
regex = re.compile(r'XXX')
with open('1.txt', encoding='utf8') as f:
    lst = f.readlines()
for i in range(len(lst)):
    lst[i] = re.sub(regex, r'YYY', lst[i])
with open('1.txt', 'w', encoding='utf8') as f:
    f.writelines(lst)
\end{verbatim}







Який рядок треба записати на місце \kw{XXX} ?\par
\vspace{5mm}


Який рядок треба записати на місце \kw{YYY} ?\par


\newpage

{\begin {center}\textbf{ 165}\end {center}}\par
Рядки відомості через кому містять порядковий номер (ціле), пріз\-ви\-ще студен\-та, ім'я студента, номер заліковки, оцінку в балах за 100 бальною шкалою.

Білі символи навколо роздільників не допускаються.

Приклад такого рядка присвоєно змінній \kw{s}.



\begin{verbatim}
s = '13,surname,name,123456,77'
\end{verbatim}


Нижче наведено код, за виконання якого для рядка \kw{s} значенням змінної \kw{out} має стати номер заліковки.

Для наведеного прикладу значенням змінної \kw{out} має стати рядок \code{123456} .



\begin{verbatim}
res = re.fullmatch(r'XXX', s)
s = ''
out = r'YYY'
\end{verbatim}


Код має працювати не тільки для конкретного рядка з прикладу, але й для кожного рядка з відомості.




Який рядок треба записати на місце \kw{XXX} ?\par
\vspace{5mm}


Який рядок треба записати на місце \kw{YYY} ?\par


\newpage

{\begin {center}\textbf{ 166}\end {center}}\par
Файл \code{'1.txt'} містить журнал з рядками, в яких через двокрапку записано поряд\-ко\-вий номер запису в журналі (ціле), тип події (ERROR, WARNING, INFO, STATUS), ім'я користувача (складається з латиниці), час події (фор\-мат як у прикладі), код події (ціле).

Білі символи навколо роздільників не допускаються.

Приклад рядка файлу



\begin{verbatim}
13:ERROR:username:2022-05-05 13-26:404
\end{verbatim}


Нижче наведено код, який має замінити у файлі імена користувачів на рядок \code{unknown}.



\begin{verbatim}
regex = re.compile(r'XXX')
with open('1.txt', encoding='utf8') as f:
    lst = f.readlines()
for i in range(len(lst)):
    lst[i] = re.sub(regex, r'YYY', lst[i])
with open('1.txt', 'w', encoding='utf8') as f:
    f.writelines(lst)
\end{verbatim}







Який рядок треба записати на місце \kw{XXX} ?\par
\vspace{5mm}


Який рядок треба записати на місце \kw{YYY} ?\par


\newpage

{\begin {center}\textbf{ 167}\end {center}}\par
Файл \code{'1.txt'} містить відомість з рядками, в яких через двокрапку записано поряд\-ковий номер (ціле), назву предмету (знаків пунктуації не містить), прізвище студента, ім'я студента, номер заліковки, оцінку в балах за 100 бальною шкалою.

Білі символи навколо роздільників не допускаються.

Приклад рядка файлу



\begin{verbatim}
13:algebra:surname:name:123456:77
\end{verbatim}


Нижче наведено код, який шукає усі назви предметів, що зустрічаються у файлі, та зберігає їх у змінній \id{codes}.



\begin{verbatim}
codes = set()
regex = re.compile(r'XXX')
with open('1.txt', encoding='utf8') as f:
    for s in f:
        res = re.fullmatch(regex, s)
        if not res: continue
        s = ''
        codes.add(r'YYY')
\end{verbatim}





Який рядок треба записати на місце \kw{XXX} ?\par
\vspace{5mm}


Який рядок треба записати на місце \kw{YYY} ?\par


\newpage

{\begin {center}\textbf{ 168}\end {center}}\par
Файл \code{'1.txt'} містить відомість з рядками, в яких через кому записано поряд\-ковий номер (ціле), назву предмету (знаків пунктуації не містить), прізвище студента, ім'я студента, номер заліковки, оцінку в балах за 100 бальною шкалою.

Білі символи навколо роздільників не допускаються.

Приклад рядка файлу



\begin{verbatim}
13,algebra,surname,name,123456,77
\end{verbatim}


Нижче наведено код, який має в кожному рядку файлу переставити місцями назву предмету та номер заліковки.

Рядок з прикладу має бути замінений на такий



\begin{verbatim}
13,123456,surname,name,algebra,77
\end{verbatim}




Код



\begin{verbatim}
regex = re.compile(r'XXX')
with open('1.txt', encoding='utf8') as f:
    lst = f.readlines()
for i in range(len(lst)):
    lst[i] = re.sub(regex, r'YYY', lst[i])
with open('1.txt', 'w', encoding='utf8') as f:
    f.writelines(lst)
\end{verbatim}







Який рядок треба записати на місце \kw{XXX} ?\par
\vspace{5mm}


Який рядок треба записати на місце \kw{YYY} ?\par


\newpage

{\begin {center}\textbf{ 169}\end {center}}\par
Файл \code{'1.txt'} містить відомість з рядками, в яких через кому записано поряд\-ковий номер (ціле), назву предмету (знаків пунктуації не містить), прізвище студента, ім'я студента, номер заліковки, оцінку в балах за 100 бальною шкалою.

Білі символи навколо роздільників не допускаються.

Приклад рядка файлу



\begin{verbatim}
13,algebra,surname,name,123456,77
\end{verbatim}


Нижче наведено код, який має залишити у файлі в кожному рядку тільки номер заліковки та предмет.

Рядок з прикладу має бути замінений на такий



\begin{verbatim}
123456,algebra
\end{verbatim}




Код



\begin{verbatim}
regex = re.compile(r'XXX')
with open('1.txt', encoding='utf8') as f:
    lst = f.readlines()
for i in range(len(lst)):
    lst[i] = re.sub(regex, r'YYY', lst[i])
with open('1.txt', 'w', encoding='utf8') as f:
    f.writelines(lst)
\end{verbatim}









Який рядок треба записати на місце \kw{XXX} ?\par
\vspace{5mm}


Який рядок треба записати на місце \kw{YYY} ?\par


\newpage

{\begin {center}\textbf{ 170}\end {center}}\par
Файл \code{'1.txt'} містить відомість з рядками, в яких через двокрапку записано поряд\-ковий номер (ціле), назву предмету (знаків пунктуації не містить), прізвище студента, ім'я студента, номер заліковки, оцінку в балах за 100 бальною шкалою.

Білі символи навколо роздільників не допускаються.

Приклад рядка файлу



\begin{verbatim}
13:algebra:surname:name:123456:77
\end{verbatim}


Нижче наведено код, який має залишити у файлі в кожному рядку тільки номер заліковки та предмет.

Рядок з прикладу має бути замінений на такий



\begin{verbatim}
123456:algebra
\end{verbatim}




Код:



\begin{verbatim}
regex = re.compile(r'XXX')
with open('1.txt', encoding='utf8') as f:
    lst = f.readlines()
for i in range(len(lst)):
    lst[i] = re.sub(regex, r'YYY', lst[i])
with open('1.txt', 'w', encoding='utf8') as f:
    f.writelines(lst)
\end{verbatim}









Який рядок треба записати на місце \kw{XXX} ?\par
\vspace{5mm}


Який рядок треба записати на місце \kw{YYY} ?\par


\newpage

{\begin {center}\textbf{ 171}\end {center}}\par
Файл \code{'1.txt'} містить журнал з рядками, в яких через двокрапку записано поряд\-ко\-вий номер запису в журналі (ціле), тип події (ERROR, WARNING, INFO, STATUS), ім'я користувача (складається з латиниці), час події (фор\-мат як у прикладі), код події (ціле).

Білі символи навколо роздільників не допускаються.

Приклад рядка файлу



\begin{verbatim}
13:ERROR:username:2022-05-05 13-26:404
\end{verbatim}


Нижче наведено код, який шукає усі коди помилок (подія ERROR), що зустріча\-ють\-ся в журналі, та зберігає їх у змінній \kw{codes}.



\begin{verbatim}
codes = set()
regex = re.compile(r'XXX')
with open('1.txt', encoding='utf8') as f:
    for s in f:
        res = re.fullmatch(regex, s)
        if not res: continue
        s = ''
        codes.add(r'YYY')
\end{verbatim}





Який рядок треба записати на місце \kw{XXX} ?\par
\vspace{5mm}


Який рядок треба записати на місце \kw{YYY} ?\par


\newpage

{\begin {center}\textbf{ 172}\end {center}}\par
Рядки журналу через двокрапку містять ім'я користувача (складається з латиниці), час події (фор\-мат як у прикладі), тип події (ERROR, WARNING, INFO, STATUS), код події (ціле).

Білі символи навколо роздільників не допускаються.

Приклад такого рядка присвоєно змінній \kw{s}.



\begin{verbatim}
s = 'username:2022-05-05 13-26:ERROR:404'
\end{verbatim}


Нижче наведено код, за виконання якого значенням змінної \kw{out} має стати код події.

Для наведеного прикладу значенням змінної \kw{out} має стати рядок \code{404}



\begin{verbatim}
res = re.fullmatch(r'XXX', s)
s = ''
out = r'YYY'
\end{verbatim}


Код має працювати не тільки для конкретного рядка з прикладу, але й для кожного рядка з журналу.




Який рядок треба записати на місце \kw{XXX} ?\par
\vspace{5mm}


Який рядок треба записати на місце \kw{YYY} ?\par


\newpage

{\begin {center}\textbf{ 173}\end {center}}\par
Рядки журналу через кому містять ім'я користувача (складається з латиниці), час події (фор\-мат як у прикладі), тип події (ERROR, WARNING, INFO, STATUS), код події (ціле).

Білі символи навколо роздільників не допускаються.

Приклад такого рядка присвоєно змінній \kw{s}.



\begin{verbatim}
s = 'username,2022-05-05 13:26,ERROR,404'
\end{verbatim}


Нижче наведено код, за виконання якого значенням змінної \kw{out} має стати код події.

Для наведеного прикладу значенням змінної \kw{out} має стати рядок \code{404}



\begin{verbatim}
res = re.fullmatch(r'XXX', s)
s = ''
out = r'YYY'
\end{verbatim}


Код має працювати не тільки для конкретного рядка з прикладу, але й для кожного рядка з журналу.




Який рядок треба записати на місце \kw{XXX} ?\par
\vspace{5mm}


Який рядок треба записати на місце \kw{YYY} ?\par


\newpage

{\begin {center}\textbf{ 174}\end {center}}\par
Файл \code{'1.txt'} містить журнал з рядками, в яких через кому записано поряд\-ко\-вий номер запису в журналі (ціле), код події (ціле), тип події (ERROR,\\ WARNING, INFO, STATUS), ім'я користувача (складається з латиниці), час події (фор\-мат як у прикладі), опис події.

Білі символи навколо роздільників не допускаються.

Приклад рядка файлу



\begin{verbatim}
13,404,ERROR,username,2022-05-05 13:26,something happens
\end{verbatim}


Нижче наведено код, який має в кожному рядку файли переставити місцями код та тип події.

Рядок з прикладу має бути замінений на такий



\begin{verbatim}
13,ERROR,404,username,2022-05-05 13:26,something happens
\end{verbatim}




Код:



\begin{verbatim}
regex = re.compile(r'XXX')
with open('1.txt', encoding='utf8') as f:
    lst = f.readlines()
for i in range(len(lst)):
    lst[i] = re.sub(regex, r'YYY', lst[i])
with open('1.txt', 'w', encoding='utf8') as f:
    f.writelines(lst)
\end{verbatim}







Який рядок треба записати на місце \kw{XXX} ?\par
\vspace{5mm}


Який рядок треба записати на місце \kw{YYY} ?\par


\newpage

{\begin {center}\textbf{ 175}\end {center}}\par
Файл \code{'1.txt'} містить відомість з рядками, в яких через кому записано поряд\-ковий номер (ціле), назву предмету (знаків пунктуації не містить), прізвище студента, ім'я студента, номер заліковки, оцінку в балах за 100 бальною шкалою.

Білі символи навколо роздільників не допускаються.

Приклад рядка файлу



\begin{verbatim}
13,algebra,surname,name,123456,77
\end{verbatim}


Нижче наведено код, який шукає усі предмети, який складав студент із заліковкою \code{777}, та зберігає їх типи в змінній \kw{codes}.



\begin{verbatim}
codes = set()
regex = re.compile(r'XXX')
with open('1.txt', encoding='utf8') as f:
    for s in f:
        res = re.fullmatch(regex, s)
        if not res: continue
        s = ''
        codes.add(r'YYY')
\end{verbatim}





Який рядок треба записати на місце \kw{XXX} ?\par
\vspace{5mm}


Який рядок треба записати на місце \kw{YYY} ?\par


\newpage

{\begin {center}\textbf{ 176}\end {center}}\par
Файл \code{'1.txt'} містить журнал з рядками, в яких через кому записано поряд\-ко\-вий номер запису в журналі (ціле), тип події (ERROR, WARNING, INFO, STATUS), ім'я користувача (складається з латиниці), час події (фор\-мат як у прикладі), опис події.

Білі символи навколо роздільників не допускаються.

Приклад рядка файлу



\begin{verbatim}
13,ERROR,username,2022-05-05 13:26,something happens
\end{verbatim}


Нижче наведено код, який має залишити у файлі в кожному рядку тільки тип події та її опис.

Рядок з прикладу має бути замінений на такий



\begin{verbatim}
ERROR,something happens
\end{verbatim}




Код:



\begin{verbatim}
regex = re.compile(r'XXX')
with open('1.txt', encoding='utf8') as f:
    lst = f.readlines()
for i in range(len(lst)):
    lst[i] = re.sub(regex, r'YYY', lst[i])
with open('1.txt', 'w', encoding='utf8') as f:
    f.writelines(lst)
\end{verbatim}









Який рядок треба записати на місце \kw{XXX} ?\par
\vspace{5mm}


Який рядок треба записати на місце \kw{YYY} ?\par


\newpage

{\begin {center}\textbf{ 177}\end {center}}\par
Файл \code{'1.txt'} містить журнал з рядками, в яких через кому записано поряд\-ко\-вий номер запису в журналі (ціле), тип події (ERROR, WARNING, INFO, STATUS), ім'я користувача (складається з латиниці), час події (фор\-мат як у прикладі), код події (ціле).

Білі символи навколо роздільників не допускаються.

Приклад рядка файлу



\begin{verbatim}
13,ERROR,username,2022-05-05 13:26,404
\end{verbatim}


Нижче наведено код, який має замінити у файлі імена користувачів на рядок \code{unknown}.



\begin{verbatim}
regex = re.compile(r'XXX')
with open('1.txt', encoding='utf8') as f:
    lst = f.readlines()
for i in range(len(lst)):
    lst[i] = re.sub(regex, r'YYY', lst[i])
with open('1.txt', 'w', encoding='utf8') as f:
    f.writelines(lst)
\end{verbatim}







Який рядок треба записати на місце \kw{XXX} ?\par
\vspace{5mm}


Який рядок треба записати на місце \kw{YYY} ?\par


\newpage

{\begin {center}\textbf{ 178}\end {center}}\par
Файл \code{'1.txt'} містить журнал з рядками, в яких через кому записано поряд\-ко\-вий номер запису (ціле), ім'я користувача (складається з латиниці), час події (фор\-мат як у прикладі), тип події (ERROR, WARNING, INFO, STATUS), код події (ціле).

Білі символи навколо роздільників не допускаються.

Приклад рядка журналу



\begin{verbatim}
12,username,2022-05-05 13:26,ERROR,404
\end{verbatim}


Нижче наведено код, який шукає усі попередження (подія WARNING), що зустрі\-ча\-ються в журналі і відбулися під час обідньої перерви деканату (з 13,00 по 13,59 включно), та зберігає їх коди в змінній \kw{codes}.



\begin{verbatim}
codes = set()
regex = re.compile(r'XXX')
with open('1.txt', encoding='utf8') as f:
    for s in f:
        res = re.fullmatch(regex, s)
        if not res: continue
        s = ''
        codes.add(r'YYY')
\end{verbatim}







Який рядок треба записати на місце \kw{XXX} ?\par
\vspace{5mm}


Який рядок треба записати на місце \kw{YYY} ?\par


\newpage

{\begin {center}\textbf{ 179}\end {center}}\par
Файл \code{'1.txt'} містить журнал з рядками, в яких через кому записано поряд\-ко\-вий номер запису в журналі (ціле), код події (ціле), тип події (ERROR, \\WARNING, INFO, STATUS), ім'я користувача (складається з латиниці), час події (фор\-мат як у прикладі), опис події.

Білі символи навколо роздільників не допускаються.

Приклад рядка файлу



\begin{verbatim}
13,404,ERROR,username,2022-05-05 13:26,something happens
\end{verbatim}


Нижче наведено код, який має залишити у файлі в кожному рядку тільки код події та ім'я користувача. Рядок з прикладу має бути замінений на такий



\begin{verbatim}
404,username
\end{verbatim}




Код:



\begin{verbatim}
regex = re.compile(r'XXX')
with open('1.txt', encoding='utf8') as f:
    lst = f.readlines()
for i in range(len(lst)):
    lst[i] = re.sub(regex, r'YYY', lst[i])
with open('1.txt', 'w', encoding='utf8') as f:
    f.writelines(lst)
\end{verbatim}







Який рядок треба записати на місце \kw{XXX} ?\par
\vspace{5mm}


Який рядок треба записати на місце \kw{YYY} ?\par


\newpage

{\begin {center}\textbf{ 180}\end {center}}\par
Файл \code{'1.txt'} містить журнал з рядками, в яких через двокрапку записано поряд\-ко\-вий номер запису в журналі (ціле), код події (ціле), тип події (ERROR, \\WARNING, INFO, STATUS), ім'я користувача (складається з латиниці), час події (фор\-мат як у прикладі), опис події.

Білі символи навколо роздільників не допускаються.

Приклад рядка файлу



\begin{verbatim}
13:404:ERROR:username:2022-05-05 13-26:something happens
\end{verbatim}


Нижче наведено код, який має залишити у файлі в кожному рядку тільки код події та ім'я користувача. Рядок з прикладу має бути замінений на такий



\begin{verbatim}
404:username
\end{verbatim}




Код:



\begin{verbatim}
regex = re.compile(r'XXX')
with open('1.txt', encoding='utf8') as f:
    lst = f.readlines()
for i in range(len(lst)):
    lst[i] = re.sub(regex, r'YYY', lst[i])
with open('1.txt', 'w', encoding='utf8') as f:
    f.writelines(lst)
\end{verbatim}







Який рядок треба записати на місце \kw{XXX} ?\par
\vspace{5mm}


Який рядок треба записати на місце \kw{YYY} ?\par


\newpage

{\begin {center}\textbf{ 181}\end {center}}\par
Рядки журналу через кому містять тип події (ERROR, WARNING, INFO, STATUS), ім'я користувача (складається з латиниці), час події (фор\-мат як у прикладі), код події (ціле).

Білі символи навколо роздільників не допускаються.

Приклад такого рядка присвоєно змінній \kw{s}.



\begin{verbatim}
s = 'ERROR,username,2022-05-05 13:26,404'
\end{verbatim}


Нижче наведено код, за виконання якого значенням змінної \kw{out} має стати ім'я користувача.

Для наведеного прикладу значенням змінної \kw{out} має стати рядок \code{username}



\begin{verbatim}
res = re.fullmatch(r'XXX', s)
s = ''
out = r'YYY'
\end{verbatim}


Код має працювати не тільки для конкретного рядка з прикладу, але й для кожного рядка з журналу.




Який рядок треба записати на місце \kw{XXX} ?\par
\vspace{5mm}


Який рядок треба записати на місце \kw{YYY} ?\par


\newpage

{\begin {center}\textbf{ 182}\end {center}}\par
Файл \code{'1.txt'} містить відомість з рядками, в яких через двокрапку записа\-но поряд\-ковий номер (ціле), назву предмету (знаків пунктуації не містить), прізвище студента, ім'я студента, номер заліковки, оцінку в балах за 100 бальною шкалою.

Білі символи навколо роздільників не допускаються.

Приклад рядка файлу



\begin{verbatim}
13:algebra:surname:name:123456:77
\end{verbatim}


Нижче наведено код, який має залишити у файлі в кожному рядку тільки номер заліковки, назву предмету та оцінку в балах за 100-бальною шкалою (у заданій послі\-дов\-ності).

Рядок з прикладу має бути замінений на такий



\begin{verbatim}
123456:algebra:77
\end{verbatim}




Код:



\begin{verbatim}
regex = re.compile(r'XXX')
with open('1.txt', encoding='utf8') as f:
    lst = f.readlines()
for i in range(len(lst)):
    lst[i] = re.sub(regex, r'YYY', lst[i])
with open('1.txt', 'w', encoding='utf8') as f:
    f.writelines(lst)
\end{verbatim}







Який рядок треба записати на місце \kw{XXX} ?\par
\vspace{5mm}


Який рядок треба записати на місце \kw{YYY} ?\par


\newpage

{\begin {center}\textbf{ 183}\end {center}}\par
Файл \code{'1.txt'} містить журнал з рядками, в яких через двокрапку записано поряд\-ко\-вий номер запису в журналі (ціле), тип події (ERROR, WARNING, INFO, STATUS), ім'я користувача (складається з латиниці), час події (фор\-мат як у прикладі), код події (ціле).

Білі символи навколо роздільників не допускаються.

Приклад рядка файлу



\begin{verbatim}
13:ERROR:username:2022-05-05 13-26:404
\end{verbatim}


Нижче наведено код, який має замінити у файлі імена користувачів на рядок \code{unknown}.



\begin{verbatim}
regex = re.compile(r'XXX')
with open('1.txt', encoding='utf8') as f:
    lst = f.readlines()
for i in range(len(lst)):
    lst[i] = re.sub(regex, r'YYY', lst[i])
with open('1.txt', 'w', encoding='utf8') as f:
    f.writelines(lst)
\end{verbatim}







Який рядок треба записати на місце \kw{XXX} ?\par
\vspace{5mm}


Який рядок треба записати на місце \kw{YYY} ?\par


\newpage

{\begin {center}\textbf{ 184}\end {center}}\par
Файл \code{'1.txt'} містить відомість з рядками, в яких через двокрапку записано поряд\-ковий номер (ціле), назву предмету (знаків пунктуації не містить), прізвище студента, ім'я студента, номер заліковки, оцінку в балах за 100 бальною шкалою.

Білі символи навколо роздільників не допускаються.

Приклад рядка файлу



\begin{verbatim}
13:algebra:surname:name:123456:77
\end{verbatim}


Нижче наведено код, який шукає усі назви предметів, що зустрічаються у файлі, та зберігає їх у змінній \id{codes}.



\begin{verbatim}
codes = set()
regex = re.compile(r'XXX')
with open('1.txt', encoding='utf8') as f:
    for s in f:
        res = re.fullmatch(regex, s)
        if not res: continue
        s = ''
        codes.add(r'YYY')
\end{verbatim}





Який рядок треба записати на місце \kw{XXX} ?\par
\vspace{5mm}


Який рядок треба записати на місце \kw{YYY} ?\par


\newpage

{\begin {center}\textbf{ 185}\end {center}}\par
Файл \code{'1.txt'} містить відомість з рядками, в яких через двокрапку записа\-но назву предмету (знаків пунктуації не містить), прізвище студента, ім'я студента, номер залі\-ков\-ки, оцінку в балах за 100 бальною шкалою.

Білі символи навколо роздільників не допускаються.

Приклад рядка файлу



\begin{verbatim}
algebra:surname:name:123456:77
\end{verbatim}


Нижче наведено код, який має замінити у файлі усі номери заліковок на рядок \code{***}.



\begin{verbatim}
regex = re.compile(r'XXX')
with open('1.txt', encoding='utf8') as f:
    lst = f.readlines()
for i in range(len(lst)):
    lst[i] = re.sub(regex, r'YYY', lst[i])
with open('1.txt', 'w', encoding='utf8') as f:
    f.writelines(lst)
\end{verbatim}







Який рядок треба записати на місце \kw{XXX} ?\par
\vspace{5mm}


Який рядок треба записати на місце \kw{YYY} ?\par


\newpage

{\begin {center}\textbf{ 186}\end {center}}\par
Файл \code{'1.txt'} містить відомість з рядками, в яких через двокрапку записано поряд\-ковий номер (ціле), назву предмету (знаків пунктуації не містить), прізвище студента, ім'я студента, номер заліковки, оцінку в балах за 100 бальною шкалою.

Білі символи навколо роздільників не допускаються.

Приклад рядка файлу



\begin{verbatim}
13:algebra:surname:name:123456:77
\end{verbatim}


Нижче наведено код, який шукає усі предмети, який складав студент із заліковкою \code{777}, та зберігає їх типи в змінній \kw{codes}.



\begin{verbatim}
codes = set()
regex = re.compile(r'XXX')
with open('1.txt', encoding='utf8') as f:
    for s in f:
        res = re.fullmatch(regex, s)
        if not res: continue
        s = ''
        codes.add(r'YYY')
\end{verbatim}





Який рядок треба записати на місце \kw{XXX} ?\par
\vspace{5mm}


Який рядок треба записати на місце \kw{YYY} ?\par


\newpage

{\begin {center}\textbf{ 187}\end {center}}\par
Файл \code{'1.txt'} містить журнал з рядками, в яких через двокрапку записано поряд\-ко\-вий номер запису в журналі (ціле), код події (ціле), тип події (ERROR,\\ WARNING, INFO, STATUS), ім'я користувача (складається з латиниці), час події (фор\-мат як у прикладі), опис події.

Білі символи навколо роздільників не допускаються.

Приклад рядка файлу



\begin{verbatim}
13:404:ERROR:username:2022-05-05 13:26:something happens
\end{verbatim}


Нижче наведено код, який має в кожному рядку файли переставити місцями код та тип події.

Рядок з прикладу має бути замінений на такий



\begin{verbatim}
13:ERROR:404:username:2022-05-05 13:26:something happens
\end{verbatim}




Код:



\begin{verbatim}
regex = re.compile(r'XXX')
with open('1.txt', encoding='utf8') as f:
    lst = f.readlines()
for i in range(len(lst)):
    lst[i] = re.sub(regex, r'YYY', lst[i])
with open('1.txt', 'w', encoding='utf8') as f:
    f.writelines(lst)
\end{verbatim}







Який рядок треба записати на місце \kw{XXX} ?\par
\vspace{5mm}


Який рядок треба записати на місце \kw{YYY} ?\par


\newpage

{\begin {center}\textbf{ 188}\end {center}}\par
Файл \code{'1.txt'} містить журнал з рядками, в яких через двокрапку записано поряд\-ко\-вий номер запису (ціле), ім'я користувача (складається з латиниці), час події (фор\-мат як у прикладі), тип події (ERROR, WARNING, INFO, STATUS), код події (ціле).

Білі символи навколо роздільників не допускаються.

Приклад рядка журналу



\begin{verbatim}
12:username:2022-05-05 13-26:ERROR:404
\end{verbatim}


Нижче наведено код, який шукає усі події, що відбулись\\ для користувача \code{anonymous}, та зберігає їх типи в змінній \kw{codes}.



\begin{verbatim}
codes = set()
regex = re.compile(r'XXX')
with open('1.txt', encoding='utf8') as f:
    for s in f:
        res = re.fullmatch(regex, s)
        if not res: continue
        s = ''
        codes.add(r'YYY')
\end{verbatim}





Який рядок треба записати на місце \kw{XXX} ?\par
\vspace{5mm}


Який рядок треба записати на місце \kw{YYY} ?\par


\newpage

{\begin {center}\textbf{ 189}\end {center}}\par
Файл \code{'1.txt'} містить відомість з рядками, в яких через двокрапку записано поряд\-ковий номер (ціле), назву предмету (знаків пунктуації не містить), прізвище студента, ім'я студента, номер заліковки, оцінку в балах за 100 бальною шкалою.

Білі символи навколо роздільників не допускаються.

Приклад рядка файлу



\begin{verbatim}
13:algebra:surname:name:123456:77
\end{verbatim}


Нижче наведено код, який має залишити у файлі в кожному рядку тільки номер заліковки та предмет.

Рядок з прикладу має бути замінений на такий



\begin{verbatim}
123456:algebra
\end{verbatim}




Код:



\begin{verbatim}
regex = re.compile(r'XXX')
with open('1.txt', encoding='utf8') as f:
    lst = f.readlines()
for i in range(len(lst)):
    lst[i] = re.sub(regex, r'YYY', lst[i])
with open('1.txt', 'w', encoding='utf8') as f:
    f.writelines(lst)
\end{verbatim}









Який рядок треба записати на місце \kw{XXX} ?\par
\vspace{5mm}


Який рядок треба записати на місце \kw{YYY} ?\par


\newpage

{\begin {center}\textbf{ 190}\end {center}}\par
Файл \code{'1.txt'} містить відомість з рядками, в яких через кому записано поряд\-ковий номер (ціле), назву предмету (знаків пунктуації не містить), прізвище студента, ім'я студента, номер заліковки, оцінку в балах за 100 бальною шкалою.

Білі символи навколо роздільників не допускаються.

Приклад рядка файлу



\begin{verbatim}
13,algebra,surname,name,123456,77
\end{verbatim}


Нижче наведено код, який шукає усі назви предметів, що зустрічаються у файлі, та зберігає їх у змінній \id{codes}.



\begin{verbatim}
codes = set()
regex = re.compile(r'XXX')
with open('1.txt', encoding='utf8') as f:
    for s in f:
        res = re.fullmatch(regex, s)
        if not res: continue
        s = ''
        codes.add(r'YYY')
\end{verbatim}





Який рядок треба записати на місце \kw{XXX} ?\par
\vspace{5mm}


Який рядок треба записати на місце \kw{YYY} ?\par


\newpage

{\begin {center}\textbf{ 191}\end {center}}\par
Файл \code{'1.txt'} містить відомість з рядками, в яких через двокрапку записано поряд\-ковий номер (ціле), назву предмету (знаків пунктуації не містить), прізвище студента, ім'я студента, номер заліковки, оцінку в балах за 100 бальною шкалою.

Білі символи навколо роздільників не допускаються.

Приклад рядка файлу



\begin{verbatim}
13:algebra:surname:name:123456:77
\end{verbatim}


Нижче наведено код, який має в кожному рядку файлу переставити місцями назву предмету та номер заліковки.

Рядок з прикладу має бути замінений на такий



\begin{verbatim}
13:123456:surname:name:algebra:77
\end{verbatim}




Код:



\begin{verbatim}
regex = re.compile(r'XXX')
with open('1.txt', encoding='utf8') as f:
    lst = f.readlines()
for i in range(len(lst)):
    lst[i] = re.sub(regex, r'YYY', lst[i])
with open('1.txt', 'w', encoding='utf8') as f:
    f.writelines(lst)
\end{verbatim}







Який рядок треба записати на місце \kw{XXX} ?\par
\vspace{5mm}


Який рядок треба записати на місце \kw{YYY} ?\par


\newpage

{\begin {center}\textbf{ 192}\end {center}}\par
Рядки відомості через кому містять порядковий номер (ціле), назву предмету (зна\-ків пунктуації не містить), прізвище студента, ім'я студента, оцінку в балах за 100 бальною шкалою.

Білі символи навколо роздільників не допускаються.

Приклад такого рядка присвоєно змінній \kw{s}.



\begin{verbatim}
s = '13,algebra,surname,name,77'
\end{verbatim}


Нижче наведено код, за виконання якого значенням змінної \kw{out} має стати прізвище студента.

Для наведеного прикладу значенням змінної \kw{out} має стати рядок \code{surname}



\begin{verbatim}
res = re.fullmatch(r'XXX', s)
s = ''
out = r'YYY'
\end{verbatim}


Код має працювати не тільки для конкретного рядка з прикладу, але й для кожного рядка з відомості.




Який рядок треба записати на місце \kw{XXX} ?\par
\vspace{5mm}


Який рядок треба записати на місце \kw{YYY} ?\par


\newpage

{\begin {center}\textbf{ 193}\end {center}}\par
Файл \code{'1.txt'} містить відомість з рядками, в яких через кому записа\-но поряд\-ковий номер (ціле), назву предмету (знаків пунктуації не містить), прізвище студента, ім'я студента, номер заліковки, оцінку в балах за 100 бальною шкалою.

Білі символи навколо роздільників не допускаються.

Приклад рядка файлу



\begin{verbatim}
13,algebra,surname,name,123456,77
\end{verbatim}


Нижче наведено код, який має залишити у файлі в кожному рядку тільки номер заліковки, назву предмету та оцінку в балах за 100-бальною шкалою (у заданій послі\-дов\-ності).

Рядок з прикладу має бути замінений на такий



\begin{verbatim}
123456,algebra,77
\end{verbatim}




Код



\begin{verbatim}
regex = re.compile(r'XXX')
with open('1.txt', encoding='utf8') as f:
    lst = f.readlines()
for i in range(len(lst)):
    lst[i] = re.sub(regex, r'YYY', lst[i])
with open('1.txt', 'w', encoding='utf8') as f:
    f.writelines(lst)
\end{verbatim}







Який рядок треба записати на місце \kw{XXX} ?\par
\vspace{5mm}


Який рядок треба записати на місце \kw{YYY} ?\par


\newpage

{\begin {center}\textbf{ 194}\end {center}}\par
Файл \code{'1.txt'} містить журнал з рядками, в яких через кому записано поряд\-ко\-вий номер запису (ціле), ім'я користувача (складається з латиниці), час події (фор\-мат як у прикладі), тип події (ERROR, WARNING, INFO, STATUS), код події (ціле).

Білі символи навколо роздільників не допускаються.

Приклад рядка журналу



\begin{verbatim}
12,username,2022-05-05 13:26,ERROR,404
\end{verbatim}


Нижче наведено код, який шукає усі події, що відбулись\\ для користувача \code{anonymous}, та зберігає їх типи в змінній \kw{codes}.



\begin{verbatim}
codes = set()
regex = re.compile(r'XXX')
with open('1.txt', encoding='utf8') as f:
    for s in f:
        res = re.fullmatch(regex, s)
        if not res: continue
        s = ''
        codes.add(r'YYY')
\end{verbatim}





Який рядок треба записати на місце \kw{XXX} ?\par
\vspace{5mm}


Який рядок треба записати на місце \kw{YYY} ?\par


\newpage

{\begin {center}\textbf{ 195}\end {center}}\par
Файл \code{'1.txt'} містить відомість з рядками, в яких через кому записа\-но назву предмету (знаків пунктуації не містить), прізвище студента, ім'я студента, номер залі\-ков\-ки, оцінку в балах за 100 бальною шкалою.

Білі символи навколо роздільників не допускаються.

Приклад рядка файлу



\begin{verbatim}
algebra,surname,name,123456,77
\end{verbatim}


Нижче наведено код, який має замінити у файлі усі номери заліковок на рядок \code{***}.



\begin{verbatim}
regex = re.compile(r'XXX')
with open('1.txt', encoding='utf8') as f:
    lst = f.readlines()
for i in range(len(lst)):
    lst[i] = re.sub(regex, r'YYY', lst[i])
with open('1.txt', 'w', encoding='utf8') as f:
    f.writelines(lst)
\end{verbatim}







Який рядок треба записати на місце \kw{XXX} ?\par
\vspace{5mm}


Який рядок треба записати на місце \kw{YYY} ?\par


\newpage

{\begin {center}\textbf{ 196}\end {center}}\par
Рядки журналу через двокрапку містять тип події (ERROR, WARNING, INFO, STATUS), ім'я користувача (складається з латиниці), час події (фор\-мат як у прикладі), код події (ціле).

Білі символи навколо роздільників не допускаються.

Приклад такого рядка присвоєно змінній \kw{s}.



\begin{verbatim}
s = 'ERROR:username:2022-05-05 13-26:404'
\end{verbatim}


Нижче наведено код, за виконання якого значенням змінної \kw{out} має стати ім'я користувача.

Для наведеного прикладу значенням змінної \kw{out} має стати рядок \code{username}



\begin{verbatim}
res = re.fullmatch(r'XXX', s)
s = ''
out = r'YYY'
\end{verbatim}


Код має працювати не тільки для конкретного рядка з прикладу, але й для кожного рядка з журналу.




Який рядок треба записати на місце \kw{XXX} ?\par
\vspace{5mm}


Який рядок треба записати на місце \kw{YYY} ?\par


\newpage

{\begin {center}\textbf{ 197}\end {center}}\par
Файл \code{'1.txt'} містить журнал з рядками, в яких через кому записано поряд\-ко\-вий номер запису в журналі (ціле), тип події (ERROR, WARNING, INFO, STATUS), ім'я користувача (складається з латиниці), час події (фор\-мат як у прикладі), код події (ціле).

Білі символи навколо роздільників не допускаються.

Приклад рядка файлу



\begin{verbatim}
13,ERROR,username,2022-05-05 13:26,404
\end{verbatim}


Нижче наведено код, який шукає усі коди помилок (подія ERROR), що зустріча\-ють\-ся в журналі, та зберігає їх у змінній \kw{codes}.



\begin{verbatim}
codes = set()
regex = re.compile(r'XXX')
with open('1.txt', encoding='utf8') as f:
    for s in f:
        res = re.fullmatch(regex, s)
        if not res: continue
        s = ''
        codes.add(r'YYY')
\end{verbatim}





Який рядок треба записати на місце \kw{XXX} ?\par
\vspace{5mm}


Який рядок треба записати на місце \kw{YYY} ?\par


\newpage

{\begin {center}\textbf{ 198}\end {center}}\par
Рядки відомості через двокрапку містять порядковий номер (ціле), назву предмету (знаків пунктуації не містить), прізвище студента, ім'я студента, оцінку в балах за 100 бальною шкалою.

Білі символи навколо роздільників не допускаються.

Приклад такого рядка присвоєно змінній \kw{s}.



\begin{verbatim}
s = '13:algebra:surname:name:77'
\end{verbatim}


Нижче наведено код, за виконання якого значенням змінної \kw{out} має стати прізвище студента.

Для наведеного прикладу значенням змінної \kw{out} має стати рядок \code{surname}



\begin{verbatim}
res = re.fullmatch(r'XXX', s)
s = ''
out = r'YYY'
\end{verbatim}


Код має працювати не тільки для конкретного рядка з прикладу, але й для кожного рядка з відомості.




Який рядок треба записати на місце \kw{XXX} ?\par
\vspace{5mm}


Який рядок треба записати на місце \kw{YYY} ?\par


\newpage

{\begin {center}\textbf{ 199}\end {center}}\par
Файл \code{'1.txt'} містить відомість з рядками, в яких через кому записано поряд\-ковий номер (ціле), назву предмету (знаків пунктуації не містить), прізвище студента, ім'я студента, номер заліковки, оцінку в балах за 100 бальною шкалою.

Білі символи навколо роздільників не допускаються.

Приклад рядка файлу



\begin{verbatim}
13,algebra,surname,name,123456,77
\end{verbatim}


Нижче наведено код, який має в кожному рядку файлу переставити місцями назву предмету та номер заліковки.

Рядок з прикладу має бути замінений на такий



\begin{verbatim}
13,123456,surname,name,algebra,77
\end{verbatim}




Код



\begin{verbatim}
regex = re.compile(r'XXX')
with open('1.txt', encoding='utf8') as f:
    lst = f.readlines()
for i in range(len(lst)):
    lst[i] = re.sub(regex, r'YYY', lst[i])
with open('1.txt', 'w', encoding='utf8') as f:
    f.writelines(lst)
\end{verbatim}







Який рядок треба записати на місце \kw{XXX} ?\par
\vspace{5mm}


Який рядок треба записати на місце \kw{YYY} ?\par


\newpage

{\begin {center}\textbf{ 200}\end {center}}\par
Файл \code{'1.txt'} містить відомість з рядками, в яких через кому записано поряд\-ковий номер (ціле), назву предмету (знаків пунктуації не містить), прізвище студента, ім'я студента, номер заліковки, оцінку в балах за 100 бальною шкалою.

Білі символи навколо роздільників не допускаються.

Приклад рядка файлу



\begin{verbatim}
13,algebra,surname,name,123456,77
\end{verbatim}


Нижче наведено код, який має залишити у файлі в кожному рядку тільки номер заліковки та предмет.

Рядок з прикладу має бути замінений на такий



\begin{verbatim}
123456,algebra
\end{verbatim}




Код



\begin{verbatim}
regex = re.compile(r'XXX')
with open('1.txt', encoding='utf8') as f:
    lst = f.readlines()
for i in range(len(lst)):
    lst[i] = re.sub(regex, r'YYY', lst[i])
with open('1.txt', 'w', encoding='utf8') as f:
    f.writelines(lst)
\end{verbatim}









Який рядок треба записати на місце \kw{XXX} ?\par
\vspace{5mm}


Який рядок треба записати на місце \kw{YYY} ?\par


\newpage

{\begin {center}\textbf{ 201}\end {center}}\par
Файл \code{'1.txt'} містить журнал з рядками, в яких через двокрапку записано поряд\-ко\-вий номер запису в журналі (ціле), тип події (ERROR, WARNING, INFO, STATUS), ім'я користувача (складається з латиниці), час події (фор\-мат як у прикладі), опис події.

Білі символи навколо роздільників не допускаються.

Приклад рядка файлу



\begin{verbatim}
13:ERROR:username:2022-05-05 13-26:something happens
\end{verbatim}


Нижче наведено код, який має залишити у файлі в кожному рядку тільки тип події та її опис.

Рядок з прикладу має бути замінений на такий



\begin{verbatim}
ERROR:something happens
\end{verbatim}




Код:



\begin{verbatim}
regex = re.compile(r'XXX')
with open('1.txt', encoding='utf8') as f:
    lst = f.readlines()
for i in range(len(lst)):
    lst[i] = re.sub(regex, r'YYY', lst[i])
with open('1.txt', 'w', encoding='utf8') as f:
    f.writelines(lst)
\end{verbatim}









Який рядок треба записати на місце \kw{XXX} ?\par
\vspace{5mm}


Який рядок треба записати на місце \kw{YYY} ?\par


\newpage

{\begin {center}\textbf{ 202}\end {center}}\par
Рядки відомості через кому містять порядковий номер (ціле), пріз\-ви\-ще студен\-та, ім'я студента, номер заліковки, оцінку в балах за 100 бальною шкалою.

Білі символи навколо роздільників не допускаються.

Приклад такого рядка присвоєно змінній \kw{s}.



\begin{verbatim}
s = '13,surname,name,123456,77'
\end{verbatim}


Нижче наведено код, за виконання якого для рядка \kw{s} значенням змінної \kw{out} має стати номер заліковки.

Для наведеного прикладу значенням змінної \kw{out} має стати рядок \code{123456} .



\begin{verbatim}
res = re.fullmatch(r'XXX', s)
s = ''
out = r'YYY'
\end{verbatim}


Код має працювати не тільки для конкретного рядка з прикладу, але й для кожного рядка з відомості.




Який рядок треба записати на місце \kw{XXX} ?\par
\vspace{5mm}


Який рядок треба записати на місце \kw{YYY} ?\par


\newpage

{\begin {center}\textbf{ 203}\end {center}}\par
Рядки відомості через двокрапку містять порядковий номер (ціле), пріз\-ви\-ще студен\-та, ім'я студента, номер заліковки, оцінку в балах за 100 бальною шкалою.

Білі символи навколо роздільників не допускаються.

Приклад такого рядка присвоєно змінній \kw{s}.



\begin{verbatim}
s = '13:surname:name:123456:77'
\end{verbatim}


Нижче наведено код, за виконання якого для рядка \kw{s} значенням змінної \kw{out} має стати номер заліковки.

Для наведеного прикладу значенням змінної \kw{out} має стати рядок \code{123456} .



\begin{verbatim}
res = re.fullmatch(r'XXX', s)
s = ''
out = r'YYY'
\end{verbatim}


Код має працювати не тільки для конкретного рядка з прикладу, але й для кожного рядка з відомості.




Який рядок треба записати на місце \kw{XXX} ?\par
\vspace{5mm}


Який рядок треба записати на місце \kw{YYY} ?\par


\newpage

{\begin {center}\textbf{ 204}\end {center}}\par
Файл \code{'1.txt'} містить журнал з рядками, в яких через двокрапку записано поряд\-ко\-вий номер запису (ціле), ім'я користувача (складається з латиниці), час події (фор\-мат як у прикладі), тип події (ERROR, WARNING, INFO, STATUS), код події (ціле).

Білі символи навколо роздільників не допускаються.

Приклад рядка журналу



\begin{verbatim}
12:username:2022-05-05 13-26:ERROR:404
\end{verbatim}


Нижче наведено код, який шукає усі попередження (подія WARNING), що зустрі\-ча\-ються в журналі і відбулися під час обідньої перерви деканату (з 13:00 по 13:59 включно), та зберігає їх коди в змінній \kw{codes}.



\begin{verbatim}
codes = set()
regex = re.compile(r'XXX')
with open('1.txt', encoding='utf8') as f:
    for s in f:
        res = re.fullmatch(regex, s)
        if not res: continue
        s = ''
        codes.add(r'YYY')
\end{verbatim}







Який рядок треба записати на місце \kw{XXX} ?\par
\vspace{5mm}


Який рядок треба записати на місце \kw{YYY} ?\par


\newpage

{\begin {center}\textbf{ 205}\end {center}}\par
Файл \code{'1.txt'} містить відомість з рядками, в яких через двокрапку записано поряд\-ковий номер (ціле), назву предмету (знаків пунктуації не містить), прізвище студента, ім'я студента, номер заліковки, оцінку в балах за 100 бальною шкалою.

Білі символи навколо роздільників не допускаються.

Приклад рядка файлу



\begin{verbatim}
13:algebra:surname:name:123456:77
\end{verbatim}


Нижче наведено код, який має залишити у файлі в кожному рядку тільки номер заліковки та предмет.

Рядок з прикладу має бути замінений на такий



\begin{verbatim}
123456:algebra
\end{verbatim}




Код:



\begin{verbatim}
regex = re.compile(r'XXX')
with open('1.txt', encoding='utf8') as f:
    lst = f.readlines()
for i in range(len(lst)):
    lst[i] = re.sub(regex, r'YYY', lst[i])
with open('1.txt', 'w', encoding='utf8') as f:
    f.writelines(lst)
\end{verbatim}









Який рядок треба записати на місце \kw{XXX} ?\par
\vspace{5mm}


Який рядок треба записати на місце \kw{YYY} ?\par


\newpage

{\begin {center}\textbf{ 206}\end {center}}\par
Файл \code{'1.txt'} містить журнал з рядками, в яких через кому записано поряд\-ко\-вий номер запису (ціле), ім'я користувача (складається з латиниці), час події (фор\-мат як у прикладі), тип події (ERROR, WARNING, INFO, STATUS), код події (ціле).

Білі символи навколо роздільників не допускаються.

Приклад рядка журналу



\begin{verbatim}
12,username,2022-05-05 13:26,ERROR,404
\end{verbatim}


Нижче наведено код, який шукає усі попередження (подія WARNING), що зустрі\-ча\-ються в журналі і відбулися під час обідньої перерви деканату (з 13,00 по 13,59 включно), та зберігає їх коди в змінній \kw{codes}.



\begin{verbatim}
codes = set()
regex = re.compile(r'XXX')
with open('1.txt', encoding='utf8') as f:
    for s in f:
        res = re.fullmatch(regex, s)
        if not res: continue
        s = ''
        codes.add(r'YYY')
\end{verbatim}







Який рядок треба записати на місце \kw{XXX} ?\par
\vspace{5mm}


Який рядок треба записати на місце \kw{YYY} ?\par


\newpage

{\begin {center}\textbf{ 207}\end {center}}\par
Файл \code{'1.txt'} містить відомість з рядками, в яких через кому записано поряд\-ковий номер (ціле), назву предмету (знаків пунктуації не містить), прізвище студента, ім'я студента, номер заліковки, оцінку в балах за 100 бальною шкалою.

Білі символи навколо роздільників не допускаються.

Приклад рядка файлу



\begin{verbatim}
13,algebra,surname,name,123456,77
\end{verbatim}


Нижче наведено код, який має в кожному рядку файлу переставити місцями назву предмету та номер заліковки.

Рядок з прикладу має бути замінений на такий



\begin{verbatim}
13,123456,surname,name,algebra,77
\end{verbatim}




Код



\begin{verbatim}
regex = re.compile(r'XXX')
with open('1.txt', encoding='utf8') as f:
    lst = f.readlines()
for i in range(len(lst)):
    lst[i] = re.sub(regex, r'YYY', lst[i])
with open('1.txt', 'w', encoding='utf8') as f:
    f.writelines(lst)
\end{verbatim}







Який рядок треба записати на місце \kw{XXX} ?\par
\vspace{5mm}


Який рядок треба записати на місце \kw{YYY} ?\par


\newpage

{\begin {center}\textbf{ 208}\end {center}}\par
Файл \code{'1.txt'} містить відомість з рядками, в яких через двокрапку записа\-но назву предмету (знаків пунктуації не містить), прізвище студента, ім'я студента, номер залі\-ков\-ки, оцінку в балах за 100 бальною шкалою.

Білі символи навколо роздільників не допускаються.

Приклад рядка файлу



\begin{verbatim}
algebra:surname:name:123456:77
\end{verbatim}


Нижче наведено код, який має замінити у файлі усі номери заліковок на рядок \code{***}.



\begin{verbatim}
regex = re.compile(r'XXX')
with open('1.txt', encoding='utf8') as f:
    lst = f.readlines()
for i in range(len(lst)):
    lst[i] = re.sub(regex, r'YYY', lst[i])
with open('1.txt', 'w', encoding='utf8') as f:
    f.writelines(lst)
\end{verbatim}







Який рядок треба записати на місце \kw{XXX} ?\par
\vspace{5mm}


Який рядок треба записати на місце \kw{YYY} ?\par


\newpage

{\begin {center}\textbf{ 209}\end {center}}\par
Файл \code{'1.txt'} містить журнал з рядками, в яких через кому записано поряд\-ко\-вий номер запису в журналі (ціле), тип події (ERROR, WARNING, INFO, STATUS), ім'я користувача (складається з латиниці), час події (фор\-мат як у прикладі), код події (ціле).

Білі символи навколо роздільників не допускаються.

Приклад рядка файлу



\begin{verbatim}
13,ERROR,username,2022-05-05 13:26,404
\end{verbatim}


Нижче наведено код, який шукає усі коди помилок (подія ERROR), що зустріча\-ють\-ся в журналі, та зберігає їх у змінній \kw{codes}.



\begin{verbatim}
codes = set()
regex = re.compile(r'XXX')
with open('1.txt', encoding='utf8') as f:
    for s in f:
        res = re.fullmatch(regex, s)
        if not res: continue
        s = ''
        codes.add(r'YYY')
\end{verbatim}





Який рядок треба записати на місце \kw{XXX} ?\par
\vspace{5mm}


Який рядок треба записати на місце \kw{YYY} ?\par


\newpage

{\begin {center}\textbf{ 210}\end {center}}\par
Рядки відомості через кому містять порядковий номер (ціле), назву предмету (зна\-ків пунктуації не містить), прізвище студента, ім'я студента, оцінку в балах за 100 бальною шкалою.

Білі символи навколо роздільників не допускаються.

Приклад такого рядка присвоєно змінній \kw{s}.



\begin{verbatim}
s = '13,algebra,surname,name,77'
\end{verbatim}


Нижче наведено код, за виконання якого значенням змінної \kw{out} має стати прізвище студента.

Для наведеного прикладу значенням змінної \kw{out} має стати рядок \code{surname}



\begin{verbatim}
res = re.fullmatch(r'XXX', s)
s = ''
out = r'YYY'
\end{verbatim}


Код має працювати не тільки для конкретного рядка з прикладу, але й для кожного рядка з відомості.




Який рядок треба записати на місце \kw{XXX} ?\par
\vspace{5mm}


Який рядок треба записати на місце \kw{YYY} ?\par


\newpage

{\begin {center}\textbf{ 211}\end {center}}\par
Файл \code{'1.txt'} містить відомість з рядками, в яких через кому записано поряд\-ковий номер (ціле), назву предмету (знаків пунктуації не містить), прізвище студента, ім'я студента, номер заліковки, оцінку в балах за 100 бальною шкалою.

Білі символи навколо роздільників не допускаються.

Приклад рядка файлу



\begin{verbatim}
13,algebra,surname,name,123456,77
\end{verbatim}


Нижче наведено код, який шукає усі назви предметів, що зустрічаються у файлі, та зберігає їх у змінній \id{codes}.



\begin{verbatim}
codes = set()
regex = re.compile(r'XXX')
with open('1.txt', encoding='utf8') as f:
    for s in f:
        res = re.fullmatch(regex, s)
        if not res: continue
        s = ''
        codes.add(r'YYY')
\end{verbatim}





Який рядок треба записати на місце \kw{XXX} ?\par
\vspace{5mm}


Який рядок треба записати на місце \kw{YYY} ?\par


\newpage

{\begin {center}\textbf{ 212}\end {center}}\par
Файл \code{'1.txt'} містить журнал з рядками, в яких через двокрапку записано поряд\-ко\-вий номер запису в журналі (ціле), тип події (ERROR, WARNING, INFO, STATUS), ім'я користувача (складається з латиниці), час події (фор\-мат як у прикладі), код події (ціле).

Білі символи навколо роздільників не допускаються.

Приклад рядка файлу



\begin{verbatim}
13:ERROR:username:2022-05-05 13-26:404
\end{verbatim}


Нижче наведено код, який має замінити у файлі імена користувачів на рядок \code{unknown}.



\begin{verbatim}
regex = re.compile(r'XXX')
with open('1.txt', encoding='utf8') as f:
    lst = f.readlines()
for i in range(len(lst)):
    lst[i] = re.sub(regex, r'YYY', lst[i])
with open('1.txt', 'w', encoding='utf8') as f:
    f.writelines(lst)
\end{verbatim}







Який рядок треба записати на місце \kw{XXX} ?\par
\vspace{5mm}


Який рядок треба записати на місце \kw{YYY} ?\par


\newpage

{\begin {center}\textbf{ 213}\end {center}}\par
Файл \code{'1.txt'} містить журнал з рядками, в яких через двокрапку записано поряд\-ко\-вий номер запису в журналі (ціле), тип події (ERROR, WARNING, INFO, STATUS), ім'я користувача (складається з латиниці), час події (фор\-мат як у прикладі), опис події.

Білі символи навколо роздільників не допускаються.

Приклад рядка файлу



\begin{verbatim}
13:ERROR:username:2022-05-05 13-26:something happens
\end{verbatim}


Нижче наведено код, який має залишити у файлі в кожному рядку тільки тип події та її опис.

Рядок з прикладу має бути замінений на такий



\begin{verbatim}
ERROR:something happens
\end{verbatim}




Код:



\begin{verbatim}
regex = re.compile(r'XXX')
with open('1.txt', encoding='utf8') as f:
    lst = f.readlines()
for i in range(len(lst)):
    lst[i] = re.sub(regex, r'YYY', lst[i])
with open('1.txt', 'w', encoding='utf8') as f:
    f.writelines(lst)
\end{verbatim}









Який рядок треба записати на місце \kw{XXX} ?\par
\vspace{5mm}


Який рядок треба записати на місце \kw{YYY} ?\par


\newpage

{\begin {center}\textbf{ 214}\end {center}}\par
Файл \code{'1.txt'} містить журнал з рядками, в яких через двокрапку записано поряд\-ко\-вий номер запису (ціле), ім'я користувача (складається з латиниці), час події (фор\-мат як у прикладі), тип події (ERROR, WARNING, INFO, STATUS), код події (ціле).

Білі символи навколо роздільників не допускаються.

Приклад рядка журналу



\begin{verbatim}
12:username:2022-05-05 13-26:ERROR:404
\end{verbatim}


Нижче наведено код, який шукає усі попередження (подія WARNING), що зустрі\-ча\-ються в журналі і відбулися під час обідньої перерви деканату (з 13:00 по 13:59 включно), та зберігає їх коди в змінній \kw{codes}.



\begin{verbatim}
codes = set()
regex = re.compile(r'XXX')
with open('1.txt', encoding='utf8') as f:
    for s in f:
        res = re.fullmatch(regex, s)
        if not res: continue
        s = ''
        codes.add(r'YYY')
\end{verbatim}







Який рядок треба записати на місце \kw{XXX} ?\par
\vspace{5mm}


Який рядок треба записати на місце \kw{YYY} ?\par


\newpage

{\begin {center}\textbf{ 215}\end {center}}\par
Рядки журналу через кому містять тип події (ERROR, WARNING, INFO, STATUS), ім'я користувача (складається з латиниці), час події (фор\-мат як у прикладі), код події (ціле).

Білі символи навколо роздільників не допускаються.

Приклад такого рядка присвоєно змінній \kw{s}.



\begin{verbatim}
s = 'ERROR,username,2022-05-05 13:26,404'
\end{verbatim}


Нижче наведено код, за виконання якого значенням змінної \kw{out} має стати ім'я користувача.

Для наведеного прикладу значенням змінної \kw{out} має стати рядок \code{username}



\begin{verbatim}
res = re.fullmatch(r'XXX', s)
s = ''
out = r'YYY'
\end{verbatim}


Код має працювати не тільки для конкретного рядка з прикладу, але й для кожного рядка з журналу.




Який рядок треба записати на місце \kw{XXX} ?\par
\vspace{5mm}


Який рядок треба записати на місце \kw{YYY} ?\par


\newpage

{\begin {center}\textbf{ 216}\end {center}}\par
Файл \code{'1.txt'} містить відомість з рядками, в яких через двокрапку записа\-но поряд\-ковий номер (ціле), назву предмету (знаків пунктуації не містить), прізвище студента, ім'я студента, номер заліковки, оцінку в балах за 100 бальною шкалою.

Білі символи навколо роздільників не допускаються.

Приклад рядка файлу



\begin{verbatim}
13:algebra:surname:name:123456:77
\end{verbatim}


Нижче наведено код, який має залишити у файлі в кожному рядку тільки номер заліковки, назву предмету та оцінку в балах за 100-бальною шкалою (у заданій послі\-дов\-ності).

Рядок з прикладу має бути замінений на такий



\begin{verbatim}
123456:algebra:77
\end{verbatim}




Код:



\begin{verbatim}
regex = re.compile(r'XXX')
with open('1.txt', encoding='utf8') as f:
    lst = f.readlines()
for i in range(len(lst)):
    lst[i] = re.sub(regex, r'YYY', lst[i])
with open('1.txt', 'w', encoding='utf8') as f:
    f.writelines(lst)
\end{verbatim}







Який рядок треба записати на місце \kw{XXX} ?\par
\vspace{5mm}


Який рядок треба записати на місце \kw{YYY} ?\par


\newpage

{\begin {center}\textbf{ 217}\end {center}}\par
Рядки журналу через кому містять ім'я користувача (складається з латиниці), час події (фор\-мат як у прикладі), тип події (ERROR, WARNING, INFO, STATUS), код події (ціле).

Білі символи навколо роздільників не допускаються.

Приклад такого рядка присвоєно змінній \kw{s}.



\begin{verbatim}
s = 'username,2022-05-05 13:26,ERROR,404'
\end{verbatim}


Нижче наведено код, за виконання якого значенням змінної \kw{out} має стати код події.

Для наведеного прикладу значенням змінної \kw{out} має стати рядок \code{404}



\begin{verbatim}
res = re.fullmatch(r'XXX', s)
s = ''
out = r'YYY'
\end{verbatim}


Код має працювати не тільки для конкретного рядка з прикладу, але й для кожного рядка з журналу.




Який рядок треба записати на місце \kw{XXX} ?\par
\vspace{5mm}


Який рядок треба записати на місце \kw{YYY} ?\par


\newpage

{\begin {center}\textbf{ 218}\end {center}}\par
Файл \code{'1.txt'} містить відомість з рядками, в яких через кому записано поряд\-ковий номер (ціле), назву предмету (знаків пунктуації не містить), прізвище студента, ім'я студента, номер заліковки, оцінку в балах за 100 бальною шкалою.

Білі символи навколо роздільників не допускаються.

Приклад рядка файлу



\begin{verbatim}
13,algebra,surname,name,123456,77
\end{verbatim}


Нижче наведено код, який шукає усі предмети, який складав студент із заліковкою \code{777}, та зберігає їх типи в змінній \kw{codes}.



\begin{verbatim}
codes = set()
regex = re.compile(r'XXX')
with open('1.txt', encoding='utf8') as f:
    for s in f:
        res = re.fullmatch(regex, s)
        if not res: continue
        s = ''
        codes.add(r'YYY')
\end{verbatim}





Який рядок треба записати на місце \kw{XXX} ?\par
\vspace{5mm}


Який рядок треба записати на місце \kw{YYY} ?\par


\newpage

{\begin {center}\textbf{ 219}\end {center}}\par
Файл \code{'1.txt'} містить журнал з рядками, в яких через двокрапку записано поряд\-ко\-вий номер запису в журналі (ціле), код події (ціле), тип події (ERROR,\\ WARNING, INFO, STATUS), ім'я користувача (складається з латиниці), час події (фор\-мат як у прикладі), опис події.

Білі символи навколо роздільників не допускаються.

Приклад рядка файлу



\begin{verbatim}
13:404:ERROR:username:2022-05-05 13:26:something happens
\end{verbatim}


Нижче наведено код, який має в кожному рядку файли переставити місцями код та тип події.

Рядок з прикладу має бути замінений на такий



\begin{verbatim}
13:ERROR:404:username:2022-05-05 13:26:something happens
\end{verbatim}




Код:



\begin{verbatim}
regex = re.compile(r'XXX')
with open('1.txt', encoding='utf8') as f:
    lst = f.readlines()
for i in range(len(lst)):
    lst[i] = re.sub(regex, r'YYY', lst[i])
with open('1.txt', 'w', encoding='utf8') as f:
    f.writelines(lst)
\end{verbatim}







Який рядок треба записати на місце \kw{XXX} ?\par
\vspace{5mm}


Який рядок треба записати на місце \kw{YYY} ?\par


\newpage

{\begin {center}\textbf{ 220}\end {center}}\par
Файл \code{'1.txt'} містить журнал з рядками, в яких через кому записано поряд\-ко\-вий номер запису в журналі (ціле), код події (ціле), тип події (ERROR, \\WARNING, INFO, STATUS), ім'я користувача (складається з латиниці), час події (фор\-мат як у прикладі), опис події.

Білі символи навколо роздільників не допускаються.

Приклад рядка файлу



\begin{verbatim}
13,404,ERROR,username,2022-05-05 13:26,something happens
\end{verbatim}


Нижче наведено код, який має залишити у файлі в кожному рядку тільки код події та ім'я користувача. Рядок з прикладу має бути замінений на такий



\begin{verbatim}
404,username
\end{verbatim}




Код:



\begin{verbatim}
regex = re.compile(r'XXX')
with open('1.txt', encoding='utf8') as f:
    lst = f.readlines()
for i in range(len(lst)):
    lst[i] = re.sub(regex, r'YYY', lst[i])
with open('1.txt', 'w', encoding='utf8') as f:
    f.writelines(lst)
\end{verbatim}







Який рядок треба записати на місце \kw{XXX} ?\par
\vspace{5mm}


Який рядок треба записати на місце \kw{YYY} ?\par


\newpage

{\begin {center}\textbf{ 221}\end {center}}\par
Рядки журналу через двокрапку містять тип події (ERROR, WARNING, INFO, STATUS), ім'я користувача (складається з латиниці), час події (фор\-мат як у прикладі), код події (ціле).

Білі символи навколо роздільників не допускаються.

Приклад такого рядка присвоєно змінній \kw{s}.



\begin{verbatim}
s = 'ERROR:username:2022-05-05 13-26:404'
\end{verbatim}


Нижче наведено код, за виконання якого значенням змінної \kw{out} має стати ім'я користувача.

Для наведеного прикладу значенням змінної \kw{out} має стати рядок \code{username}



\begin{verbatim}
res = re.fullmatch(r'XXX', s)
s = ''
out = r'YYY'
\end{verbatim}


Код має працювати не тільки для конкретного рядка з прикладу, але й для кожного рядка з журналу.




Який рядок треба записати на місце \kw{XXX} ?\par
\vspace{5mm}


Який рядок треба записати на місце \kw{YYY} ?\par


\newpage

{\begin {center}\textbf{ 222}\end {center}}\par
Файл \code{'1.txt'} містить журнал з рядками, в яких через двокрапку записано поряд\-ко\-вий номер запису в журналі (ціле), тип події (ERROR, WARNING, INFO, STATUS), ім'я користувача (складається з латиниці), час події (фор\-мат як у прикладі), код події (ціле).

Білі символи навколо роздільників не допускаються.

Приклад рядка файлу



\begin{verbatim}
13:ERROR:username:2022-05-05 13-26:404
\end{verbatim}


Нижче наведено код, який шукає усі коди помилок (подія ERROR), що зустріча\-ють\-ся в журналі, та зберігає їх у змінній \kw{codes}.



\begin{verbatim}
codes = set()
regex = re.compile(r'XXX')
with open('1.txt', encoding='utf8') as f:
    for s in f:
        res = re.fullmatch(regex, s)
        if not res: continue
        s = ''
        codes.add(r'YYY')
\end{verbatim}





Який рядок треба записати на місце \kw{XXX} ?\par
\vspace{5mm}


Який рядок треба записати на місце \kw{YYY} ?\par


\newpage

{\begin {center}\textbf{ 223}\end {center}}\par
Файл \code{'1.txt'} містить журнал з рядками, в яких через кому записано поряд\-ко\-вий номер запису в журналі (ціле), тип події (ERROR, WARNING, INFO, STATUS), ім'я користувача (складається з латиниці), час події (фор\-мат як у прикладі), код події (ціле).

Білі символи навколо роздільників не допускаються.

Приклад рядка файлу



\begin{verbatim}
13,ERROR,username,2022-05-05 13:26,404
\end{verbatim}


Нижче наведено код, який має замінити у файлі імена користувачів на рядок \code{unknown}.



\begin{verbatim}
regex = re.compile(r'XXX')
with open('1.txt', encoding='utf8') as f:
    lst = f.readlines()
for i in range(len(lst)):
    lst[i] = re.sub(regex, r'YYY', lst[i])
with open('1.txt', 'w', encoding='utf8') as f:
    f.writelines(lst)
\end{verbatim}







Який рядок треба записати на місце \kw{XXX} ?\par
\vspace{5mm}


Який рядок треба записати на місце \kw{YYY} ?\par


\newpage

{\begin {center}\textbf{ 224}\end {center}}\par
Файл \code{'1.txt'} містить відомість з рядками, в яких через двокрапку записано поряд\-ковий номер (ціле), назву предмету (знаків пунктуації не містить), прізвище студента, ім'я студента, номер заліковки, оцінку в балах за 100 бальною шкалою.

Білі символи навколо роздільників не допускаються.

Приклад рядка файлу



\begin{verbatim}
13:algebra:surname:name:123456:77
\end{verbatim}


Нижче наведено код, який шукає усі назви предметів, що зустрічаються у файлі, та зберігає їх у змінній \id{codes}.



\begin{verbatim}
codes = set()
regex = re.compile(r'XXX')
with open('1.txt', encoding='utf8') as f:
    for s in f:
        res = re.fullmatch(regex, s)
        if not res: continue
        s = ''
        codes.add(r'YYY')
\end{verbatim}





Який рядок треба записати на місце \kw{XXX} ?\par
\vspace{5mm}


Який рядок треба записати на місце \kw{YYY} ?\par


\newpage

{\begin {center}\textbf{ 225}\end {center}}\par
Рядки відомості через двокрапку містять порядковий номер (ціле), пріз\-ви\-ще студен\-та, ім'я студента, номер заліковки, оцінку в балах за 100 бальною шкалою.

Білі символи навколо роздільників не допускаються.

Приклад такого рядка присвоєно змінній \kw{s}.



\begin{verbatim}
s = '13:surname:name:123456:77'
\end{verbatim}


Нижче наведено код, за виконання якого для рядка \kw{s} значенням змінної \kw{out} має стати номер заліковки.

Для наведеного прикладу значенням змінної \kw{out} має стати рядок \code{123456} .



\begin{verbatim}
res = re.fullmatch(r'XXX', s)
s = ''
out = r'YYY'
\end{verbatim}


Код має працювати не тільки для конкретного рядка з прикладу, але й для кожного рядка з відомості.




Який рядок треба записати на місце \kw{XXX} ?\par
\vspace{5mm}


Який рядок треба записати на місце \kw{YYY} ?\par


\newpage

{\begin {center}\textbf{ 226}\end {center}}\par
Файл \code{'1.txt'} містить журнал з рядками, в яких через кому записано поряд\-ко\-вий номер запису в журналі (ціле), код події (ціле), тип події (ERROR,\\ WARNING, INFO, STATUS), ім'я користувача (складається з латиниці), час події (фор\-мат як у прикладі), опис події.

Білі символи навколо роздільників не допускаються.

Приклад рядка файлу



\begin{verbatim}
13,404,ERROR,username,2022-05-05 13:26,something happens
\end{verbatim}


Нижче наведено код, який має в кожному рядку файли переставити місцями код та тип події.

Рядок з прикладу має бути замінений на такий



\begin{verbatim}
13,ERROR,404,username,2022-05-05 13:26,something happens
\end{verbatim}




Код:



\begin{verbatim}
regex = re.compile(r'XXX')
with open('1.txt', encoding='utf8') as f:
    lst = f.readlines()
for i in range(len(lst)):
    lst[i] = re.sub(regex, r'YYY', lst[i])
with open('1.txt', 'w', encoding='utf8') as f:
    f.writelines(lst)
\end{verbatim}







Який рядок треба записати на місце \kw{XXX} ?\par
\vspace{5mm}


Який рядок треба записати на місце \kw{YYY} ?\par


\newpage

{\begin {center}\textbf{ 227}\end {center}}\par
Файл \code{'1.txt'} містить відомість з рядками, в яких через кому записа\-но поряд\-ковий номер (ціле), назву предмету (знаків пунктуації не містить), прізвище студента, ім'я студента, номер заліковки, оцінку в балах за 100 бальною шкалою.

Білі символи навколо роздільників не допускаються.

Приклад рядка файлу



\begin{verbatim}
13,algebra,surname,name,123456,77
\end{verbatim}


Нижче наведено код, який має залишити у файлі в кожному рядку тільки номер заліковки, назву предмету та оцінку в балах за 100-бальною шкалою (у заданій послі\-дов\-ності).

Рядок з прикладу має бути замінений на такий



\begin{verbatim}
123456,algebra,77
\end{verbatim}




Код



\begin{verbatim}
regex = re.compile(r'XXX')
with open('1.txt', encoding='utf8') as f:
    lst = f.readlines()
for i in range(len(lst)):
    lst[i] = re.sub(regex, r'YYY', lst[i])
with open('1.txt', 'w', encoding='utf8') as f:
    f.writelines(lst)
\end{verbatim}







Який рядок треба записати на місце \kw{XXX} ?\par
\vspace{5mm}


Який рядок треба записати на місце \kw{YYY} ?\par


\newpage

{\begin {center}\textbf{ 228}\end {center}}\par
Файл \code{'1.txt'} містить журнал з рядками, в яких через кому записано поряд\-ко\-вий номер запису в журналі (ціле), тип події (ERROR, WARNING, INFO, STATUS), ім'я користувача (складається з латиниці), час події (фор\-мат як у прикладі), опис події.

Білі символи навколо роздільників не допускаються.

Приклад рядка файлу



\begin{verbatim}
13,ERROR,username,2022-05-05 13:26,something happens
\end{verbatim}


Нижче наведено код, який має залишити у файлі в кожному рядку тільки тип події та її опис.

Рядок з прикладу має бути замінений на такий



\begin{verbatim}
ERROR,something happens
\end{verbatim}




Код:



\begin{verbatim}
regex = re.compile(r'XXX')
with open('1.txt', encoding='utf8') as f:
    lst = f.readlines()
for i in range(len(lst)):
    lst[i] = re.sub(regex, r'YYY', lst[i])
with open('1.txt', 'w', encoding='utf8') as f:
    f.writelines(lst)
\end{verbatim}









Який рядок треба записати на місце \kw{XXX} ?\par
\vspace{5mm}


Який рядок треба записати на місце \kw{YYY} ?\par


\newpage

{\begin {center}\textbf{ 229}\end {center}}\par
Файл \code{'1.txt'} містить відомість з рядками, в яких через кому записано поряд\-ковий номер (ціле), назву предмету (знаків пунктуації не містить), прізвище студента, ім'я студента, номер заліковки, оцінку в балах за 100 бальною шкалою.

Білі символи навколо роздільників не допускаються.

Приклад рядка файлу



\begin{verbatim}
13,algebra,surname,name,123456,77
\end{verbatim}


Нижче наведено код, який має залишити у файлі в кожному рядку тільки номер заліковки та предмет.

Рядок з прикладу має бути замінений на такий



\begin{verbatim}
123456,algebra
\end{verbatim}




Код



\begin{verbatim}
regex = re.compile(r'XXX')
with open('1.txt', encoding='utf8') as f:
    lst = f.readlines()
for i in range(len(lst)):
    lst[i] = re.sub(regex, r'YYY', lst[i])
with open('1.txt', 'w', encoding='utf8') as f:
    f.writelines(lst)
\end{verbatim}









Який рядок треба записати на місце \kw{XXX} ?\par
\vspace{5mm}


Який рядок треба записати на місце \kw{YYY} ?\par


\newpage

{\begin {center}\textbf{ 230}\end {center}}\par
Файл \code{'1.txt'} містить відомість з рядками, в яких через двокрапку записано поряд\-ковий номер (ціле), назву предмету (знаків пунктуації не містить), прізвище студента, ім'я студента, номер заліковки, оцінку в балах за 100 бальною шкалою.

Білі символи навколо роздільників не допускаються.

Приклад рядка файлу



\begin{verbatim}
13:algebra:surname:name:123456:77
\end{verbatim}


Нижче наведено код, який має в кожному рядку файлу переставити місцями назву предмету та номер заліковки.

Рядок з прикладу має бути замінений на такий



\begin{verbatim}
13:123456:surname:name:algebra:77
\end{verbatim}




Код:



\begin{verbatim}
regex = re.compile(r'XXX')
with open('1.txt', encoding='utf8') as f:
    lst = f.readlines()
for i in range(len(lst)):
    lst[i] = re.sub(regex, r'YYY', lst[i])
with open('1.txt', 'w', encoding='utf8') as f:
    f.writelines(lst)
\end{verbatim}







Який рядок треба записати на місце \kw{XXX} ?\par
\vspace{5mm}


Який рядок треба записати на місце \kw{YYY} ?\par


\newpage

{\begin {center}\textbf{ 231}\end {center}}\par
Рядки журналу через двокрапку містять ім'я користувача (складається з латиниці), час події (фор\-мат як у прикладі), тип події (ERROR, WARNING, INFO, STATUS), код події (ціле).

Білі символи навколо роздільників не допускаються.

Приклад такого рядка присвоєно змінній \kw{s}.



\begin{verbatim}
s = 'username:2022-05-05 13-26:ERROR:404'
\end{verbatim}


Нижче наведено код, за виконання якого значенням змінної \kw{out} має стати код події.

Для наведеного прикладу значенням змінної \kw{out} має стати рядок \code{404}



\begin{verbatim}
res = re.fullmatch(r'XXX', s)
s = ''
out = r'YYY'
\end{verbatim}


Код має працювати не тільки для конкретного рядка з прикладу, але й для кожного рядка з журналу.




Який рядок треба записати на місце \kw{XXX} ?\par
\vspace{5mm}


Який рядок треба записати на місце \kw{YYY} ?\par


\newpage

{\begin {center}\textbf{ 232}\end {center}}\par
Файл \code{'1.txt'} містить журнал з рядками, в яких через двокрапку записано поряд\-ко\-вий номер запису (ціле), ім'я користувача (складається з латиниці), час події (фор\-мат як у прикладі), тип події (ERROR, WARNING, INFO, STATUS), код події (ціле).

Білі символи навколо роздільників не допускаються.

Приклад рядка журналу



\begin{verbatim}
12:username:2022-05-05 13-26:ERROR:404
\end{verbatim}


Нижче наведено код, який шукає усі події, що відбулись\\ для користувача \code{anonymous}, та зберігає їх типи в змінній \kw{codes}.



\begin{verbatim}
codes = set()
regex = re.compile(r'XXX')
with open('1.txt', encoding='utf8') as f:
    for s in f:
        res = re.fullmatch(regex, s)
        if not res: continue
        s = ''
        codes.add(r'YYY')
\end{verbatim}





Який рядок треба записати на місце \kw{XXX} ?\par
\vspace{5mm}


Який рядок треба записати на місце \kw{YYY} ?\par


\newpage

{\begin {center}\textbf{ 233}\end {center}}\par
Файл \code{'1.txt'} містить журнал з рядками, в яких через двокрапку записано поряд\-ко\-вий номер запису в журналі (ціле), код події (ціле), тип події (ERROR, \\WARNING, INFO, STATUS), ім'я користувача (складається з латиниці), час події (фор\-мат як у прикладі), опис події.

Білі символи навколо роздільників не допускаються.

Приклад рядка файлу



\begin{verbatim}
13:404:ERROR:username:2022-05-05 13-26:something happens
\end{verbatim}


Нижче наведено код, який має залишити у файлі в кожному рядку тільки код події та ім'я користувача. Рядок з прикладу має бути замінений на такий



\begin{verbatim}
404:username
\end{verbatim}




Код:



\begin{verbatim}
regex = re.compile(r'XXX')
with open('1.txt', encoding='utf8') as f:
    lst = f.readlines()
for i in range(len(lst)):
    lst[i] = re.sub(regex, r'YYY', lst[i])
with open('1.txt', 'w', encoding='utf8') as f:
    f.writelines(lst)
\end{verbatim}







Який рядок треба записати на місце \kw{XXX} ?\par
\vspace{5mm}


Який рядок треба записати на місце \kw{YYY} ?\par


\newpage

{\begin {center}\textbf{ 234}\end {center}}\par
Файл \code{'1.txt'} містить відомість з рядками, в яких через двокрапку записано поряд\-ковий номер (ціле), назву предмету (знаків пунктуації не містить), прізвище студента, ім'я студента, номер заліковки, оцінку в балах за 100 бальною шкалою.

Білі символи навколо роздільників не допускаються.

Приклад рядка файлу



\begin{verbatim}
13:algebra:surname:name:123456:77
\end{verbatim}


Нижче наведено код, який шукає усі предмети, який складав студент із заліковкою \code{777}, та зберігає їх типи в змінній \kw{codes}.



\begin{verbatim}
codes = set()
regex = re.compile(r'XXX')
with open('1.txt', encoding='utf8') as f:
    for s in f:
        res = re.fullmatch(regex, s)
        if not res: continue
        s = ''
        codes.add(r'YYY')
\end{verbatim}





Який рядок треба записати на місце \kw{XXX} ?\par
\vspace{5mm}


Який рядок треба записати на місце \kw{YYY} ?\par


\newpage

{\begin {center}\textbf{ 235}\end {center}}\par
Файл \code{'1.txt'} містить відомість з рядками, в яких через кому записа\-но назву предмету (знаків пунктуації не містить), прізвище студента, ім'я студента, номер залі\-ков\-ки, оцінку в балах за 100 бальною шкалою.

Білі символи навколо роздільників не допускаються.

Приклад рядка файлу



\begin{verbatim}
algebra,surname,name,123456,77
\end{verbatim}


Нижче наведено код, який має замінити у файлі усі номери заліковок на рядок \code{***}.



\begin{verbatim}
regex = re.compile(r'XXX')
with open('1.txt', encoding='utf8') as f:
    lst = f.readlines()
for i in range(len(lst)):
    lst[i] = re.sub(regex, r'YYY', lst[i])
with open('1.txt', 'w', encoding='utf8') as f:
    f.writelines(lst)
\end{verbatim}







Який рядок треба записати на місце \kw{XXX} ?\par
\vspace{5mm}


Який рядок треба записати на місце \kw{YYY} ?\par


\newpage

{\begin {center}\textbf{ 236}\end {center}}\par
Рядки відомості через кому містять порядковий номер (ціле), пріз\-ви\-ще студен\-та, ім'я студента, номер заліковки, оцінку в балах за 100 бальною шкалою.

Білі символи навколо роздільників не допускаються.

Приклад такого рядка присвоєно змінній \kw{s}.



\begin{verbatim}
s = '13,surname,name,123456,77'
\end{verbatim}


Нижче наведено код, за виконання якого для рядка \kw{s} значенням змінної \kw{out} має стати номер заліковки.

Для наведеного прикладу значенням змінної \kw{out} має стати рядок \code{123456} .



\begin{verbatim}
res = re.fullmatch(r'XXX', s)
s = ''
out = r'YYY'
\end{verbatim}


Код має працювати не тільки для конкретного рядка з прикладу, але й для кожного рядка з відомості.




Який рядок треба записати на місце \kw{XXX} ?\par
\vspace{5mm}


Який рядок треба записати на місце \kw{YYY} ?\par


\newpage

{\begin {center}\textbf{ 237}\end {center}}\par
Файл \code{'1.txt'} містить журнал з рядками, в яких через кому записано поряд\-ко\-вий номер запису (ціле), ім'я користувача (складається з латиниці), час події (фор\-мат як у прикладі), тип події (ERROR, WARNING, INFO, STATUS), код події (ціле).

Білі символи навколо роздільників не допускаються.

Приклад рядка журналу



\begin{verbatim}
12,username,2022-05-05 13:26,ERROR,404
\end{verbatim}


Нижче наведено код, який шукає усі події, що відбулись\\ для користувача \code{anonymous}, та зберігає їх типи в змінній \kw{codes}.



\begin{verbatim}
codes = set()
regex = re.compile(r'XXX')
with open('1.txt', encoding='utf8') as f:
    for s in f:
        res = re.fullmatch(regex, s)
        if not res: continue
        s = ''
        codes.add(r'YYY')
\end{verbatim}





Який рядок треба записати на місце \kw{XXX} ?\par
\vspace{5mm}


Який рядок треба записати на місце \kw{YYY} ?\par


\newpage

{\begin {center}\textbf{ 238}\end {center}}\par
Рядки відомості через двокрапку містять порядковий номер (ціле), назву предмету (знаків пунктуації не містить), прізвище студента, ім'я студента, оцінку в балах за 100 бальною шкалою.

Білі символи навколо роздільників не допускаються.

Приклад такого рядка присвоєно змінній \kw{s}.



\begin{verbatim}
s = '13:algebra:surname:name:77'
\end{verbatim}


Нижче наведено код, за виконання якого значенням змінної \kw{out} має стати прізвище студента.

Для наведеного прикладу значенням змінної \kw{out} має стати рядок \code{surname}



\begin{verbatim}
res = re.fullmatch(r'XXX', s)
s = ''
out = r'YYY'
\end{verbatim}


Код має працювати не тільки для конкретного рядка з прикладу, але й для кожного рядка з відомості.




Який рядок треба записати на місце \kw{XXX} ?\par
\vspace{5mm}


Який рядок треба записати на місце \kw{YYY} ?\par


\newpage

{\begin {center}\textbf{ 239}\end {center}}\par
Файл \code{'1.txt'} містить журнал з рядками, в яких через двокрапку записано поряд\-ко\-вий номер запису (ціле), ім'я користувача (складається з латиниці), час події (фор\-мат як у прикладі), тип події (ERROR, WARNING, INFO, STATUS), код події (ціле).

Білі символи навколо роздільників не допускаються.

Приклад рядка журналу



\begin{verbatim}
12:username:2022-05-05 13-26:ERROR:404
\end{verbatim}


Нижче наведено код, який шукає усі події, що відбулись\\ для користувача \code{anonymous}, та зберігає їх типи в змінній \kw{codes}.



\begin{verbatim}
codes = set()
regex = re.compile(r'XXX')
with open('1.txt', encoding='utf8') as f:
    for s in f:
        res = re.fullmatch(regex, s)
        if not res: continue
        s = ''
        codes.add(r'YYY')
\end{verbatim}





Який рядок треба записати на місце \kw{XXX} ?\par
\vspace{5mm}


Який рядок треба записати на місце \kw{YYY} ?\par


\newpage

{\begin {center}\textbf{ 240}\end {center}}\par
Файл \code{'1.txt'} містить відомість з рядками, в яких через кому записано поряд\-ковий номер (ціле), назву предмету (знаків пунктуації не містить), прізвище студента, ім'я студента, номер заліковки, оцінку в балах за 100 бальною шкалою.

Білі символи навколо роздільників не допускаються.

Приклад рядка файлу



\begin{verbatim}
13,algebra,surname,name,123456,77
\end{verbatim}


Нижче наведено код, який шукає усі предмети, який складав студент із заліковкою \code{777}, та зберігає їх типи в змінній \kw{codes}.



\begin{verbatim}
codes = set()
regex = re.compile(r'XXX')
with open('1.txt', encoding='utf8') as f:
    for s in f:
        res = re.fullmatch(regex, s)
        if not res: continue
        s = ''
        codes.add(r'YYY')
\end{verbatim}





Який рядок треба записати на місце \kw{XXX} ?\par
\vspace{5mm}


Який рядок треба записати на місце \kw{YYY} ?\par


\newpage

{\begin {center}\textbf{ 241}\end {center}}\par
Рядки журналу через двокрапку містять тип події (ERROR, WARNING, INFO, STATUS), ім'я користувача (складається з латиниці), час події (фор\-мат як у прикладі), код події (ціле).

Білі символи навколо роздільників не допускаються.

Приклад такого рядка присвоєно змінній \kw{s}.



\begin{verbatim}
s = 'ERROR:username:2022-05-05 13-26:404'
\end{verbatim}


Нижче наведено код, за виконання якого значенням змінної \kw{out} має стати ім'я користувача.

Для наведеного прикладу значенням змінної \kw{out} має стати рядок \code{username}



\begin{verbatim}
res = re.fullmatch(r'XXX', s)
s = ''
out = r'YYY'
\end{verbatim}


Код має працювати не тільки для конкретного рядка з прикладу, але й для кожного рядка з журналу.




Який рядок треба записати на місце \kw{XXX} ?\par
\vspace{5mm}


Який рядок треба записати на місце \kw{YYY} ?\par


\newpage

{\begin {center}\textbf{ 242}\end {center}}\par
Рядки журналу через кому містять ім'я користувача (складається з латиниці), час події (фор\-мат як у прикладі), тип події (ERROR, WARNING, INFO, STATUS), код події (ціле).

Білі символи навколо роздільників не допускаються.

Приклад такого рядка присвоєно змінній \kw{s}.



\begin{verbatim}
s = 'username,2022-05-05 13:26,ERROR,404'
\end{verbatim}


Нижче наведено код, за виконання якого значенням змінної \kw{out} має стати код події.

Для наведеного прикладу значенням змінної \kw{out} має стати рядок \code{404}



\begin{verbatim}
res = re.fullmatch(r'XXX', s)
s = ''
out = r'YYY'
\end{verbatim}


Код має працювати не тільки для конкретного рядка з прикладу, але й для кожного рядка з журналу.




Який рядок треба записати на місце \kw{XXX} ?\par
\vspace{5mm}


Який рядок треба записати на місце \kw{YYY} ?\par


\newpage

{\begin {center}\textbf{ 243}\end {center}}\par
Файл \code{'1.txt'} містить відомість з рядками, в яких через кому записано поряд\-ковий номер (ціле), назву предмету (знаків пунктуації не містить), прізвище студента, ім'я студента, номер заліковки, оцінку в балах за 100 бальною шкалою.

Білі символи навколо роздільників не допускаються.

Приклад рядка файлу



\begin{verbatim}
13,algebra,surname,name,123456,77
\end{verbatim}


Нижче наведено код, який шукає усі назви предметів, що зустрічаються у файлі, та зберігає їх у змінній \id{codes}.



\begin{verbatim}
codes = set()
regex = re.compile(r'XXX')
with open('1.txt', encoding='utf8') as f:
    for s in f:
        res = re.fullmatch(regex, s)
        if not res: continue
        s = ''
        codes.add(r'YYY')
\end{verbatim}





Який рядок треба записати на місце \kw{XXX} ?\par
\vspace{5mm}


Який рядок треба записати на місце \kw{YYY} ?\par


\newpage

{\begin {center}\textbf{ 244}\end {center}}\par
Рядки журналу через кому містять тип події (ERROR, WARNING, INFO, STATUS), ім'я користувача (складається з латиниці), час події (фор\-мат як у прикладі), код події (ціле).

Білі символи навколо роздільників не допускаються.

Приклад такого рядка присвоєно змінній \kw{s}.



\begin{verbatim}
s = 'ERROR,username,2022-05-05 13:26,404'
\end{verbatim}


Нижче наведено код, за виконання якого значенням змінної \kw{out} має стати ім'я користувача.

Для наведеного прикладу значенням змінної \kw{out} має стати рядок \code{username}



\begin{verbatim}
res = re.fullmatch(r'XXX', s)
s = ''
out = r'YYY'
\end{verbatim}


Код має працювати не тільки для конкретного рядка з прикладу, але й для кожного рядка з журналу.




Який рядок треба записати на місце \kw{XXX} ?\par
\vspace{5mm}


Який рядок треба записати на місце \kw{YYY} ?\par


\newpage

{\begin {center}\textbf{ 245}\end {center}}\par
Файл \code{'1.txt'} містить журнал з рядками, в яких через двокрапку записано поряд\-ко\-вий номер запису в журналі (ціле), тип події (ERROR, WARNING, INFO, STATUS), ім'я користувача (складається з латиниці), час події (фор\-мат як у прикладі), опис події.

Білі символи навколо роздільників не допускаються.

Приклад рядка файлу



\begin{verbatim}
13:ERROR:username:2022-05-05 13-26:something happens
\end{verbatim}


Нижче наведено код, який має залишити у файлі в кожному рядку тільки тип події та її опис.

Рядок з прикладу має бути замінений на такий



\begin{verbatim}
ERROR:something happens
\end{verbatim}




Код:



\begin{verbatim}
regex = re.compile(r'XXX')
with open('1.txt', encoding='utf8') as f:
    lst = f.readlines()
for i in range(len(lst)):
    lst[i] = re.sub(regex, r'YYY', lst[i])
with open('1.txt', 'w', encoding='utf8') as f:
    f.writelines(lst)
\end{verbatim}









Який рядок треба записати на місце \kw{XXX} ?\par
\vspace{5mm}


Який рядок треба записати на місце \kw{YYY} ?\par


\newpage

{\begin {center}\textbf{ 246}\end {center}}\par
Файл \code{'1.txt'} містить відомість з рядками, в яких через кому записано поряд\-ковий номер (ціле), назву предмету (знаків пунктуації не містить), прізвище студента, ім'я студента, номер заліковки, оцінку в балах за 100 бальною шкалою.

Білі символи навколо роздільників не допускаються.

Приклад рядка файлу



\begin{verbatim}
13,algebra,surname,name,123456,77
\end{verbatim}


Нижче наведено код, який має залишити у файлі в кожному рядку тільки номер заліковки та предмет.

Рядок з прикладу має бути замінений на такий



\begin{verbatim}
123456,algebra
\end{verbatim}




Код



\begin{verbatim}
regex = re.compile(r'XXX')
with open('1.txt', encoding='utf8') as f:
    lst = f.readlines()
for i in range(len(lst)):
    lst[i] = re.sub(regex, r'YYY', lst[i])
with open('1.txt', 'w', encoding='utf8') as f:
    f.writelines(lst)
\end{verbatim}









Який рядок треба записати на місце \kw{XXX} ?\par
\vspace{5mm}


Який рядок треба записати на місце \kw{YYY} ?\par


\newpage

{\begin {center}\textbf{ 247}\end {center}}\par
Файл \code{'1.txt'} містить журнал з рядками, в яких через двокрапку записано поряд\-ко\-вий номер запису в журналі (ціле), код події (ціле), тип події (ERROR, \\WARNING, INFO, STATUS), ім'я користувача (складається з латиниці), час події (фор\-мат як у прикладі), опис події.

Білі символи навколо роздільників не допускаються.

Приклад рядка файлу



\begin{verbatim}
13:404:ERROR:username:2022-05-05 13-26:something happens
\end{verbatim}


Нижче наведено код, який має залишити у файлі в кожному рядку тільки код події та ім'я користувача. Рядок з прикладу має бути замінений на такий



\begin{verbatim}
404:username
\end{verbatim}




Код:



\begin{verbatim}
regex = re.compile(r'XXX')
with open('1.txt', encoding='utf8') as f:
    lst = f.readlines()
for i in range(len(lst)):
    lst[i] = re.sub(regex, r'YYY', lst[i])
with open('1.txt', 'w', encoding='utf8') as f:
    f.writelines(lst)
\end{verbatim}







Який рядок треба записати на місце \kw{XXX} ?\par
\vspace{5mm}


Який рядок треба записати на місце \kw{YYY} ?\par


\newpage

{\begin {center}\textbf{ 248}\end {center}}\par
Файл \code{'1.txt'} містить відомість з рядками, в яких через двокрапку записа\-но назву предмету (знаків пунктуації не містить), прізвище студента, ім'я студента, номер залі\-ков\-ки, оцінку в балах за 100 бальною шкалою.

Білі символи навколо роздільників не допускаються.

Приклад рядка файлу



\begin{verbatim}
algebra:surname:name:123456:77
\end{verbatim}


Нижче наведено код, який має замінити у файлі усі номери заліковок на рядок \code{***}.



\begin{verbatim}
regex = re.compile(r'XXX')
with open('1.txt', encoding='utf8') as f:
    lst = f.readlines()
for i in range(len(lst)):
    lst[i] = re.sub(regex, r'YYY', lst[i])
with open('1.txt', 'w', encoding='utf8') as f:
    f.writelines(lst)
\end{verbatim}







Який рядок треба записати на місце \kw{XXX} ?\par
\vspace{5mm}


Який рядок треба записати на місце \kw{YYY} ?\par


\newpage

{\begin {center}\textbf{ 249}\end {center}}\par
Рядки відомості через кому містять порядковий номер (ціле), пріз\-ви\-ще студен\-та, ім'я студента, номер заліковки, оцінку в балах за 100 бальною шкалою.

Білі символи навколо роздільників не допускаються.

Приклад такого рядка присвоєно змінній \kw{s}.



\begin{verbatim}
s = '13,surname,name,123456,77'
\end{verbatim}


Нижче наведено код, за виконання якого для рядка \kw{s} значенням змінної \kw{out} має стати номер заліковки.

Для наведеного прикладу значенням змінної \kw{out} має стати рядок \code{123456} .



\begin{verbatim}
res = re.fullmatch(r'XXX', s)
s = ''
out = r'YYY'
\end{verbatim}


Код має працювати не тільки для конкретного рядка з прикладу, але й для кожного рядка з відомості.




Який рядок треба записати на місце \kw{XXX} ?\par
\vspace{5mm}


Який рядок треба записати на місце \kw{YYY} ?\par


\newpage

{\begin {center}\textbf{ 250}\end {center}}\par
Рядки відомості через двокрапку містять порядковий номер (ціле), назву предмету (знаків пунктуації не містить), прізвище студента, ім'я студента, оцінку в балах за 100 бальною шкалою.

Білі символи навколо роздільників не допускаються.

Приклад такого рядка присвоєно змінній \kw{s}.



\begin{verbatim}
s = '13:algebra:surname:name:77'
\end{verbatim}


Нижче наведено код, за виконання якого значенням змінної \kw{out} має стати прізвище студента.

Для наведеного прикладу значенням змінної \kw{out} має стати рядок \code{surname}



\begin{verbatim}
res = re.fullmatch(r'XXX', s)
s = ''
out = r'YYY'
\end{verbatim}


Код має працювати не тільки для конкретного рядка з прикладу, але й для кожного рядка з відомості.




Який рядок треба записати на місце \kw{XXX} ?\par
\vspace{5mm}


Який рядок треба записати на місце \kw{YYY} ?\par


\newpage

{\begin {center}\textbf{ 251}\end {center}}\par
Рядки відомості через кому містять порядковий номер (ціле), назву предмету (зна\-ків пунктуації не містить), прізвище студента, ім'я студента, оцінку в балах за 100 бальною шкалою.

Білі символи навколо роздільників не допускаються.

Приклад такого рядка присвоєно змінній \kw{s}.



\begin{verbatim}
s = '13,algebra,surname,name,77'
\end{verbatim}


Нижче наведено код, за виконання якого значенням змінної \kw{out} має стати прізвище студента.

Для наведеного прикладу значенням змінної \kw{out} має стати рядок \code{surname}



\begin{verbatim}
res = re.fullmatch(r'XXX', s)
s = ''
out = r'YYY'
\end{verbatim}


Код має працювати не тільки для конкретного рядка з прикладу, але й для кожного рядка з відомості.




Який рядок треба записати на місце \kw{XXX} ?\par
\vspace{5mm}


Який рядок треба записати на місце \kw{YYY} ?\par


\newpage

{\begin {center}\textbf{ 252}\end {center}}\par
Файл \code{'1.txt'} містить журнал з рядками, в яких через кому записано поряд\-ко\-вий номер запису (ціле), ім'я користувача (складається з латиниці), час події (фор\-мат як у прикладі), тип події (ERROR, WARNING, INFO, STATUS), код події (ціле).

Білі символи навколо роздільників не допускаються.

Приклад рядка журналу



\begin{verbatim}
12,username,2022-05-05 13:26,ERROR,404
\end{verbatim}


Нижче наведено код, який шукає усі події, що відбулись\\ для користувача \code{anonymous}, та зберігає їх типи в змінній \kw{codes}.



\begin{verbatim}
codes = set()
regex = re.compile(r'XXX')
with open('1.txt', encoding='utf8') as f:
    for s in f:
        res = re.fullmatch(regex, s)
        if not res: continue
        s = ''
        codes.add(r'YYY')
\end{verbatim}





Який рядок треба записати на місце \kw{XXX} ?\par
\vspace{5mm}


Який рядок треба записати на місце \kw{YYY} ?\par


\newpage

{\begin {center}\textbf{ 253}\end {center}}\par
Файл \code{'1.txt'} містить відомість з рядками, в яких через кому записа\-но поряд\-ковий номер (ціле), назву предмету (знаків пунктуації не містить), прізвище студента, ім'я студента, номер заліковки, оцінку в балах за 100 бальною шкалою.

Білі символи навколо роздільників не допускаються.

Приклад рядка файлу



\begin{verbatim}
13,algebra,surname,name,123456,77
\end{verbatim}


Нижче наведено код, який має залишити у файлі в кожному рядку тільки номер заліковки, назву предмету та оцінку в балах за 100-бальною шкалою (у заданій послі\-дов\-ності).

Рядок з прикладу має бути замінений на такий



\begin{verbatim}
123456,algebra,77
\end{verbatim}




Код



\begin{verbatim}
regex = re.compile(r'XXX')
with open('1.txt', encoding='utf8') as f:
    lst = f.readlines()
for i in range(len(lst)):
    lst[i] = re.sub(regex, r'YYY', lst[i])
with open('1.txt', 'w', encoding='utf8') as f:
    f.writelines(lst)
\end{verbatim}







Який рядок треба записати на місце \kw{XXX} ?\par
\vspace{5mm}


Який рядок треба записати на місце \kw{YYY} ?\par


\newpage

{\begin {center}\textbf{ 254}\end {center}}\par
Файл \code{'1.txt'} містить журнал з рядками, в яких через кому записано поряд\-ко\-вий номер запису в журналі (ціле), тип події (ERROR, WARNING, INFO, STATUS), ім'я користувача (складається з латиниці), час події (фор\-мат як у прикладі), опис події.

Білі символи навколо роздільників не допускаються.

Приклад рядка файлу



\begin{verbatim}
13,ERROR,username,2022-05-05 13:26,something happens
\end{verbatim}


Нижче наведено код, який має залишити у файлі в кожному рядку тільки тип події та її опис.

Рядок з прикладу має бути замінений на такий



\begin{verbatim}
ERROR,something happens
\end{verbatim}




Код:



\begin{verbatim}
regex = re.compile(r'XXX')
with open('1.txt', encoding='utf8') as f:
    lst = f.readlines()
for i in range(len(lst)):
    lst[i] = re.sub(regex, r'YYY', lst[i])
with open('1.txt', 'w', encoding='utf8') as f:
    f.writelines(lst)
\end{verbatim}









Який рядок треба записати на місце \kw{XXX} ?\par
\vspace{5mm}


Який рядок треба записати на місце \kw{YYY} ?\par


\newpage

{\begin {center}\textbf{ 255}\end {center}}\par
Файл \code{'1.txt'} містить журнал з рядками, в яких через двокрапку записано поряд\-ко\-вий номер запису в журналі (ціле), тип події (ERROR, WARNING, INFO, STATUS), ім'я користувача (складається з латиниці), час події (фор\-мат як у прикладі), код події (ціле).

Білі символи навколо роздільників не допускаються.

Приклад рядка файлу



\begin{verbatim}
13:ERROR:username:2022-05-05 13-26:404
\end{verbatim}


Нижче наведено код, який шукає усі коди помилок (подія ERROR), що зустріча\-ють\-ся в журналі, та зберігає їх у змінній \kw{codes}.



\begin{verbatim}
codes = set()
regex = re.compile(r'XXX')
with open('1.txt', encoding='utf8') as f:
    for s in f:
        res = re.fullmatch(regex, s)
        if not res: continue
        s = ''
        codes.add(r'YYY')
\end{verbatim}





Який рядок треба записати на місце \kw{XXX} ?\par
\vspace{5mm}


Який рядок треба записати на місце \kw{YYY} ?\par


\newpage

{\begin {center}\textbf{ 256}\end {center}}\par
Файл \code{'1.txt'} містить відомість з рядками, в яких через кому записано поряд\-ковий номер (ціле), назву предмету (знаків пунктуації не містить), прізвище студента, ім'я студента, номер заліковки, оцінку в балах за 100 бальною шкалою.

Білі символи навколо роздільників не допускаються.

Приклад рядка файлу



\begin{verbatim}
13,algebra,surname,name,123456,77
\end{verbatim}


Нижче наведено код, який має в кожному рядку файлу переставити місцями назву предмету та номер заліковки.

Рядок з прикладу має бути замінений на такий



\begin{verbatim}
13,123456,surname,name,algebra,77
\end{verbatim}




Код



\begin{verbatim}
regex = re.compile(r'XXX')
with open('1.txt', encoding='utf8') as f:
    lst = f.readlines()
for i in range(len(lst)):
    lst[i] = re.sub(regex, r'YYY', lst[i])
with open('1.txt', 'w', encoding='utf8') as f:
    f.writelines(lst)
\end{verbatim}







Який рядок треба записати на місце \kw{XXX} ?\par
\vspace{5mm}


Який рядок треба записати на місце \kw{YYY} ?\par


\newpage

{\begin {center}\textbf{ 257}\end {center}}\par
Файл \code{'1.txt'} містить журнал з рядками, в яких через двокрапку записано поряд\-ко\-вий номер запису в журналі (ціле), тип події (ERROR, WARNING, INFO, STATUS), ім'я користувача (складається з латиниці), час події (фор\-мат як у прикладі), код події (ціле).

Білі символи навколо роздільників не допускаються.

Приклад рядка файлу



\begin{verbatim}
13:ERROR:username:2022-05-05 13-26:404
\end{verbatim}


Нижче наведено код, який має замінити у файлі імена користувачів на рядок \code{unknown}.



\begin{verbatim}
regex = re.compile(r'XXX')
with open('1.txt', encoding='utf8') as f:
    lst = f.readlines()
for i in range(len(lst)):
    lst[i] = re.sub(regex, r'YYY', lst[i])
with open('1.txt', 'w', encoding='utf8') as f:
    f.writelines(lst)
\end{verbatim}







Який рядок треба записати на місце \kw{XXX} ?\par
\vspace{5mm}


Який рядок треба записати на місце \kw{YYY} ?\par


\newpage

{\begin {center}\textbf{ 258}\end {center}}\par
Файл \code{'1.txt'} містить відомість з рядками, в яких через двокрапку записано поряд\-ковий номер (ціле), назву предмету (знаків пунктуації не містить), прізвище студента, ім'я студента, номер заліковки, оцінку в балах за 100 бальною шкалою.

Білі символи навколо роздільників не допускаються.

Приклад рядка файлу



\begin{verbatim}
13:algebra:surname:name:123456:77
\end{verbatim}


Нижче наведено код, який шукає усі назви предметів, що зустрічаються у файлі, та зберігає їх у змінній \id{codes}.



\begin{verbatim}
codes = set()
regex = re.compile(r'XXX')
with open('1.txt', encoding='utf8') as f:
    for s in f:
        res = re.fullmatch(regex, s)
        if not res: continue
        s = ''
        codes.add(r'YYY')
\end{verbatim}





Який рядок треба записати на місце \kw{XXX} ?\par
\vspace{5mm}


Який рядок треба записати на місце \kw{YYY} ?\par


\newpage

{\begin {center}\textbf{ 259}\end {center}}\par
Файл \code{'1.txt'} містить відомість з рядками, в яких через кому записа\-но назву предмету (знаків пунктуації не містить), прізвище студента, ім'я студента, номер залі\-ков\-ки, оцінку в балах за 100 бальною шкалою.

Білі символи навколо роздільників не допускаються.

Приклад рядка файлу



\begin{verbatim}
algebra,surname,name,123456,77
\end{verbatim}


Нижче наведено код, який має замінити у файлі усі номери заліковок на рядок \code{***}.



\begin{verbatim}
regex = re.compile(r'XXX')
with open('1.txt', encoding='utf8') as f:
    lst = f.readlines()
for i in range(len(lst)):
    lst[i] = re.sub(regex, r'YYY', lst[i])
with open('1.txt', 'w', encoding='utf8') as f:
    f.writelines(lst)
\end{verbatim}







Який рядок треба записати на місце \kw{XXX} ?\par
\vspace{5mm}


Який рядок треба записати на місце \kw{YYY} ?\par


\newpage

{\begin {center}\textbf{ 260}\end {center}}\par
Файл \code{'1.txt'} містить відомість з рядками, в яких через двокрапку записано поряд\-ковий номер (ціле), назву предмету (знаків пунктуації не містить), прізвище студента, ім'я студента, номер заліковки, оцінку в балах за 100 бальною шкалою.

Білі символи навколо роздільників не допускаються.

Приклад рядка файлу



\begin{verbatim}
13:algebra:surname:name:123456:77
\end{verbatim}


Нижче наведено код, який має залишити у файлі в кожному рядку тільки номер заліковки та предмет.

Рядок з прикладу має бути замінений на такий



\begin{verbatim}
123456:algebra
\end{verbatim}




Код:



\begin{verbatim}
regex = re.compile(r'XXX')
with open('1.txt', encoding='utf8') as f:
    lst = f.readlines()
for i in range(len(lst)):
    lst[i] = re.sub(regex, r'YYY', lst[i])
with open('1.txt', 'w', encoding='utf8') as f:
    f.writelines(lst)
\end{verbatim}









Який рядок треба записати на місце \kw{XXX} ?\par
\vspace{5mm}


Який рядок треба записати на місце \kw{YYY} ?\par


\newpage

{\begin {center}\textbf{ 261}\end {center}}\par
Файл \code{'1.txt'} містить журнал з рядками, в яких через двокрапку записано поряд\-ко\-вий номер запису (ціле), ім'я користувача (складається з латиниці), час події (фор\-мат як у прикладі), тип події (ERROR, WARNING, INFO, STATUS), код події (ціле).

Білі символи навколо роздільників не допускаються.

Приклад рядка журналу



\begin{verbatim}
12:username:2022-05-05 13-26:ERROR:404
\end{verbatim}


Нижче наведено код, який шукає усі попередження (подія WARNING), що зустрі\-ча\-ються в журналі і відбулися під час обідньої перерви деканату (з 13:00 по 13:59 включно), та зберігає їх коди в змінній \kw{codes}.



\begin{verbatim}
codes = set()
regex = re.compile(r'XXX')
with open('1.txt', encoding='utf8') as f:
    for s in f:
        res = re.fullmatch(regex, s)
        if not res: continue
        s = ''
        codes.add(r'YYY')
\end{verbatim}







Який рядок треба записати на місце \kw{XXX} ?\par
\vspace{5mm}


Який рядок треба записати на місце \kw{YYY} ?\par


\newpage

{\begin {center}\textbf{ 262}\end {center}}\par
Файл \code{'1.txt'} містить журнал з рядками, в яких через кому записано поряд\-ко\-вий номер запису в журналі (ціле), тип події (ERROR, WARNING, INFO, STATUS), ім'я користувача (складається з латиниці), час події (фор\-мат як у прикладі), код події (ціле).

Білі символи навколо роздільників не допускаються.

Приклад рядка файлу



\begin{verbatim}
13,ERROR,username,2022-05-05 13:26,404
\end{verbatim}


Нижче наведено код, який шукає усі коди помилок (подія ERROR), що зустріча\-ють\-ся в журналі, та зберігає їх у змінній \kw{codes}.



\begin{verbatim}
codes = set()
regex = re.compile(r'XXX')
with open('1.txt', encoding='utf8') as f:
    for s in f:
        res = re.fullmatch(regex, s)
        if not res: continue
        s = ''
        codes.add(r'YYY')
\end{verbatim}





Який рядок треба записати на місце \kw{XXX} ?\par
\vspace{5mm}


Який рядок треба записати на місце \kw{YYY} ?\par


\newpage

{\begin {center}\textbf{ 263}\end {center}}\par
Файл \code{'1.txt'} містить журнал з рядками, в яких через кому записано поряд\-ко\-вий номер запису (ціле), ім'я користувача (складається з латиниці), час події (фор\-мат як у прикладі), тип події (ERROR, WARNING, INFO, STATUS), код події (ціле).

Білі символи навколо роздільників не допускаються.

Приклад рядка журналу



\begin{verbatim}
12,username,2022-05-05 13:26,ERROR,404
\end{verbatim}


Нижче наведено код, який шукає усі попередження (подія WARNING), що зустрі\-ча\-ються в журналі і відбулися під час обідньої перерви деканату (з 13,00 по 13,59 включно), та зберігає їх коди в змінній \kw{codes}.



\begin{verbatim}
codes = set()
regex = re.compile(r'XXX')
with open('1.txt', encoding='utf8') as f:
    for s in f:
        res = re.fullmatch(regex, s)
        if not res: continue
        s = ''
        codes.add(r'YYY')
\end{verbatim}







Який рядок треба записати на місце \kw{XXX} ?\par
\vspace{5mm}


Який рядок треба записати на місце \kw{YYY} ?\par


\newpage

{\begin {center}\textbf{ 264}\end {center}}\par
Файл \code{'1.txt'} містить журнал з рядками, в яких через кому записано поряд\-ко\-вий номер запису в журналі (ціле), тип події (ERROR, WARNING, INFO, STATUS), ім'я користувача (складається з латиниці), час події (фор\-мат як у прикладі), код події (ціле).

Білі символи навколо роздільників не допускаються.

Приклад рядка файлу



\begin{verbatim}
13,ERROR,username,2022-05-05 13:26,404
\end{verbatim}


Нижче наведено код, який має замінити у файлі імена користувачів на рядок \code{unknown}.



\begin{verbatim}
regex = re.compile(r'XXX')
with open('1.txt', encoding='utf8') as f:
    lst = f.readlines()
for i in range(len(lst)):
    lst[i] = re.sub(regex, r'YYY', lst[i])
with open('1.txt', 'w', encoding='utf8') as f:
    f.writelines(lst)
\end{verbatim}







Який рядок треба записати на місце \kw{XXX} ?\par
\vspace{5mm}


Який рядок треба записати на місце \kw{YYY} ?\par


\newpage

{\begin {center}\textbf{ 265}\end {center}}\par
Рядки журналу через двокрапку містять ім'я користувача (складається з латиниці), час події (фор\-мат як у прикладі), тип події (ERROR, WARNING, INFO, STATUS), код події (ціле).

Білі символи навколо роздільників не допускаються.

Приклад такого рядка присвоєно змінній \kw{s}.



\begin{verbatim}
s = 'username:2022-05-05 13-26:ERROR:404'
\end{verbatim}


Нижче наведено код, за виконання якого значенням змінної \kw{out} має стати код події.

Для наведеного прикладу значенням змінної \kw{out} має стати рядок \code{404}



\begin{verbatim}
res = re.fullmatch(r'XXX', s)
s = ''
out = r'YYY'
\end{verbatim}


Код має працювати не тільки для конкретного рядка з прикладу, але й для кожного рядка з журналу.




Який рядок треба записати на місце \kw{XXX} ?\par
\vspace{5mm}


Який рядок треба записати на місце \kw{YYY} ?\par


\newpage

{\begin {center}\textbf{ 266}\end {center}}\par
Рядки відомості через двокрапку містять порядковий номер (ціле), пріз\-ви\-ще студен\-та, ім'я студента, номер заліковки, оцінку в балах за 100 бальною шкалою.

Білі символи навколо роздільників не допускаються.

Приклад такого рядка присвоєно змінній \kw{s}.



\begin{verbatim}
s = '13:surname:name:123456:77'
\end{verbatim}


Нижче наведено код, за виконання якого для рядка \kw{s} значенням змінної \kw{out} має стати номер заліковки.

Для наведеного прикладу значенням змінної \kw{out} має стати рядок \code{123456} .



\begin{verbatim}
res = re.fullmatch(r'XXX', s)
s = ''
out = r'YYY'
\end{verbatim}


Код має працювати не тільки для конкретного рядка з прикладу, але й для кожного рядка з відомості.




Який рядок треба записати на місце \kw{XXX} ?\par
\vspace{5mm}


Який рядок треба записати на місце \kw{YYY} ?\par


\newpage

{\begin {center}\textbf{ 267}\end {center}}\par
Файл \code{'1.txt'} містить журнал з рядками, в яких через двокрапку записано поряд\-ко\-вий номер запису в журналі (ціле), код події (ціле), тип події (ERROR,\\ WARNING, INFO, STATUS), ім'я користувача (складається з латиниці), час події (фор\-мат як у прикладі), опис події.

Білі символи навколо роздільників не допускаються.

Приклад рядка файлу



\begin{verbatim}
13:404:ERROR:username:2022-05-05 13:26:something happens
\end{verbatim}


Нижче наведено код, який має в кожному рядку файли переставити місцями код та тип події.

Рядок з прикладу має бути замінений на такий



\begin{verbatim}
13:ERROR:404:username:2022-05-05 13:26:something happens
\end{verbatim}




Код:



\begin{verbatim}
regex = re.compile(r'XXX')
with open('1.txt', encoding='utf8') as f:
    lst = f.readlines()
for i in range(len(lst)):
    lst[i] = re.sub(regex, r'YYY', lst[i])
with open('1.txt', 'w', encoding='utf8') as f:
    f.writelines(lst)
\end{verbatim}







Який рядок треба записати на місце \kw{XXX} ?\par
\vspace{5mm}


Який рядок треба записати на місце \kw{YYY} ?\par


\newpage

{\begin {center}\textbf{ 268}\end {center}}\par
Файл \code{'1.txt'} містить відомість з рядками, в яких через двокрапку записано поряд\-ковий номер (ціле), назву предмету (знаків пунктуації не містить), прізвище студента, ім'я студента, номер заліковки, оцінку в балах за 100 бальною шкалою.

Білі символи навколо роздільників не допускаються.

Приклад рядка файлу



\begin{verbatim}
13:algebra:surname:name:123456:77
\end{verbatim}


Нижче наведено код, який шукає усі предмети, який складав студент із заліковкою \code{777}, та зберігає їх типи в змінній \kw{codes}.



\begin{verbatim}
codes = set()
regex = re.compile(r'XXX')
with open('1.txt', encoding='utf8') as f:
    for s in f:
        res = re.fullmatch(regex, s)
        if not res: continue
        s = ''
        codes.add(r'YYY')
\end{verbatim}





Який рядок треба записати на місце \kw{XXX} ?\par
\vspace{5mm}


Який рядок треба записати на місце \kw{YYY} ?\par


\newpage

{\begin {center}\textbf{ 269}\end {center}}\par
Файл \code{'1.txt'} містить відомість з рядками, в яких через двокрапку записа\-но поряд\-ковий номер (ціле), назву предмету (знаків пунктуації не містить), прізвище студента, ім'я студента, номер заліковки, оцінку в балах за 100 бальною шкалою.

Білі символи навколо роздільників не допускаються.

Приклад рядка файлу



\begin{verbatim}
13:algebra:surname:name:123456:77
\end{verbatim}


Нижче наведено код, який має залишити у файлі в кожному рядку тільки номер заліковки, назву предмету та оцінку в балах за 100-бальною шкалою (у заданій послі\-дов\-ності).

Рядок з прикладу має бути замінений на такий



\begin{verbatim}
123456:algebra:77
\end{verbatim}




Код:



\begin{verbatim}
regex = re.compile(r'XXX')
with open('1.txt', encoding='utf8') as f:
    lst = f.readlines()
for i in range(len(lst)):
    lst[i] = re.sub(regex, r'YYY', lst[i])
with open('1.txt', 'w', encoding='utf8') as f:
    f.writelines(lst)
\end{verbatim}







Який рядок треба записати на місце \kw{XXX} ?\par
\vspace{5mm}


Який рядок треба записати на місце \kw{YYY} ?\par


\newpage

{\begin {center}\textbf{ 270}\end {center}}\par
Файл \code{'1.txt'} містить журнал з рядками, в яких через кому записано поряд\-ко\-вий номер запису в журналі (ціле), код події (ціле), тип події (ERROR, \\WARNING, INFO, STATUS), ім'я користувача (складається з латиниці), час події (фор\-мат як у прикладі), опис події.

Білі символи навколо роздільників не допускаються.

Приклад рядка файлу



\begin{verbatim}
13,404,ERROR,username,2022-05-05 13:26,something happens
\end{verbatim}


Нижче наведено код, який має залишити у файлі в кожному рядку тільки код події та ім'я користувача. Рядок з прикладу має бути замінений на такий



\begin{verbatim}
404,username
\end{verbatim}




Код:



\begin{verbatim}
regex = re.compile(r'XXX')
with open('1.txt', encoding='utf8') as f:
    lst = f.readlines()
for i in range(len(lst)):
    lst[i] = re.sub(regex, r'YYY', lst[i])
with open('1.txt', 'w', encoding='utf8') as f:
    f.writelines(lst)
\end{verbatim}







Який рядок треба записати на місце \kw{XXX} ?\par
\vspace{5mm}


Який рядок треба записати на місце \kw{YYY} ?\par


\newpage

{\begin {center}\textbf{ 271}\end {center}}\par
Файл \code{'1.txt'} містить відомість з рядками, в яких через двокрапку записано поряд\-ковий номер (ціле), назву предмету (знаків пунктуації не містить), прізвище студента, ім'я студента, номер заліковки, оцінку в балах за 100 бальною шкалою.

Білі символи навколо роздільників не допускаються.

Приклад рядка файлу



\begin{verbatim}
13:algebra:surname:name:123456:77
\end{verbatim}


Нижче наведено код, який має в кожному рядку файлу переставити місцями назву предмету та номер заліковки.

Рядок з прикладу має бути замінений на такий



\begin{verbatim}
13:123456:surname:name:algebra:77
\end{verbatim}




Код:



\begin{verbatim}
regex = re.compile(r'XXX')
with open('1.txt', encoding='utf8') as f:
    lst = f.readlines()
for i in range(len(lst)):
    lst[i] = re.sub(regex, r'YYY', lst[i])
with open('1.txt', 'w', encoding='utf8') as f:
    f.writelines(lst)
\end{verbatim}







Який рядок треба записати на місце \kw{XXX} ?\par
\vspace{5mm}


Який рядок треба записати на місце \kw{YYY} ?\par


\newpage

{\begin {center}\textbf{ 272}\end {center}}\par
Файл \code{'1.txt'} містить журнал з рядками, в яких через кому записано поряд\-ко\-вий номер запису в журналі (ціле), код події (ціле), тип події (ERROR,\\ WARNING, INFO, STATUS), ім'я користувача (складається з латиниці), час події (фор\-мат як у прикладі), опис події.

Білі символи навколо роздільників не допускаються.

Приклад рядка файлу



\begin{verbatim}
13,404,ERROR,username,2022-05-05 13:26,something happens
\end{verbatim}


Нижче наведено код, який має в кожному рядку файли переставити місцями код та тип події.

Рядок з прикладу має бути замінений на такий



\begin{verbatim}
13,ERROR,404,username,2022-05-05 13:26,something happens
\end{verbatim}




Код:



\begin{verbatim}
regex = re.compile(r'XXX')
with open('1.txt', encoding='utf8') as f:
    lst = f.readlines()
for i in range(len(lst)):
    lst[i] = re.sub(regex, r'YYY', lst[i])
with open('1.txt', 'w', encoding='utf8') as f:
    f.writelines(lst)
\end{verbatim}







Який рядок треба записати на місце \kw{XXX} ?\par
\vspace{5mm}


Який рядок треба записати на місце \kw{YYY} ?\par


\newpage

{\begin {center}\textbf{ 273}\end {center}}\par
Файл \code{'1.txt'} містить журнал з рядками, в яких через двокрапку записано поряд\-ко\-вий номер запису в журналі (ціле), код події (ціле), тип події (ERROR, \\WARNING, INFO, STATUS), ім'я користувача (складається з латиниці), час події (фор\-мат як у прикладі), опис події.

Білі символи навколо роздільників не допускаються.

Приклад рядка файлу



\begin{verbatim}
13:404:ERROR:username:2022-05-05 13-26:something happens
\end{verbatim}


Нижче наведено код, який має залишити у файлі в кожному рядку тільки код події та ім'я користувача. Рядок з прикладу має бути замінений на такий



\begin{verbatim}
404:username
\end{verbatim}




Код:



\begin{verbatim}
regex = re.compile(r'XXX')
with open('1.txt', encoding='utf8') as f:
    lst = f.readlines()
for i in range(len(lst)):
    lst[i] = re.sub(regex, r'YYY', lst[i])
with open('1.txt', 'w', encoding='utf8') as f:
    f.writelines(lst)
\end{verbatim}







Який рядок треба записати на місце \kw{XXX} ?\par
\vspace{5mm}


Який рядок треба записати на місце \kw{YYY} ?\par


\newpage

{\begin {center}\textbf{ 274}\end {center}}\par
Файл \code{'1.txt'} містить журнал з рядками, в яких через двокрапку записано поряд\-ко\-вий номер запису в журналі (ціле), тип події (ERROR, WARNING, INFO, STATUS), ім'я користувача (складається з латиниці), час події (фор\-мат як у прикладі), код події (ціле).

Білі символи навколо роздільників не допускаються.

Приклад рядка файлу



\begin{verbatim}
13:ERROR:username:2022-05-05 13-26:404
\end{verbatim}


Нижче наведено код, який шукає усі коди помилок (подія ERROR), що зустріча\-ють\-ся в журналі, та зберігає їх у змінній \kw{codes}.



\begin{verbatim}
codes = set()
regex = re.compile(r'XXX')
with open('1.txt', encoding='utf8') as f:
    for s in f:
        res = re.fullmatch(regex, s)
        if not res: continue
        s = ''
        codes.add(r'YYY')
\end{verbatim}





Який рядок треба записати на місце \kw{XXX} ?\par
\vspace{5mm}


Який рядок треба записати на місце \kw{YYY} ?\par


\newpage

{\begin {center}\textbf{ 275}\end {center}}\par
Файл \code{'1.txt'} містить відомість з рядками, в яких через двокрапку записано поряд\-ковий номер (ціле), назву предмету (знаків пунктуації не містить), прізвище студента, ім'я студента, номер заліковки, оцінку в балах за 100 бальною шкалою.

Білі символи навколо роздільників не допускаються.

Приклад рядка файлу



\begin{verbatim}
13:algebra:surname:name:123456:77
\end{verbatim}


Нижче наведено код, який шукає усі назви предметів, що зустрічаються у файлі, та зберігає їх у змінній \id{codes}.



\begin{verbatim}
codes = set()
regex = re.compile(r'XXX')
with open('1.txt', encoding='utf8') as f:
    for s in f:
        res = re.fullmatch(regex, s)
        if not res: continue
        s = ''
        codes.add(r'YYY')
\end{verbatim}





Який рядок треба записати на місце \kw{XXX} ?\par
\vspace{5mm}


Який рядок треба записати на місце \kw{YYY} ?\par


\newpage

{\begin {center}\textbf{ 276}\end {center}}\par
Файл \code{'1.txt'} містить відомість з рядками, в яких через двокрапку записано поряд\-ковий номер (ціле), назву предмету (знаків пунктуації не містить), прізвище студента, ім'я студента, номер заліковки, оцінку в балах за 100 бальною шкалою.

Білі символи навколо роздільників не допускаються.

Приклад рядка файлу



\begin{verbatim}
13:algebra:surname:name:123456:77
\end{verbatim}


Нижче наведено код, який має залишити у файлі в кожному рядку тільки номер заліковки та предмет.

Рядок з прикладу має бути замінений на такий



\begin{verbatim}
123456:algebra
\end{verbatim}




Код:



\begin{verbatim}
regex = re.compile(r'XXX')
with open('1.txt', encoding='utf8') as f:
    lst = f.readlines()
for i in range(len(lst)):
    lst[i] = re.sub(regex, r'YYY', lst[i])
with open('1.txt', 'w', encoding='utf8') as f:
    f.writelines(lst)
\end{verbatim}









Який рядок треба записати на місце \kw{XXX} ?\par
\vspace{5mm}


Який рядок треба записати на місце \kw{YYY} ?\par


\newpage

{\begin {center}\textbf{ 277}\end {center}}\par
Рядки журналу через двокрапку містять ім'я користувача (складається з латиниці), час події (фор\-мат як у прикладі), тип події (ERROR, WARNING, INFO, STATUS), код події (ціле).

Білі символи навколо роздільників не допускаються.

Приклад такого рядка присвоєно змінній \kw{s}.



\begin{verbatim}
s = 'username:2022-05-05 13-26:ERROR:404'
\end{verbatim}


Нижче наведено код, за виконання якого значенням змінної \kw{out} має стати код події.

Для наведеного прикладу значенням змінної \kw{out} має стати рядок \code{404}



\begin{verbatim}
res = re.fullmatch(r'XXX', s)
s = ''
out = r'YYY'
\end{verbatim}


Код має працювати не тільки для конкретного рядка з прикладу, але й для кожного рядка з журналу.




Який рядок треба записати на місце \kw{XXX} ?\par
\vspace{5mm}


Який рядок треба записати на місце \kw{YYY} ?\par


\newpage

{\begin {center}\textbf{ 278}\end {center}}\par
Файл \code{'1.txt'} містить відомість з рядками, в яких через кому записа\-но назву предмету (знаків пунктуації не містить), прізвище студента, ім'я студента, номер залі\-ков\-ки, оцінку в балах за 100 бальною шкалою.

Білі символи навколо роздільників не допускаються.

Приклад рядка файлу



\begin{verbatim}
algebra,surname,name,123456,77
\end{verbatim}


Нижче наведено код, який має замінити у файлі усі номери заліковок на рядок \code{***}.



\begin{verbatim}
regex = re.compile(r'XXX')
with open('1.txt', encoding='utf8') as f:
    lst = f.readlines()
for i in range(len(lst)):
    lst[i] = re.sub(regex, r'YYY', lst[i])
with open('1.txt', 'w', encoding='utf8') as f:
    f.writelines(lst)
\end{verbatim}







Який рядок треба записати на місце \kw{XXX} ?\par
\vspace{5mm}


Який рядок треба записати на місце \kw{YYY} ?\par


\newpage

{\begin {center}\textbf{ 279}\end {center}}\par
Рядки журналу через кому містять ім'я користувача (складається з латиниці), час події (фор\-мат як у прикладі), тип події (ERROR, WARNING, INFO, STATUS), код події (ціле).

Білі символи навколо роздільників не допускаються.

Приклад такого рядка присвоєно змінній \kw{s}.



\begin{verbatim}
s = 'username,2022-05-05 13:26,ERROR,404'
\end{verbatim}


Нижче наведено код, за виконання якого значенням змінної \kw{out} має стати код події.

Для наведеного прикладу значенням змінної \kw{out} має стати рядок \code{404}



\begin{verbatim}
res = re.fullmatch(r'XXX', s)
s = ''
out = r'YYY'
\end{verbatim}


Код має працювати не тільки для конкретного рядка з прикладу, але й для кожного рядка з журналу.




Який рядок треба записати на місце \kw{XXX} ?\par
\vspace{5mm}


Який рядок треба записати на місце \kw{YYY} ?\par


\newpage

{\begin {center}\textbf{ 280}\end {center}}\par
Файл \code{'1.txt'} містить журнал з рядками, в яких через кому записано поряд\-ко\-вий номер запису в журналі (ціле), тип події (ERROR, WARNING, INFO, STATUS), ім'я користувача (складається з латиниці), час події (фор\-мат як у прикладі), опис події.

Білі символи навколо роздільників не допускаються.

Приклад рядка файлу



\begin{verbatim}
13,ERROR,username,2022-05-05 13:26,something happens
\end{verbatim}


Нижче наведено код, який має залишити у файлі в кожному рядку тільки тип події та її опис.

Рядок з прикладу має бути замінений на такий



\begin{verbatim}
ERROR,something happens
\end{verbatim}




Код:



\begin{verbatim}
regex = re.compile(r'XXX')
with open('1.txt', encoding='utf8') as f:
    lst = f.readlines()
for i in range(len(lst)):
    lst[i] = re.sub(regex, r'YYY', lst[i])
with open('1.txt', 'w', encoding='utf8') as f:
    f.writelines(lst)
\end{verbatim}









Який рядок треба записати на місце \kw{XXX} ?\par
\vspace{5mm}


Який рядок треба записати на місце \kw{YYY} ?\par


\newpage

{\begin {center}\textbf{ 281}\end {center}}\par
Файл \code{'1.txt'} містить журнал з рядками, в яких через двокрапку записано поряд\-ко\-вий номер запису в журналі (ціле), код події (ціле), тип події (ERROR,\\ WARNING, INFO, STATUS), ім'я користувача (складається з латиниці), час події (фор\-мат як у прикладі), опис події.

Білі символи навколо роздільників не допускаються.

Приклад рядка файлу



\begin{verbatim}
13:404:ERROR:username:2022-05-05 13:26:something happens
\end{verbatim}


Нижче наведено код, який має в кожному рядку файли переставити місцями код та тип події.

Рядок з прикладу має бути замінений на такий



\begin{verbatim}
13:ERROR:404:username:2022-05-05 13:26:something happens
\end{verbatim}




Код:



\begin{verbatim}
regex = re.compile(r'XXX')
with open('1.txt', encoding='utf8') as f:
    lst = f.readlines()
for i in range(len(lst)):
    lst[i] = re.sub(regex, r'YYY', lst[i])
with open('1.txt', 'w', encoding='utf8') as f:
    f.writelines(lst)
\end{verbatim}







Який рядок треба записати на місце \kw{XXX} ?\par
\vspace{5mm}


Який рядок треба записати на місце \kw{YYY} ?\par


\newpage

{\begin {center}\textbf{ 282}\end {center}}\par
Файл \code{'1.txt'} містить журнал з рядками, в яких через двокрапку записано поряд\-ко\-вий номер запису в журналі (ціле), тип події (ERROR, WARNING, INFO, STATUS), ім'я користувача (складається з латиниці), час події (фор\-мат як у прикладі), код події (ціле).

Білі символи навколо роздільників не допускаються.

Приклад рядка файлу



\begin{verbatim}
13:ERROR:username:2022-05-05 13-26:404
\end{verbatim}


Нижче наведено код, який має замінити у файлі імена користувачів на рядок \code{unknown}.



\begin{verbatim}
regex = re.compile(r'XXX')
with open('1.txt', encoding='utf8') as f:
    lst = f.readlines()
for i in range(len(lst)):
    lst[i] = re.sub(regex, r'YYY', lst[i])
with open('1.txt', 'w', encoding='utf8') as f:
    f.writelines(lst)
\end{verbatim}







Який рядок треба записати на місце \kw{XXX} ?\par
\vspace{5mm}


Який рядок треба записати на місце \kw{YYY} ?\par


\newpage

{\begin {center}\textbf{ 283}\end {center}}\par
Файл \code{'1.txt'} містить відомість з рядками, в яких через двокрапку записано поряд\-ковий номер (ціле), назву предмету (знаків пунктуації не містить), прізвище студента, ім'я студента, номер заліковки, оцінку в балах за 100 бальною шкалою.

Білі символи навколо роздільників не допускаються.

Приклад рядка файлу



\begin{verbatim}
13:algebra:surname:name:123456:77
\end{verbatim}


Нижче наведено код, який має в кожному рядку файлу переставити місцями назву предмету та номер заліковки.

Рядок з прикладу має бути замінений на такий



\begin{verbatim}
13:123456:surname:name:algebra:77
\end{verbatim}




Код:



\begin{verbatim}
regex = re.compile(r'XXX')
with open('1.txt', encoding='utf8') as f:
    lst = f.readlines()
for i in range(len(lst)):
    lst[i] = re.sub(regex, r'YYY', lst[i])
with open('1.txt', 'w', encoding='utf8') as f:
    f.writelines(lst)
\end{verbatim}







Який рядок треба записати на місце \kw{XXX} ?\par
\vspace{5mm}


Який рядок треба записати на місце \kw{YYY} ?\par


\newpage

{\begin {center}\textbf{ 284}\end {center}}\par
Файл \code{'1.txt'} містить журнал з рядками, в яких через кому записано поряд\-ко\-вий номер запису в журналі (ціле), тип події (ERROR, WARNING, INFO, STATUS), ім'я користувача (складається з латиниці), час події (фор\-мат як у прикладі), код події (ціле).

Білі символи навколо роздільників не допускаються.

Приклад рядка файлу



\begin{verbatim}
13,ERROR,username,2022-05-05 13:26,404
\end{verbatim}


Нижче наведено код, який шукає усі коди помилок (подія ERROR), що зустріча\-ють\-ся в журналі, та зберігає їх у змінній \kw{codes}.



\begin{verbatim}
codes = set()
regex = re.compile(r'XXX')
with open('1.txt', encoding='utf8') as f:
    for s in f:
        res = re.fullmatch(regex, s)
        if not res: continue
        s = ''
        codes.add(r'YYY')
\end{verbatim}





Який рядок треба записати на місце \kw{XXX} ?\par
\vspace{5mm}


Який рядок треба записати на місце \kw{YYY} ?\par


\newpage

{\begin {center}\textbf{ 285}\end {center}}\par
Файл \code{'1.txt'} містить відомість з рядками, в яких через кому записано поряд\-ковий номер (ціле), назву предмету (знаків пунктуації не містить), прізвище студента, ім'я студента, номер заліковки, оцінку в балах за 100 бальною шкалою.

Білі символи навколо роздільників не допускаються.

Приклад рядка файлу



\begin{verbatim}
13,algebra,surname,name,123456,77
\end{verbatim}


Нижче наведено код, який має залишити у файлі в кожному рядку тільки номер заліковки та предмет.

Рядок з прикладу має бути замінений на такий



\begin{verbatim}
123456,algebra
\end{verbatim}




Код



\begin{verbatim}
regex = re.compile(r'XXX')
with open('1.txt', encoding='utf8') as f:
    lst = f.readlines()
for i in range(len(lst)):
    lst[i] = re.sub(regex, r'YYY', lst[i])
with open('1.txt', 'w', encoding='utf8') as f:
    f.writelines(lst)
\end{verbatim}









Який рядок треба записати на місце \kw{XXX} ?\par
\vspace{5mm}


Який рядок треба записати на місце \kw{YYY} ?\par


\newpage

{\begin {center}\textbf{ 286}\end {center}}\par
Файл \code{'1.txt'} містить відомість з рядками, в яких через двокрапку записано поряд\-ковий номер (ціле), назву предмету (знаків пунктуації не містить), прізвище студента, ім'я студента, номер заліковки, оцінку в балах за 100 бальною шкалою.

Білі символи навколо роздільників не допускаються.

Приклад рядка файлу



\begin{verbatim}
13:algebra:surname:name:123456:77
\end{verbatim}


Нижче наведено код, який шукає усі предмети, який складав студент із заліковкою \code{777}, та зберігає їх типи в змінній \kw{codes}.



\begin{verbatim}
codes = set()
regex = re.compile(r'XXX')
with open('1.txt', encoding='utf8') as f:
    for s in f:
        res = re.fullmatch(regex, s)
        if not res: continue
        s = ''
        codes.add(r'YYY')
\end{verbatim}





Який рядок треба записати на місце \kw{XXX} ?\par
\vspace{5mm}


Який рядок треба записати на місце \kw{YYY} ?\par


\newpage

{\begin {center}\textbf{ 287}\end {center}}\par
Файл \code{'1.txt'} містить журнал з рядками, в яких через двокрапку записано поряд\-ко\-вий номер запису (ціле), ім'я користувача (складається з латиниці), час події (фор\-мат як у прикладі), тип події (ERROR, WARNING, INFO, STATUS), код події (ціле).

Білі символи навколо роздільників не допускаються.

Приклад рядка журналу



\begin{verbatim}
12:username:2022-05-05 13-26:ERROR:404
\end{verbatim}


Нижче наведено код, який шукає усі попередження (подія WARNING), що зустрі\-ча\-ються в журналі і відбулися під час обідньої перерви деканату (з 13:00 по 13:59 включно), та зберігає їх коди в змінній \kw{codes}.



\begin{verbatim}
codes = set()
regex = re.compile(r'XXX')
with open('1.txt', encoding='utf8') as f:
    for s in f:
        res = re.fullmatch(regex, s)
        if not res: continue
        s = ''
        codes.add(r'YYY')
\end{verbatim}







Який рядок треба записати на місце \kw{XXX} ?\par
\vspace{5mm}


Який рядок треба записати на місце \kw{YYY} ?\par


\newpage

{\begin {center}\textbf{ 288}\end {center}}\par
Файл \code{'1.txt'} містить журнал з рядками, в яких через кому записано поряд\-ко\-вий номер запису (ціле), ім'я користувача (складається з латиниці), час події (фор\-мат як у прикладі), тип події (ERROR, WARNING, INFO, STATUS), код події (ціле).

Білі символи навколо роздільників не допускаються.

Приклад рядка журналу



\begin{verbatim}
12,username,2022-05-05 13:26,ERROR,404
\end{verbatim}


Нижче наведено код, який шукає усі події, що відбулись\\ для користувача \code{anonymous}, та зберігає їх типи в змінній \kw{codes}.



\begin{verbatim}
codes = set()
regex = re.compile(r'XXX')
with open('1.txt', encoding='utf8') as f:
    for s in f:
        res = re.fullmatch(regex, s)
        if not res: continue
        s = ''
        codes.add(r'YYY')
\end{verbatim}





Який рядок треба записати на місце \kw{XXX} ?\par
\vspace{5mm}


Який рядок треба записати на місце \kw{YYY} ?\par


\newpage

{\begin {center}\textbf{ 289}\end {center}}\par
Файл \code{'1.txt'} містить журнал з рядками, в яких через кому записано поряд\-ко\-вий номер запису в журналі (ціле), тип події (ERROR, WARNING, INFO, STATUS), ім'я користувача (складається з латиниці), час події (фор\-мат як у прикладі), код події (ціле).

Білі символи навколо роздільників не допускаються.

Приклад рядка файлу



\begin{verbatim}
13,ERROR,username,2022-05-05 13:26,404
\end{verbatim}


Нижче наведено код, який має замінити у файлі імена користувачів на рядок \code{unknown}.



\begin{verbatim}
regex = re.compile(r'XXX')
with open('1.txt', encoding='utf8') as f:
    lst = f.readlines()
for i in range(len(lst)):
    lst[i] = re.sub(regex, r'YYY', lst[i])
with open('1.txt', 'w', encoding='utf8') as f:
    f.writelines(lst)
\end{verbatim}







Який рядок треба записати на місце \kw{XXX} ?\par
\vspace{5mm}


Який рядок треба записати на місце \kw{YYY} ?\par


\newpage

{\begin {center}\textbf{ 290}\end {center}}\par
Рядки відомості через двокрапку містять порядковий номер (ціле), назву предмету (знаків пунктуації не містить), прізвище студента, ім'я студента, оцінку в балах за 100 бальною шкалою.

Білі символи навколо роздільників не допускаються.

Приклад такого рядка присвоєно змінній \kw{s}.



\begin{verbatim}
s = '13:algebra:surname:name:77'
\end{verbatim}


Нижче наведено код, за виконання якого значенням змінної \kw{out} має стати прізвище студента.

Для наведеного прикладу значенням змінної \kw{out} має стати рядок \code{surname}



\begin{verbatim}
res = re.fullmatch(r'XXX', s)
s = ''
out = r'YYY'
\end{verbatim}


Код має працювати не тільки для конкретного рядка з прикладу, але й для кожного рядка з відомості.




Який рядок треба записати на місце \kw{XXX} ?\par
\vspace{5mm}


Який рядок треба записати на місце \kw{YYY} ?\par


\newpage

{\begin {center}\textbf{ 291}\end {center}}\par
Файл \code{'1.txt'} містить журнал з рядками, в яких через двокрапку записано поряд\-ко\-вий номер запису в журналі (ціле), тип події (ERROR, WARNING, INFO, STATUS), ім'я користувача (складається з латиниці), час події (фор\-мат як у прикладі), опис події.

Білі символи навколо роздільників не допускаються.

Приклад рядка файлу



\begin{verbatim}
13:ERROR:username:2022-05-05 13-26:something happens
\end{verbatim}


Нижче наведено код, який має залишити у файлі в кожному рядку тільки тип події та її опис.

Рядок з прикладу має бути замінений на такий



\begin{verbatim}
ERROR:something happens
\end{verbatim}




Код:



\begin{verbatim}
regex = re.compile(r'XXX')
with open('1.txt', encoding='utf8') as f:
    lst = f.readlines()
for i in range(len(lst)):
    lst[i] = re.sub(regex, r'YYY', lst[i])
with open('1.txt', 'w', encoding='utf8') as f:
    f.writelines(lst)
\end{verbatim}









Який рядок треба записати на місце \kw{XXX} ?\par
\vspace{5mm}


Який рядок треба записати на місце \kw{YYY} ?\par


\newpage

{\begin {center}\textbf{ 292}\end {center}}\par
Рядки журналу через кому містять тип події (ERROR, WARNING, INFO, STATUS), ім'я користувача (складається з латиниці), час події (фор\-мат як у прикладі), код події (ціле).

Білі символи навколо роздільників не допускаються.

Приклад такого рядка присвоєно змінній \kw{s}.



\begin{verbatim}
s = 'ERROR,username,2022-05-05 13:26,404'
\end{verbatim}


Нижче наведено код, за виконання якого значенням змінної \kw{out} має стати ім'я користувача.

Для наведеного прикладу значенням змінної \kw{out} має стати рядок \code{username}



\begin{verbatim}
res = re.fullmatch(r'XXX', s)
s = ''
out = r'YYY'
\end{verbatim}


Код має працювати не тільки для конкретного рядка з прикладу, але й для кожного рядка з журналу.




Який рядок треба записати на місце \kw{XXX} ?\par
\vspace{5mm}


Який рядок треба записати на місце \kw{YYY} ?\par


\newpage

{\begin {center}\textbf{ 293}\end {center}}\par
Файл \code{'1.txt'} містить відомість з рядками, в яких через кому записа\-но поряд\-ковий номер (ціле), назву предмету (знаків пунктуації не містить), прізвище студента, ім'я студента, номер заліковки, оцінку в балах за 100 бальною шкалою.

Білі символи навколо роздільників не допускаються.

Приклад рядка файлу



\begin{verbatim}
13,algebra,surname,name,123456,77
\end{verbatim}


Нижче наведено код, який має залишити у файлі в кожному рядку тільки номер заліковки, назву предмету та оцінку в балах за 100-бальною шкалою (у заданій послі\-дов\-ності).

Рядок з прикладу має бути замінений на такий



\begin{verbatim}
123456,algebra,77
\end{verbatim}




Код



\begin{verbatim}
regex = re.compile(r'XXX')
with open('1.txt', encoding='utf8') as f:
    lst = f.readlines()
for i in range(len(lst)):
    lst[i] = re.sub(regex, r'YYY', lst[i])
with open('1.txt', 'w', encoding='utf8') as f:
    f.writelines(lst)
\end{verbatim}







Який рядок треба записати на місце \kw{XXX} ?\par
\vspace{5mm}


Який рядок треба записати на місце \kw{YYY} ?\par


\newpage

{\begin {center}\textbf{ 294}\end {center}}\par
Файл \code{'1.txt'} містить журнал з рядками, в яких через двокрапку записано поряд\-ко\-вий номер запису (ціле), ім'я користувача (складається з латиниці), час події (фор\-мат як у прикладі), тип події (ERROR, WARNING, INFO, STATUS), код події (ціле).

Білі символи навколо роздільників не допускаються.

Приклад рядка журналу



\begin{verbatim}
12:username:2022-05-05 13-26:ERROR:404
\end{verbatim}


Нижче наведено код, який шукає усі події, що відбулись\\ для користувача \code{anonymous}, та зберігає їх типи в змінній \kw{codes}.



\begin{verbatim}
codes = set()
regex = re.compile(r'XXX')
with open('1.txt', encoding='utf8') as f:
    for s in f:
        res = re.fullmatch(regex, s)
        if not res: continue
        s = ''
        codes.add(r'YYY')
\end{verbatim}





Який рядок треба записати на місце \kw{XXX} ?\par
\vspace{5mm}


Який рядок треба записати на місце \kw{YYY} ?\par


\newpage

{\begin {center}\textbf{ 295}\end {center}}\par
Рядки журналу через двокрапку містять тип події (ERROR, WARNING, INFO, STATUS), ім'я користувача (складається з латиниці), час події (фор\-мат як у прикладі), код події (ціле).

Білі символи навколо роздільників не допускаються.

Приклад такого рядка присвоєно змінній \kw{s}.



\begin{verbatim}
s = 'ERROR:username:2022-05-05 13-26:404'
\end{verbatim}


Нижче наведено код, за виконання якого значенням змінної \kw{out} має стати ім'я користувача.

Для наведеного прикладу значенням змінної \kw{out} має стати рядок \code{username}



\begin{verbatim}
res = re.fullmatch(r'XXX', s)
s = ''
out = r'YYY'
\end{verbatim}


Код має працювати не тільки для конкретного рядка з прикладу, але й для кожного рядка з журналу.




Який рядок треба записати на місце \kw{XXX} ?\par
\vspace{5mm}


Який рядок треба записати на місце \kw{YYY} ?\par


\newpage

{\begin {center}\textbf{ 296}\end {center}}\par
Файл \code{'1.txt'} містить відомість з рядками, в яких через двокрапку записа\-но назву предмету (знаків пунктуації не містить), прізвище студента, ім'я студента, номер залі\-ков\-ки, оцінку в балах за 100 бальною шкалою.

Білі символи навколо роздільників не допускаються.

Приклад рядка файлу



\begin{verbatim}
algebra:surname:name:123456:77
\end{verbatim}


Нижче наведено код, який має замінити у файлі усі номери заліковок на рядок \code{***}.



\begin{verbatim}
regex = re.compile(r'XXX')
with open('1.txt', encoding='utf8') as f:
    lst = f.readlines()
for i in range(len(lst)):
    lst[i] = re.sub(regex, r'YYY', lst[i])
with open('1.txt', 'w', encoding='utf8') as f:
    f.writelines(lst)
\end{verbatim}







Який рядок треба записати на місце \kw{XXX} ?\par
\vspace{5mm}


Який рядок треба записати на місце \kw{YYY} ?\par


\newpage

{\begin {center}\textbf{ 297}\end {center}}\par
Файл \code{'1.txt'} містить журнал з рядками, в яких через кому записано поряд\-ко\-вий номер запису в журналі (ціле), код події (ціле), тип події (ERROR,\\ WARNING, INFO, STATUS), ім'я користувача (складається з латиниці), час події (фор\-мат як у прикладі), опис події.

Білі символи навколо роздільників не допускаються.

Приклад рядка файлу



\begin{verbatim}
13,404,ERROR,username,2022-05-05 13:26,something happens
\end{verbatim}


Нижче наведено код, який має в кожному рядку файли переставити місцями код та тип події.

Рядок з прикладу має бути замінений на такий



\begin{verbatim}
13,ERROR,404,username,2022-05-05 13:26,something happens
\end{verbatim}




Код:



\begin{verbatim}
regex = re.compile(r'XXX')
with open('1.txt', encoding='utf8') as f:
    lst = f.readlines()
for i in range(len(lst)):
    lst[i] = re.sub(regex, r'YYY', lst[i])
with open('1.txt', 'w', encoding='utf8') as f:
    f.writelines(lst)
\end{verbatim}







Який рядок треба записати на місце \kw{XXX} ?\par
\vspace{5mm}


Який рядок треба записати на місце \kw{YYY} ?\par


\newpage

{\begin {center}\textbf{ 298}\end {center}}\par
Файл \code{'1.txt'} містить відомість з рядками, в яких через кому записано поряд\-ковий номер (ціле), назву предмету (знаків пунктуації не містить), прізвище студента, ім'я студента, номер заліковки, оцінку в балах за 100 бальною шкалою.

Білі символи навколо роздільників не допускаються.

Приклад рядка файлу



\begin{verbatim}
13,algebra,surname,name,123456,77
\end{verbatim}


Нижче наведено код, який шукає усі назви предметів, що зустрічаються у файлі, та зберігає їх у змінній \id{codes}.



\begin{verbatim}
codes = set()
regex = re.compile(r'XXX')
with open('1.txt', encoding='utf8') as f:
    for s in f:
        res = re.fullmatch(regex, s)
        if not res: continue
        s = ''
        codes.add(r'YYY')
\end{verbatim}





Який рядок треба записати на місце \kw{XXX} ?\par
\vspace{5mm}


Який рядок треба записати на місце \kw{YYY} ?\par


\newpage

{\begin {center}\textbf{ 299}\end {center}}\par
Рядки відомості через двокрапку містять порядковий номер (ціле), пріз\-ви\-ще студен\-та, ім'я студента, номер заліковки, оцінку в балах за 100 бальною шкалою.

Білі символи навколо роздільників не допускаються.

Приклад такого рядка присвоєно змінній \kw{s}.



\begin{verbatim}
s = '13:surname:name:123456:77'
\end{verbatim}


Нижче наведено код, за виконання якого для рядка \kw{s} значенням змінної \kw{out} має стати номер заліковки.

Для наведеного прикладу значенням змінної \kw{out} має стати рядок \code{123456} .



\begin{verbatim}
res = re.fullmatch(r'XXX', s)
s = ''
out = r'YYY'
\end{verbatim}


Код має працювати не тільки для конкретного рядка з прикладу, але й для кожного рядка з відомості.




Який рядок треба записати на місце \kw{XXX} ?\par
\vspace{5mm}


Який рядок треба записати на місце \kw{YYY} ?\par


\newpage

{\begin {center}\textbf{ 300}\end {center}}\par
Файл \code{'1.txt'} містить відомість з рядками, в яких через кому записано поряд\-ковий номер (ціле), назву предмету (знаків пунктуації не містить), прізвище студента, ім'я студента, номер заліковки, оцінку в балах за 100 бальною шкалою.

Білі символи навколо роздільників не допускаються.

Приклад рядка файлу



\begin{verbatim}
13,algebra,surname,name,123456,77
\end{verbatim}


Нижче наведено код, який шукає усі предмети, який складав студент із заліковкою \code{777}, та зберігає їх типи в змінній \kw{codes}.



\begin{verbatim}
codes = set()
regex = re.compile(r'XXX')
with open('1.txt', encoding='utf8') as f:
    for s in f:
        res = re.fullmatch(regex, s)
        if not res: continue
        s = ''
        codes.add(r'YYY')
\end{verbatim}





Який рядок треба записати на місце \kw{XXX} ?\par
\vspace{5mm}


Який рядок треба записати на місце \kw{YYY} ?\par


\newpage

{\begin {center}\textbf{ 301}\end {center}}\par
Файл \code{'1.txt'} містить відомість з рядками, в яких через двокрапку записа\-но поряд\-ковий номер (ціле), назву предмету (знаків пунктуації не містить), прізвище студента, ім'я студента, номер заліковки, оцінку в балах за 100 бальною шкалою.

Білі символи навколо роздільників не допускаються.

Приклад рядка файлу



\begin{verbatim}
13:algebra:surname:name:123456:77
\end{verbatim}


Нижче наведено код, який має залишити у файлі в кожному рядку тільки номер заліковки, назву предмету та оцінку в балах за 100-бальною шкалою (у заданій послі\-дов\-ності).

Рядок з прикладу має бути замінений на такий



\begin{verbatim}
123456:algebra:77
\end{verbatim}




Код:



\begin{verbatim}
regex = re.compile(r'XXX')
with open('1.txt', encoding='utf8') as f:
    lst = f.readlines()
for i in range(len(lst)):
    lst[i] = re.sub(regex, r'YYY', lst[i])
with open('1.txt', 'w', encoding='utf8') as f:
    f.writelines(lst)
\end{verbatim}







Який рядок треба записати на місце \kw{XXX} ?\par
\vspace{5mm}


Який рядок треба записати на місце \kw{YYY} ?\par


\newpage

{\begin {center}\textbf{ 302}\end {center}}\par
Файл \code{'1.txt'} містить журнал з рядками, в яких через кому записано поряд\-ко\-вий номер запису (ціле), ім'я користувача (складається з латиниці), час події (фор\-мат як у прикладі), тип події (ERROR, WARNING, INFO, STATUS), код події (ціле).

Білі символи навколо роздільників не допускаються.

Приклад рядка журналу



\begin{verbatim}
12,username,2022-05-05 13:26,ERROR,404
\end{verbatim}


Нижче наведено код, який шукає усі попередження (подія WARNING), що зустрі\-ча\-ються в журналі і відбулися під час обідньої перерви деканату (з 13,00 по 13,59 включно), та зберігає їх коди в змінній \kw{codes}.



\begin{verbatim}
codes = set()
regex = re.compile(r'XXX')
with open('1.txt', encoding='utf8') as f:
    for s in f:
        res = re.fullmatch(regex, s)
        if not res: continue
        s = ''
        codes.add(r'YYY')
\end{verbatim}







Який рядок треба записати на місце \kw{XXX} ?\par
\vspace{5mm}


Який рядок треба записати на місце \kw{YYY} ?\par


\newpage

{\begin {center}\textbf{ 303}\end {center}}\par
Файл \code{'1.txt'} містить відомість з рядками, в яких через кому записано поряд\-ковий номер (ціле), назву предмету (знаків пунктуації не містить), прізвище студента, ім'я студента, номер заліковки, оцінку в балах за 100 бальною шкалою.

Білі символи навколо роздільників не допускаються.

Приклад рядка файлу



\begin{verbatim}
13,algebra,surname,name,123456,77
\end{verbatim}


Нижче наведено код, який має в кожному рядку файлу переставити місцями назву предмету та номер заліковки.

Рядок з прикладу має бути замінений на такий



\begin{verbatim}
13,123456,surname,name,algebra,77
\end{verbatim}




Код



\begin{verbatim}
regex = re.compile(r'XXX')
with open('1.txt', encoding='utf8') as f:
    lst = f.readlines()
for i in range(len(lst)):
    lst[i] = re.sub(regex, r'YYY', lst[i])
with open('1.txt', 'w', encoding='utf8') as f:
    f.writelines(lst)
\end{verbatim}







Який рядок треба записати на місце \kw{XXX} ?\par
\vspace{5mm}


Який рядок треба записати на місце \kw{YYY} ?\par


\newpage

{\begin {center}\textbf{ 304}\end {center}}\par
Файл \code{'1.txt'} містить журнал з рядками, в яких через кому записано поряд\-ко\-вий номер запису в журналі (ціле), код події (ціле), тип події (ERROR, \\WARNING, INFO, STATUS), ім'я користувача (складається з латиниці), час події (фор\-мат як у прикладі), опис події.

Білі символи навколо роздільників не допускаються.

Приклад рядка файлу



\begin{verbatim}
13,404,ERROR,username,2022-05-05 13:26,something happens
\end{verbatim}


Нижче наведено код, який має залишити у файлі в кожному рядку тільки код події та ім'я користувача. Рядок з прикладу має бути замінений на такий



\begin{verbatim}
404,username
\end{verbatim}




Код:



\begin{verbatim}
regex = re.compile(r'XXX')
with open('1.txt', encoding='utf8') as f:
    lst = f.readlines()
for i in range(len(lst)):
    lst[i] = re.sub(regex, r'YYY', lst[i])
with open('1.txt', 'w', encoding='utf8') as f:
    f.writelines(lst)
\end{verbatim}







Який рядок треба записати на місце \kw{XXX} ?\par
\vspace{5mm}


Який рядок треба записати на місце \kw{YYY} ?\par


\newpage

{\begin {center}\textbf{ 305}\end {center}}\par
Рядки відомості через кому містять порядковий номер (ціле), назву предмету (зна\-ків пунктуації не містить), прізвище студента, ім'я студента, оцінку в балах за 100 бальною шкалою.

Білі символи навколо роздільників не допускаються.

Приклад такого рядка присвоєно змінній \kw{s}.



\begin{verbatim}
s = '13,algebra,surname,name,77'
\end{verbatim}


Нижче наведено код, за виконання якого значенням змінної \kw{out} має стати прізвище студента.

Для наведеного прикладу значенням змінної \kw{out} має стати рядок \code{surname}



\begin{verbatim}
res = re.fullmatch(r'XXX', s)
s = ''
out = r'YYY'
\end{verbatim}


Код має працювати не тільки для конкретного рядка з прикладу, але й для кожного рядка з відомості.




Який рядок треба записати на місце \kw{XXX} ?\par
\vspace{5mm}


Який рядок треба записати на місце \kw{YYY} ?\par


\newpage

{\begin {center}\textbf{ 306}\end {center}}\par
Рядки відомості через кому містять порядковий номер (ціле), пріз\-ви\-ще студен\-та, ім'я студента, номер заліковки, оцінку в балах за 100 бальною шкалою.

Білі символи навколо роздільників не допускаються.

Приклад такого рядка присвоєно змінній \kw{s}.



\begin{verbatim}
s = '13,surname,name,123456,77'
\end{verbatim}


Нижче наведено код, за виконання якого для рядка \kw{s} значенням змінної \kw{out} має стати номер заліковки.

Для наведеного прикладу значенням змінної \kw{out} має стати рядок \code{123456} .



\begin{verbatim}
res = re.fullmatch(r'XXX', s)
s = ''
out = r'YYY'
\end{verbatim}


Код має працювати не тільки для конкретного рядка з прикладу, але й для кожного рядка з відомості.




Який рядок треба записати на місце \kw{XXX} ?\par
\vspace{5mm}


Який рядок треба записати на місце \kw{YYY} ?\par


\newpage

{\begin {center}\textbf{ 307}\end {center}}\par
Файл \code{'1.txt'} містить журнал з рядками, в яких через двокрапку записано поряд\-ко\-вий номер запису в журналі (ціле), тип події (ERROR, WARNING, INFO, STATUS), ім'я користувача (складається з латиниці), час події (фор\-мат як у прикладі), опис події.

Білі символи навколо роздільників не допускаються.

Приклад рядка файлу



\begin{verbatim}
13:ERROR:username:2022-05-05 13-26:something happens
\end{verbatim}


Нижче наведено код, який має залишити у файлі в кожному рядку тільки тип події та її опис.

Рядок з прикладу має бути замінений на такий



\begin{verbatim}
ERROR:something happens
\end{verbatim}




Код:



\begin{verbatim}
regex = re.compile(r'XXX')
with open('1.txt', encoding='utf8') as f:
    lst = f.readlines()
for i in range(len(lst)):
    lst[i] = re.sub(regex, r'YYY', lst[i])
with open('1.txt', 'w', encoding='utf8') as f:
    f.writelines(lst)
\end{verbatim}









Який рядок треба записати на місце \kw{XXX} ?\par
\vspace{5mm}


Який рядок треба записати на місце \kw{YYY} ?\par


\newpage

{\begin {center}\textbf{ 308}\end {center}}\par
Файл \code{'1.txt'} містить журнал з рядками, в яких через кому записано поряд\-ко\-вий номер запису в журналі (ціле), код події (ціле), тип події (ERROR,\\ WARNING, INFO, STATUS), ім'я користувача (складається з латиниці), час події (фор\-мат як у прикладі), опис події.

Білі символи навколо роздільників не допускаються.

Приклад рядка файлу



\begin{verbatim}
13,404,ERROR,username,2022-05-05 13:26,something happens
\end{verbatim}


Нижче наведено код, який має в кожному рядку файли переставити місцями код та тип події.

Рядок з прикладу має бути замінений на такий



\begin{verbatim}
13,ERROR,404,username,2022-05-05 13:26,something happens
\end{verbatim}




Код:



\begin{verbatim}
regex = re.compile(r'XXX')
with open('1.txt', encoding='utf8') as f:
    lst = f.readlines()
for i in range(len(lst)):
    lst[i] = re.sub(regex, r'YYY', lst[i])
with open('1.txt', 'w', encoding='utf8') as f:
    f.writelines(lst)
\end{verbatim}







Який рядок треба записати на місце \kw{XXX} ?\par
\vspace{5mm}


Який рядок треба записати на місце \kw{YYY} ?\par


\newpage

{\begin {center}\textbf{ 309}\end {center}}\par
Рядки відомості через двокрапку містять порядковий номер (ціле), пріз\-ви\-ще студен\-та, ім'я студента, номер заліковки, оцінку в балах за 100 бальною шкалою.

Білі символи навколо роздільників не допускаються.

Приклад такого рядка присвоєно змінній \kw{s}.



\begin{verbatim}
s = '13:surname:name:123456:77'
\end{verbatim}


Нижче наведено код, за виконання якого для рядка \kw{s} значенням змінної \kw{out} має стати номер заліковки.

Для наведеного прикладу значенням змінної \kw{out} має стати рядок \code{123456} .



\begin{verbatim}
res = re.fullmatch(r'XXX', s)
s = ''
out = r'YYY'
\end{verbatim}


Код має працювати не тільки для конкретного рядка з прикладу, але й для кожного рядка з відомості.




Який рядок треба записати на місце \kw{XXX} ?\par
\vspace{5mm}


Який рядок треба записати на місце \kw{YYY} ?\par


\newpage

{\begin {center}\textbf{ 310}\end {center}}\par
Файл \code{'1.txt'} містить відомість з рядками, в яких через кому записано поряд\-ковий номер (ціле), назву предмету (знаків пунктуації не містить), прізвище студента, ім'я студента, номер заліковки, оцінку в балах за 100 бальною шкалою.

Білі символи навколо роздільників не допускаються.

Приклад рядка файлу



\begin{verbatim}
13,algebra,surname,name,123456,77
\end{verbatim}


Нижче наведено код, який має в кожному рядку файлу переставити місцями назву предмету та номер заліковки.

Рядок з прикладу має бути замінений на такий



\begin{verbatim}
13,123456,surname,name,algebra,77
\end{verbatim}




Код



\begin{verbatim}
regex = re.compile(r'XXX')
with open('1.txt', encoding='utf8') as f:
    lst = f.readlines()
for i in range(len(lst)):
    lst[i] = re.sub(regex, r'YYY', lst[i])
with open('1.txt', 'w', encoding='utf8') as f:
    f.writelines(lst)
\end{verbatim}







Який рядок треба записати на місце \kw{XXX} ?\par
\vspace{5mm}


Який рядок треба записати на місце \kw{YYY} ?\par


\newpage

{\begin {center}\textbf{ 311}\end {center}}\par
Рядки відомості через кому містять порядковий номер (ціле), пріз\-ви\-ще студен\-та, ім'я студента, номер заліковки, оцінку в балах за 100 бальною шкалою.

Білі символи навколо роздільників не допускаються.

Приклад такого рядка присвоєно змінній \kw{s}.



\begin{verbatim}
s = '13,surname,name,123456,77'
\end{verbatim}


Нижче наведено код, за виконання якого для рядка \kw{s} значенням змінної \kw{out} має стати номер заліковки.

Для наведеного прикладу значенням змінної \kw{out} має стати рядок \code{123456} .



\begin{verbatim}
res = re.fullmatch(r'XXX', s)
s = ''
out = r'YYY'
\end{verbatim}


Код має працювати не тільки для конкретного рядка з прикладу, але й для кожного рядка з відомості.




Який рядок треба записати на місце \kw{XXX} ?\par
\vspace{5mm}


Який рядок треба записати на місце \kw{YYY} ?\par


\newpage

{\begin {center}\textbf{ 312}\end {center}}\par
Файл \code{'1.txt'} містить журнал з рядками, в яких через двокрапку записано поряд\-ко\-вий номер запису в журналі (ціле), код події (ціле), тип події (ERROR,\\ WARNING, INFO, STATUS), ім'я користувача (складається з латиниці), час події (фор\-мат як у прикладі), опис події.

Білі символи навколо роздільників не допускаються.

Приклад рядка файлу



\begin{verbatim}
13:404:ERROR:username:2022-05-05 13:26:something happens
\end{verbatim}


Нижче наведено код, який має в кожному рядку файли переставити місцями код та тип події.

Рядок з прикладу має бути замінений на такий



\begin{verbatim}
13:ERROR:404:username:2022-05-05 13:26:something happens
\end{verbatim}




Код:



\begin{verbatim}
regex = re.compile(r'XXX')
with open('1.txt', encoding='utf8') as f:
    lst = f.readlines()
for i in range(len(lst)):
    lst[i] = re.sub(regex, r'YYY', lst[i])
with open('1.txt', 'w', encoding='utf8') as f:
    f.writelines(lst)
\end{verbatim}







Який рядок треба записати на місце \kw{XXX} ?\par
\vspace{5mm}


Який рядок треба записати на місце \kw{YYY} ?\par


\newpage

{\begin {center}\textbf{ 313}\end {center}}\par
Файл \code{'1.txt'} містить журнал з рядками, в яких через двокрапку записано поряд\-ко\-вий номер запису в журналі (ціле), тип події (ERROR, WARNING, INFO, STATUS), ім'я користувача (складається з латиниці), час події (фор\-мат як у прикладі), код події (ціле).

Білі символи навколо роздільників не допускаються.

Приклад рядка файлу



\begin{verbatim}
13:ERROR:username:2022-05-05 13-26:404
\end{verbatim}


Нижче наведено код, який має замінити у файлі імена користувачів на рядок \code{unknown}.



\begin{verbatim}
regex = re.compile(r'XXX')
with open('1.txt', encoding='utf8') as f:
    lst = f.readlines()
for i in range(len(lst)):
    lst[i] = re.sub(regex, r'YYY', lst[i])
with open('1.txt', 'w', encoding='utf8') as f:
    f.writelines(lst)
\end{verbatim}







Який рядок треба записати на місце \kw{XXX} ?\par
\vspace{5mm}


Який рядок треба записати на місце \kw{YYY} ?\par


\newpage

{\begin {center}\textbf{ 314}\end {center}}\par
Файл \code{'1.txt'} містить журнал з рядками, в яких через кому записано поряд\-ко\-вий номер запису (ціле), ім'я користувача (складається з латиниці), час події (фор\-мат як у прикладі), тип події (ERROR, WARNING, INFO, STATUS), код події (ціле).

Білі символи навколо роздільників не допускаються.

Приклад рядка журналу



\begin{verbatim}
12,username,2022-05-05 13:26,ERROR,404
\end{verbatim}


Нижче наведено код, який шукає усі попередження (подія WARNING), що зустрі\-ча\-ються в журналі і відбулися під час обідньої перерви деканату (з 13,00 по 13,59 включно), та зберігає їх коди в змінній \kw{codes}.



\begin{verbatim}
codes = set()
regex = re.compile(r'XXX')
with open('1.txt', encoding='utf8') as f:
    for s in f:
        res = re.fullmatch(regex, s)
        if not res: continue
        s = ''
        codes.add(r'YYY')
\end{verbatim}







Який рядок треба записати на місце \kw{XXX} ?\par
\vspace{5mm}


Який рядок треба записати на місце \kw{YYY} ?\par


\newpage

{\begin {center}\textbf{ 315}\end {center}}\par
Файл \code{'1.txt'} містить відомість з рядками, в яких через двокрапку записано поряд\-ковий номер (ціле), назву предмету (знаків пунктуації не містить), прізвище студента, ім'я студента, номер заліковки, оцінку в балах за 100 бальною шкалою.

Білі символи навколо роздільників не допускаються.

Приклад рядка файлу



\begin{verbatim}
13:algebra:surname:name:123456:77
\end{verbatim}


Нижче наведено код, який шукає усі предмети, який складав студент із заліковкою \code{777}, та зберігає їх типи в змінній \kw{codes}.



\begin{verbatim}
codes = set()
regex = re.compile(r'XXX')
with open('1.txt', encoding='utf8') as f:
    for s in f:
        res = re.fullmatch(regex, s)
        if not res: continue
        s = ''
        codes.add(r'YYY')
\end{verbatim}





Який рядок треба записати на місце \kw{XXX} ?\par
\vspace{5mm}


Який рядок треба записати на місце \kw{YYY} ?\par


\newpage

{\begin {center}\textbf{ 316}\end {center}}\par
Файл \code{'1.txt'} містить журнал з рядками, в яких через кому записано поряд\-ко\-вий номер запису (ціле), ім'я користувача (складається з латиниці), час події (фор\-мат як у прикладі), тип події (ERROR, WARNING, INFO, STATUS), код події (ціле).

Білі символи навколо роздільників не допускаються.

Приклад рядка журналу



\begin{verbatim}
12,username,2022-05-05 13:26,ERROR,404
\end{verbatim}


Нижче наведено код, який шукає усі події, що відбулись\\ для користувача \code{anonymous}, та зберігає їх типи в змінній \kw{codes}.



\begin{verbatim}
codes = set()
regex = re.compile(r'XXX')
with open('1.txt', encoding='utf8') as f:
    for s in f:
        res = re.fullmatch(regex, s)
        if not res: continue
        s = ''
        codes.add(r'YYY')
\end{verbatim}





Який рядок треба записати на місце \kw{XXX} ?\par
\vspace{5mm}


Який рядок треба записати на місце \kw{YYY} ?\par


\newpage

{\begin {center}\textbf{ 317}\end {center}}\par
Файл \code{'1.txt'} містить журнал з рядками, в яких через двокрапку записано поряд\-ко\-вий номер запису (ціле), ім'я користувача (складається з латиниці), час події (фор\-мат як у прикладі), тип події (ERROR, WARNING, INFO, STATUS), код події (ціле).

Білі символи навколо роздільників не допускаються.

Приклад рядка журналу



\begin{verbatim}
12:username:2022-05-05 13-26:ERROR:404
\end{verbatim}


Нижче наведено код, який шукає усі події, що відбулись\\ для користувача \code{anonymous}, та зберігає їх типи в змінній \kw{codes}.



\begin{verbatim}
codes = set()
regex = re.compile(r'XXX')
with open('1.txt', encoding='utf8') as f:
    for s in f:
        res = re.fullmatch(regex, s)
        if not res: continue
        s = ''
        codes.add(r'YYY')
\end{verbatim}





Який рядок треба записати на місце \kw{XXX} ?\par
\vspace{5mm}


Який рядок треба записати на місце \kw{YYY} ?\par


\newpage

{\begin {center}\textbf{ 318}\end {center}}\par
Файл \code{'1.txt'} містить відомість з рядками, в яких через кому записано поряд\-ковий номер (ціле), назву предмету (знаків пунктуації не містить), прізвище студента, ім'я студента, номер заліковки, оцінку в балах за 100 бальною шкалою.

Білі символи навколо роздільників не допускаються.

Приклад рядка файлу



\begin{verbatim}
13,algebra,surname,name,123456,77
\end{verbatim}


Нижче наведено код, який має залишити у файлі в кожному рядку тільки номер заліковки та предмет.

Рядок з прикладу має бути замінений на такий



\begin{verbatim}
123456,algebra
\end{verbatim}




Код



\begin{verbatim}
regex = re.compile(r'XXX')
with open('1.txt', encoding='utf8') as f:
    lst = f.readlines()
for i in range(len(lst)):
    lst[i] = re.sub(regex, r'YYY', lst[i])
with open('1.txt', 'w', encoding='utf8') as f:
    f.writelines(lst)
\end{verbatim}









Який рядок треба записати на місце \kw{XXX} ?\par
\vspace{5mm}


Який рядок треба записати на місце \kw{YYY} ?\par


\newpage

{\begin {center}\textbf{ 319}\end {center}}\par
Файл \code{'1.txt'} містить журнал з рядками, в яких через кому записано поряд\-ко\-вий номер запису в журналі (ціле), тип події (ERROR, WARNING, INFO, STATUS), ім'я користувача (складається з латиниці), час події (фор\-мат як у прикладі), опис події.

Білі символи навколо роздільників не допускаються.

Приклад рядка файлу



\begin{verbatim}
13,ERROR,username,2022-05-05 13:26,something happens
\end{verbatim}


Нижче наведено код, який має залишити у файлі в кожному рядку тільки тип події та її опис.

Рядок з прикладу має бути замінений на такий



\begin{verbatim}
ERROR,something happens
\end{verbatim}




Код:



\begin{verbatim}
regex = re.compile(r'XXX')
with open('1.txt', encoding='utf8') as f:
    lst = f.readlines()
for i in range(len(lst)):
    lst[i] = re.sub(regex, r'YYY', lst[i])
with open('1.txt', 'w', encoding='utf8') as f:
    f.writelines(lst)
\end{verbatim}









Який рядок треба записати на місце \kw{XXX} ?\par
\vspace{5mm}


Який рядок треба записати на місце \kw{YYY} ?\par


\newpage

{\begin {center}\textbf{ 320}\end {center}}\par
Файл \code{'1.txt'} містить журнал з рядками, в яких через двокрапку записано поряд\-ко\-вий номер запису в журналі (ціле), тип події (ERROR, WARNING, INFO, STATUS), ім'я користувача (складається з латиниці), час події (фор\-мат як у прикладі), код події (ціле).

Білі символи навколо роздільників не допускаються.

Приклад рядка файлу



\begin{verbatim}
13:ERROR:username:2022-05-05 13-26:404
\end{verbatim}


Нижче наведено код, який шукає усі коди помилок (подія ERROR), що зустріча\-ють\-ся в журналі, та зберігає їх у змінній \kw{codes}.



\begin{verbatim}
codes = set()
regex = re.compile(r'XXX')
with open('1.txt', encoding='utf8') as f:
    for s in f:
        res = re.fullmatch(regex, s)
        if not res: continue
        s = ''
        codes.add(r'YYY')
\end{verbatim}





Який рядок треба записати на місце \kw{XXX} ?\par
\vspace{5mm}


Який рядок треба записати на місце \kw{YYY} ?\par


\newpage

{\begin {center}\textbf{ 321}\end {center}}\par
Рядки відомості через кому містять порядковий номер (ціле), назву предмету (зна\-ків пунктуації не містить), прізвище студента, ім'я студента, оцінку в балах за 100 бальною шкалою.

Білі символи навколо роздільників не допускаються.

Приклад такого рядка присвоєно змінній \kw{s}.



\begin{verbatim}
s = '13,algebra,surname,name,77'
\end{verbatim}


Нижче наведено код, за виконання якого значенням змінної \kw{out} має стати прізвище студента.

Для наведеного прикладу значенням змінної \kw{out} має стати рядок \code{surname}



\begin{verbatim}
res = re.fullmatch(r'XXX', s)
s = ''
out = r'YYY'
\end{verbatim}


Код має працювати не тільки для конкретного рядка з прикладу, але й для кожного рядка з відомості.




Який рядок треба записати на місце \kw{XXX} ?\par
\vspace{5mm}


Який рядок треба записати на місце \kw{YYY} ?\par


\newpage

{\begin {center}\textbf{ 322}\end {center}}\par
Рядки відомості через двокрапку містять порядковий номер (ціле), назву предмету (знаків пунктуації не містить), прізвище студента, ім'я студента, оцінку в балах за 100 бальною шкалою.

Білі символи навколо роздільників не допускаються.

Приклад такого рядка присвоєно змінній \kw{s}.



\begin{verbatim}
s = '13:algebra:surname:name:77'
\end{verbatim}


Нижче наведено код, за виконання якого значенням змінної \kw{out} має стати прізвище студента.

Для наведеного прикладу значенням змінної \kw{out} має стати рядок \code{surname}



\begin{verbatim}
res = re.fullmatch(r'XXX', s)
s = ''
out = r'YYY'
\end{verbatim}


Код має працювати не тільки для конкретного рядка з прикладу, але й для кожного рядка з відомості.




Який рядок треба записати на місце \kw{XXX} ?\par
\vspace{5mm}


Який рядок треба записати на місце \kw{YYY} ?\par


\newpage

{\begin {center}\textbf{ 323}\end {center}}\par
Рядки журналу через двокрапку містять тип події (ERROR, WARNING, INFO, STATUS), ім'я користувача (складається з латиниці), час події (фор\-мат як у прикладі), код події (ціле).

Білі символи навколо роздільників не допускаються.

Приклад такого рядка присвоєно змінній \kw{s}.



\begin{verbatim}
s = 'ERROR:username:2022-05-05 13-26:404'
\end{verbatim}


Нижче наведено код, за виконання якого значенням змінної \kw{out} має стати ім'я користувача.

Для наведеного прикладу значенням змінної \kw{out} має стати рядок \code{username}



\begin{verbatim}
res = re.fullmatch(r'XXX', s)
s = ''
out = r'YYY'
\end{verbatim}


Код має працювати не тільки для конкретного рядка з прикладу, але й для кожного рядка з журналу.




Який рядок треба записати на місце \kw{XXX} ?\par
\vspace{5mm}


Який рядок треба записати на місце \kw{YYY} ?\par


\newpage

{\begin {center}\textbf{ 324}\end {center}}\par
Файл \code{'1.txt'} містить журнал з рядками, в яких через двокрапку записано поряд\-ко\-вий номер запису (ціле), ім'я користувача (складається з латиниці), час події (фор\-мат як у прикладі), тип події (ERROR, WARNING, INFO, STATUS), код події (ціле).

Білі символи навколо роздільників не допускаються.

Приклад рядка журналу



\begin{verbatim}
12:username:2022-05-05 13-26:ERROR:404
\end{verbatim}


Нижче наведено код, який шукає усі попередження (подія WARNING), що зустрі\-ча\-ються в журналі і відбулися під час обідньої перерви деканату (з 13:00 по 13:59 включно), та зберігає їх коди в змінній \kw{codes}.



\begin{verbatim}
codes = set()
regex = re.compile(r'XXX')
with open('1.txt', encoding='utf8') as f:
    for s in f:
        res = re.fullmatch(regex, s)
        if not res: continue
        s = ''
        codes.add(r'YYY')
\end{verbatim}







Який рядок треба записати на місце \kw{XXX} ?\par
\vspace{5mm}


Який рядок треба записати на місце \kw{YYY} ?\par


\newpage

{\begin {center}\textbf{ 325}\end {center}}\par
Рядки журналу через кому містять ім'я користувача (складається з латиниці), час події (фор\-мат як у прикладі), тип події (ERROR, WARNING, INFO, STATUS), код події (ціле).

Білі символи навколо роздільників не допускаються.

Приклад такого рядка присвоєно змінній \kw{s}.



\begin{verbatim}
s = 'username,2022-05-05 13:26,ERROR,404'
\end{verbatim}


Нижче наведено код, за виконання якого значенням змінної \kw{out} має стати код події.

Для наведеного прикладу значенням змінної \kw{out} має стати рядок \code{404}



\begin{verbatim}
res = re.fullmatch(r'XXX', s)
s = ''
out = r'YYY'
\end{verbatim}


Код має працювати не тільки для конкретного рядка з прикладу, але й для кожного рядка з журналу.




Який рядок треба записати на місце \kw{XXX} ?\par
\vspace{5mm}


Який рядок треба записати на місце \kw{YYY} ?\par


\newpage

{\begin {center}\textbf{ 326}\end {center}}\par
Файл \code{'1.txt'} містить журнал з рядками, в яких через двокрапку записано поряд\-ко\-вий номер запису в журналі (ціле), код події (ціле), тип події (ERROR, \\WARNING, INFO, STATUS), ім'я користувача (складається з латиниці), час події (фор\-мат як у прикладі), опис події.

Білі символи навколо роздільників не допускаються.

Приклад рядка файлу



\begin{verbatim}
13:404:ERROR:username:2022-05-05 13-26:something happens
\end{verbatim}


Нижче наведено код, який має залишити у файлі в кожному рядку тільки код події та ім'я користувача. Рядок з прикладу має бути замінений на такий



\begin{verbatim}
404:username
\end{verbatim}




Код:



\begin{verbatim}
regex = re.compile(r'XXX')
with open('1.txt', encoding='utf8') as f:
    lst = f.readlines()
for i in range(len(lst)):
    lst[i] = re.sub(regex, r'YYY', lst[i])
with open('1.txt', 'w', encoding='utf8') as f:
    f.writelines(lst)
\end{verbatim}







Який рядок треба записати на місце \kw{XXX} ?\par
\vspace{5mm}


Який рядок треба записати на місце \kw{YYY} ?\par


\newpage

{\begin {center}\textbf{ 327}\end {center}}\par
Файл \code{'1.txt'} містить журнал з рядками, в яких через кому записано поряд\-ко\-вий номер запису в журналі (ціле), тип події (ERROR, WARNING, INFO, STATUS), ім'я користувача (складається з латиниці), час події (фор\-мат як у прикладі), код події (ціле).

Білі символи навколо роздільників не допускаються.

Приклад рядка файлу



\begin{verbatim}
13,ERROR,username,2022-05-05 13:26,404
\end{verbatim}


Нижче наведено код, який шукає усі коди помилок (подія ERROR), що зустріча\-ють\-ся в журналі, та зберігає їх у змінній \kw{codes}.



\begin{verbatim}
codes = set()
regex = re.compile(r'XXX')
with open('1.txt', encoding='utf8') as f:
    for s in f:
        res = re.fullmatch(regex, s)
        if not res: continue
        s = ''
        codes.add(r'YYY')
\end{verbatim}





Який рядок треба записати на місце \kw{XXX} ?\par
\vspace{5mm}


Який рядок треба записати на місце \kw{YYY} ?\par


\newpage

{\begin {center}\textbf{ 328}\end {center}}\par
Файл \code{'1.txt'} містить відомість з рядками, в яких через двокрапку записа\-но назву предмету (знаків пунктуації не містить), прізвище студента, ім'я студента, номер залі\-ков\-ки, оцінку в балах за 100 бальною шкалою.

Білі символи навколо роздільників не допускаються.

Приклад рядка файлу



\begin{verbatim}
algebra:surname:name:123456:77
\end{verbatim}


Нижче наведено код, який має замінити у файлі усі номери заліковок на рядок \code{***}.



\begin{verbatim}
regex = re.compile(r'XXX')
with open('1.txt', encoding='utf8') as f:
    lst = f.readlines()
for i in range(len(lst)):
    lst[i] = re.sub(regex, r'YYY', lst[i])
with open('1.txt', 'w', encoding='utf8') as f:
    f.writelines(lst)
\end{verbatim}







Який рядок треба записати на місце \kw{XXX} ?\par
\vspace{5mm}


Який рядок треба записати на місце \kw{YYY} ?\par


\newpage

{\begin {center}\textbf{ 329}\end {center}}\par
Файл \code{'1.txt'} містить відомість з рядками, в яких через двокрапку записано поряд\-ковий номер (ціле), назву предмету (знаків пунктуації не містить), прізвище студента, ім'я студента, номер заліковки, оцінку в балах за 100 бальною шкалою.

Білі символи навколо роздільників не допускаються.

Приклад рядка файлу



\begin{verbatim}
13:algebra:surname:name:123456:77
\end{verbatim}


Нижче наведено код, який має залишити у файлі в кожному рядку тільки номер заліковки та предмет.

Рядок з прикладу має бути замінений на такий



\begin{verbatim}
123456:algebra
\end{verbatim}




Код:



\begin{verbatim}
regex = re.compile(r'XXX')
with open('1.txt', encoding='utf8') as f:
    lst = f.readlines()
for i in range(len(lst)):
    lst[i] = re.sub(regex, r'YYY', lst[i])
with open('1.txt', 'w', encoding='utf8') as f:
    f.writelines(lst)
\end{verbatim}









Який рядок треба записати на місце \kw{XXX} ?\par
\vspace{5mm}


Який рядок треба записати на місце \kw{YYY} ?\par


\newpage

{\begin {center}\textbf{ 330}\end {center}}\par
Файл \code{'1.txt'} містить відомість з рядками, в яких через кому записа\-но назву предмету (знаків пунктуації не містить), прізвище студента, ім'я студента, номер залі\-ков\-ки, оцінку в балах за 100 бальною шкалою.

Білі символи навколо роздільників не допускаються.

Приклад рядка файлу



\begin{verbatim}
algebra,surname,name,123456,77
\end{verbatim}


Нижче наведено код, який має замінити у файлі усі номери заліковок на рядок \code{***}.



\begin{verbatim}
regex = re.compile(r'XXX')
with open('1.txt', encoding='utf8') as f:
    lst = f.readlines()
for i in range(len(lst)):
    lst[i] = re.sub(regex, r'YYY', lst[i])
with open('1.txt', 'w', encoding='utf8') as f:
    f.writelines(lst)
\end{verbatim}







Який рядок треба записати на місце \kw{XXX} ?\par
\vspace{5mm}


Який рядок треба записати на місце \kw{YYY} ?\par


\newpage

{\begin {center}\textbf{ 331}\end {center}}\par
Рядки журналу через кому містять тип події (ERROR, WARNING, INFO, STATUS), ім'я користувача (складається з латиниці), час події (фор\-мат як у прикладі), код події (ціле).

Білі символи навколо роздільників не допускаються.

Приклад такого рядка присвоєно змінній \kw{s}.



\begin{verbatim}
s = 'ERROR,username,2022-05-05 13:26,404'
\end{verbatim}


Нижче наведено код, за виконання якого значенням змінної \kw{out} має стати ім'я користувача.

Для наведеного прикладу значенням змінної \kw{out} має стати рядок \code{username}



\begin{verbatim}
res = re.fullmatch(r'XXX', s)
s = ''
out = r'YYY'
\end{verbatim}


Код має працювати не тільки для конкретного рядка з прикладу, але й для кожного рядка з журналу.




Який рядок треба записати на місце \kw{XXX} ?\par
\vspace{5mm}


Який рядок треба записати на місце \kw{YYY} ?\par


\newpage

{\begin {center}\textbf{ 332}\end {center}}\par
Рядки журналу через двокрапку містять ім'я користувача (складається з латиниці), час події (фор\-мат як у прикладі), тип події (ERROR, WARNING, INFO, STATUS), код події (ціле).

Білі символи навколо роздільників не допускаються.

Приклад такого рядка присвоєно змінній \kw{s}.



\begin{verbatim}
s = 'username:2022-05-05 13-26:ERROR:404'
\end{verbatim}


Нижче наведено код, за виконання якого значенням змінної \kw{out} має стати код події.

Для наведеного прикладу значенням змінної \kw{out} має стати рядок \code{404}



\begin{verbatim}
res = re.fullmatch(r'XXX', s)
s = ''
out = r'YYY'
\end{verbatim}


Код має працювати не тільки для конкретного рядка з прикладу, але й для кожного рядка з журналу.




Який рядок треба записати на місце \kw{XXX} ?\par
\vspace{5mm}


Який рядок треба записати на місце \kw{YYY} ?\par


\newpage

{\begin {center}\textbf{ 333}\end {center}}\par
Файл \code{'1.txt'} містить відомість з рядками, в яких через кому записа\-но поряд\-ковий номер (ціле), назву предмету (знаків пунктуації не містить), прізвище студента, ім'я студента, номер заліковки, оцінку в балах за 100 бальною шкалою.

Білі символи навколо роздільників не допускаються.

Приклад рядка файлу



\begin{verbatim}
13,algebra,surname,name,123456,77
\end{verbatim}


Нижче наведено код, який має залишити у файлі в кожному рядку тільки номер заліковки, назву предмету та оцінку в балах за 100-бальною шкалою (у заданій послі\-дов\-ності).

Рядок з прикладу має бути замінений на такий



\begin{verbatim}
123456,algebra,77
\end{verbatim}




Код



\begin{verbatim}
regex = re.compile(r'XXX')
with open('1.txt', encoding='utf8') as f:
    lst = f.readlines()
for i in range(len(lst)):
    lst[i] = re.sub(regex, r'YYY', lst[i])
with open('1.txt', 'w', encoding='utf8') as f:
    f.writelines(lst)
\end{verbatim}







Який рядок треба записати на місце \kw{XXX} ?\par
\vspace{5mm}


Який рядок треба записати на місце \kw{YYY} ?\par


\newpage

{\begin {center}\textbf{ 334}\end {center}}\par
Файл \code{'1.txt'} містить відомість з рядками, в яких через двокрапку записа\-но поряд\-ковий номер (ціле), назву предмету (знаків пунктуації не містить), прізвище студента, ім'я студента, номер заліковки, оцінку в балах за 100 бальною шкалою.

Білі символи навколо роздільників не допускаються.

Приклад рядка файлу



\begin{verbatim}
13:algebra:surname:name:123456:77
\end{verbatim}


Нижче наведено код, який має залишити у файлі в кожному рядку тільки номер заліковки, назву предмету та оцінку в балах за 100-бальною шкалою (у заданій послі\-дов\-ності).

Рядок з прикладу має бути замінений на такий



\begin{verbatim}
123456:algebra:77
\end{verbatim}




Код:



\begin{verbatim}
regex = re.compile(r'XXX')
with open('1.txt', encoding='utf8') as f:
    lst = f.readlines()
for i in range(len(lst)):
    lst[i] = re.sub(regex, r'YYY', lst[i])
with open('1.txt', 'w', encoding='utf8') as f:
    f.writelines(lst)
\end{verbatim}







Який рядок треба записати на місце \kw{XXX} ?\par
\vspace{5mm}


Який рядок треба записати на місце \kw{YYY} ?\par


\newpage

{\begin {center}\textbf{ 335}\end {center}}\par
Файл \code{'1.txt'} містить відомість з рядками, в яких через двокрапку записано поряд\-ковий номер (ціле), назву предмету (знаків пунктуації не містить), прізвище студента, ім'я студента, номер заліковки, оцінку в балах за 100 бальною шкалою.

Білі символи навколо роздільників не допускаються.

Приклад рядка файлу



\begin{verbatim}
13:algebra:surname:name:123456:77
\end{verbatim}


Нижче наведено код, який шукає усі назви предметів, що зустрічаються у файлі, та зберігає їх у змінній \id{codes}.



\begin{verbatim}
codes = set()
regex = re.compile(r'XXX')
with open('1.txt', encoding='utf8') as f:
    for s in f:
        res = re.fullmatch(regex, s)
        if not res: continue
        s = ''
        codes.add(r'YYY')
\end{verbatim}





Який рядок треба записати на місце \kw{XXX} ?\par
\vspace{5mm}


Який рядок треба записати на місце \kw{YYY} ?\par


\newpage

{\begin {center}\textbf{ 336}\end {center}}\par
Файл \code{'1.txt'} містить журнал з рядками, в яких через кому записано поряд\-ко\-вий номер запису в журналі (ціле), код події (ціле), тип події (ERROR, \\WARNING, INFO, STATUS), ім'я користувача (складається з латиниці), час події (фор\-мат як у прикладі), опис події.

Білі символи навколо роздільників не допускаються.

Приклад рядка файлу



\begin{verbatim}
13,404,ERROR,username,2022-05-05 13:26,something happens
\end{verbatim}


Нижче наведено код, який має залишити у файлі в кожному рядку тільки код події та ім'я користувача. Рядок з прикладу має бути замінений на такий



\begin{verbatim}
404,username
\end{verbatim}




Код:



\begin{verbatim}
regex = re.compile(r'XXX')
with open('1.txt', encoding='utf8') as f:
    lst = f.readlines()
for i in range(len(lst)):
    lst[i] = re.sub(regex, r'YYY', lst[i])
with open('1.txt', 'w', encoding='utf8') as f:
    f.writelines(lst)
\end{verbatim}







Який рядок треба записати на місце \kw{XXX} ?\par
\vspace{5mm}


Який рядок треба записати на місце \kw{YYY} ?\par


\newpage

{\begin {center}\textbf{ 337}\end {center}}\par
Файл \code{'1.txt'} містить відомість з рядками, в яких через кому записано поряд\-ковий номер (ціле), назву предмету (знаків пунктуації не містить), прізвище студента, ім'я студента, номер заліковки, оцінку в балах за 100 бальною шкалою.

Білі символи навколо роздільників не допускаються.

Приклад рядка файлу



\begin{verbatim}
13,algebra,surname,name,123456,77
\end{verbatim}


Нижче наведено код, який шукає усі назви предметів, що зустрічаються у файлі, та зберігає їх у змінній \id{codes}.



\begin{verbatim}
codes = set()
regex = re.compile(r'XXX')
with open('1.txt', encoding='utf8') as f:
    for s in f:
        res = re.fullmatch(regex, s)
        if not res: continue
        s = ''
        codes.add(r'YYY')
\end{verbatim}





Який рядок треба записати на місце \kw{XXX} ?\par
\vspace{5mm}


Який рядок треба записати на місце \kw{YYY} ?\par


\newpage

{\begin {center}\textbf{ 338}\end {center}}\par
Файл \code{'1.txt'} містить відомість з рядками, в яких через кому записано поряд\-ковий номер (ціле), назву предмету (знаків пунктуації не містить), прізвище студента, ім'я студента, номер заліковки, оцінку в балах за 100 бальною шкалою.

Білі символи навколо роздільників не допускаються.

Приклад рядка файлу



\begin{verbatim}
13,algebra,surname,name,123456,77
\end{verbatim}


Нижче наведено код, який шукає усі предмети, який складав студент із заліковкою \code{777}, та зберігає їх типи в змінній \kw{codes}.



\begin{verbatim}
codes = set()
regex = re.compile(r'XXX')
with open('1.txt', encoding='utf8') as f:
    for s in f:
        res = re.fullmatch(regex, s)
        if not res: continue
        s = ''
        codes.add(r'YYY')
\end{verbatim}





Який рядок треба записати на місце \kw{XXX} ?\par
\vspace{5mm}


Який рядок треба записати на місце \kw{YYY} ?\par


\newpage

{\begin {center}\textbf{ 339}\end {center}}\par
Файл \code{'1.txt'} містить журнал з рядками, в яких через кому записано поряд\-ко\-вий номер запису в журналі (ціле), тип події (ERROR, WARNING, INFO, STATUS), ім'я користувача (складається з латиниці), час події (фор\-мат як у прикладі), код події (ціле).

Білі символи навколо роздільників не допускаються.

Приклад рядка файлу



\begin{verbatim}
13,ERROR,username,2022-05-05 13:26,404
\end{verbatim}


Нижче наведено код, який має замінити у файлі імена користувачів на рядок \code{unknown}.



\begin{verbatim}
regex = re.compile(r'XXX')
with open('1.txt', encoding='utf8') as f:
    lst = f.readlines()
for i in range(len(lst)):
    lst[i] = re.sub(regex, r'YYY', lst[i])
with open('1.txt', 'w', encoding='utf8') as f:
    f.writelines(lst)
\end{verbatim}







Який рядок треба записати на місце \kw{XXX} ?\par
\vspace{5mm}


Який рядок треба записати на місце \kw{YYY} ?\par


\newpage

{\begin {center}\textbf{ 340}\end {center}}\par
Файл \code{'1.txt'} містить відомість з рядками, в яких через двокрапку записано поряд\-ковий номер (ціле), назву предмету (знаків пунктуації не містить), прізвище студента, ім'я студента, номер заліковки, оцінку в балах за 100 бальною шкалою.

Білі символи навколо роздільників не допускаються.

Приклад рядка файлу



\begin{verbatim}
13:algebra:surname:name:123456:77
\end{verbatim}


Нижче наведено код, який має в кожному рядку файлу переставити місцями назву предмету та номер заліковки.

Рядок з прикладу має бути замінений на такий



\begin{verbatim}
13:123456:surname:name:algebra:77
\end{verbatim}




Код:



\begin{verbatim}
regex = re.compile(r'XXX')
with open('1.txt', encoding='utf8') as f:
    lst = f.readlines()
for i in range(len(lst)):
    lst[i] = re.sub(regex, r'YYY', lst[i])
with open('1.txt', 'w', encoding='utf8') as f:
    f.writelines(lst)
\end{verbatim}







Який рядок треба записати на місце \kw{XXX} ?\par
\vspace{5mm}


Який рядок треба записати на місце \kw{YYY} ?\par


\newpage

{\begin {center}\textbf{ 341}\end {center}}\par
Файл \code{'1.txt'} містить відомість з рядками, в яких через двокрапку записа\-но назву предмету (знаків пунктуації не містить), прізвище студента, ім'я студента, номер залі\-ков\-ки, оцінку в балах за 100 бальною шкалою.

Білі символи навколо роздільників не допускаються.

Приклад рядка файлу



\begin{verbatim}
algebra:surname:name:123456:77
\end{verbatim}


Нижче наведено код, який має замінити у файлі усі номери заліковок на рядок \code{***}.



\begin{verbatim}
regex = re.compile(r'XXX')
with open('1.txt', encoding='utf8') as f:
    lst = f.readlines()
for i in range(len(lst)):
    lst[i] = re.sub(regex, r'YYY', lst[i])
with open('1.txt', 'w', encoding='utf8') as f:
    f.writelines(lst)
\end{verbatim}







Який рядок треба записати на місце \kw{XXX} ?\par
\vspace{5mm}


Який рядок треба записати на місце \kw{YYY} ?\par


\newpage

{\begin {center}\textbf{ 342}\end {center}}\par
Рядки журналу через кому містять ім'я користувача (складається з латиниці), час події (фор\-мат як у прикладі), тип події (ERROR, WARNING, INFO, STATUS), код події (ціле).

Білі символи навколо роздільників не допускаються.

Приклад такого рядка присвоєно змінній \kw{s}.



\begin{verbatim}
s = 'username,2022-05-05 13:26,ERROR,404'
\end{verbatim}


Нижче наведено код, за виконання якого значенням змінної \kw{out} має стати код події.

Для наведеного прикладу значенням змінної \kw{out} має стати рядок \code{404}



\begin{verbatim}
res = re.fullmatch(r'XXX', s)
s = ''
out = r'YYY'
\end{verbatim}


Код має працювати не тільки для конкретного рядка з прикладу, але й для кожного рядка з журналу.




Який рядок треба записати на місце \kw{XXX} ?\par
\vspace{5mm}


Який рядок треба записати на місце \kw{YYY} ?\par


\newpage

{\begin {center}\textbf{ 343}\end {center}}\par
Файл \code{'1.txt'} містить журнал з рядками, в яких через кому записано поряд\-ко\-вий номер запису (ціле), ім'я користувача (складається з латиниці), час події (фор\-мат як у прикладі), тип події (ERROR, WARNING, INFO, STATUS), код події (ціле).

Білі символи навколо роздільників не допускаються.

Приклад рядка журналу



\begin{verbatim}
12,username,2022-05-05 13:26,ERROR,404
\end{verbatim}


Нижче наведено код, який шукає усі попередження (подія WARNING), що зустрі\-ча\-ються в журналі і відбулися під час обідньої перерви деканату (з 13,00 по 13,59 включно), та зберігає їх коди в змінній \kw{codes}.



\begin{verbatim}
codes = set()
regex = re.compile(r'XXX')
with open('1.txt', encoding='utf8') as f:
    for s in f:
        res = re.fullmatch(regex, s)
        if not res: continue
        s = ''
        codes.add(r'YYY')
\end{verbatim}







Який рядок треба записати на місце \kw{XXX} ?\par
\vspace{5mm}


Який рядок треба записати на місце \kw{YYY} ?\par


\newpage

{\begin {center}\textbf{ 344}\end {center}}\par
Рядки журналу через двокрапку містять ім'я користувача (складається з латиниці), час події (фор\-мат як у прикладі), тип події (ERROR, WARNING, INFO, STATUS), код події (ціле).

Білі символи навколо роздільників не допускаються.

Приклад такого рядка присвоєно змінній \kw{s}.



\begin{verbatim}
s = 'username:2022-05-05 13-26:ERROR:404'
\end{verbatim}


Нижче наведено код, за виконання якого значенням змінної \kw{out} має стати код події.

Для наведеного прикладу значенням змінної \kw{out} має стати рядок \code{404}



\begin{verbatim}
res = re.fullmatch(r'XXX', s)
s = ''
out = r'YYY'
\end{verbatim}


Код має працювати не тільки для конкретного рядка з прикладу, але й для кожного рядка з журналу.




Який рядок треба записати на місце \kw{XXX} ?\par
\vspace{5mm}


Який рядок треба записати на місце \kw{YYY} ?\par


\newpage

{\begin {center}\textbf{ 345}\end {center}}\par
Файл \code{'1.txt'} містить журнал з рядками, в яких через двокрапку записано поряд\-ко\-вий номер запису в журналі (ціле), тип події (ERROR, WARNING, INFO, STATUS), ім'я користувача (складається з латиниці), час події (фор\-мат як у прикладі), код події (ціле).

Білі символи навколо роздільників не допускаються.

Приклад рядка файлу



\begin{verbatim}
13:ERROR:username:2022-05-05 13-26:404
\end{verbatim}


Нижче наведено код, який має замінити у файлі імена користувачів на рядок \code{unknown}.



\begin{verbatim}
regex = re.compile(r'XXX')
with open('1.txt', encoding='utf8') as f:
    lst = f.readlines()
for i in range(len(lst)):
    lst[i] = re.sub(regex, r'YYY', lst[i])
with open('1.txt', 'w', encoding='utf8') as f:
    f.writelines(lst)
\end{verbatim}







Який рядок треба записати на місце \kw{XXX} ?\par
\vspace{5mm}


Який рядок треба записати на місце \kw{YYY} ?\par


\newpage

{\begin {center}\textbf{ 346}\end {center}}\par
Рядки журналу через кому містять тип події (ERROR, WARNING, INFO, STATUS), ім'я користувача (складається з латиниці), час події (фор\-мат як у прикладі), код події (ціле).

Білі символи навколо роздільників не допускаються.

Приклад такого рядка присвоєно змінній \kw{s}.



\begin{verbatim}
s = 'ERROR,username,2022-05-05 13:26,404'
\end{verbatim}


Нижче наведено код, за виконання якого значенням змінної \kw{out} має стати ім'я користувача.

Для наведеного прикладу значенням змінної \kw{out} має стати рядок \code{username}



\begin{verbatim}
res = re.fullmatch(r'XXX', s)
s = ''
out = r'YYY'
\end{verbatim}


Код має працювати не тільки для конкретного рядка з прикладу, але й для кожного рядка з журналу.




Який рядок треба записати на місце \kw{XXX} ?\par
\vspace{5mm}


Який рядок треба записати на місце \kw{YYY} ?\par


\newpage

{\begin {center}\textbf{ 347}\end {center}}\par
Файл \code{'1.txt'} містить відомість з рядками, в яких через двокрапку записано поряд\-ковий номер (ціле), назву предмету (знаків пунктуації не містить), прізвище студента, ім'я студента, номер заліковки, оцінку в балах за 100 бальною шкалою.

Білі символи навколо роздільників не допускаються.

Приклад рядка файлу



\begin{verbatim}
13:algebra:surname:name:123456:77
\end{verbatim}


Нижче наведено код, який має в кожному рядку файлу переставити місцями назву предмету та номер заліковки.

Рядок з прикладу має бути замінений на такий



\begin{verbatim}
13:123456:surname:name:algebra:77
\end{verbatim}




Код:



\begin{verbatim}
regex = re.compile(r'XXX')
with open('1.txt', encoding='utf8') as f:
    lst = f.readlines()
for i in range(len(lst)):
    lst[i] = re.sub(regex, r'YYY', lst[i])
with open('1.txt', 'w', encoding='utf8') as f:
    f.writelines(lst)
\end{verbatim}







Який рядок треба записати на місце \kw{XXX} ?\par
\vspace{5mm}


Який рядок треба записати на місце \kw{YYY} ?\par


\newpage

{\begin {center}\textbf{ 348}\end {center}}\par
Рядки відомості через кому містять порядковий номер (ціле), пріз\-ви\-ще студен\-та, ім'я студента, номер заліковки, оцінку в балах за 100 бальною шкалою.

Білі символи навколо роздільників не допускаються.

Приклад такого рядка присвоєно змінній \kw{s}.



\begin{verbatim}
s = '13,surname,name,123456,77'
\end{verbatim}


Нижче наведено код, за виконання якого для рядка \kw{s} значенням змінної \kw{out} має стати номер заліковки.

Для наведеного прикладу значенням змінної \kw{out} має стати рядок \code{123456} .



\begin{verbatim}
res = re.fullmatch(r'XXX', s)
s = ''
out = r'YYY'
\end{verbatim}


Код має працювати не тільки для конкретного рядка з прикладу, але й для кожного рядка з відомості.




Який рядок треба записати на місце \kw{XXX} ?\par
\vspace{5mm}


Який рядок треба записати на місце \kw{YYY} ?\par


\newpage

{\begin {center}\textbf{ 349}\end {center}}\par
Файл \code{'1.txt'} містить відомість з рядками, в яких через двокрапку записано поряд\-ковий номер (ціле), назву предмету (знаків пунктуації не містить), прізвище студента, ім'я студента, номер заліковки, оцінку в балах за 100 бальною шкалою.

Білі символи навколо роздільників не допускаються.

Приклад рядка файлу



\begin{verbatim}
13:algebra:surname:name:123456:77
\end{verbatim}


Нижче наведено код, який шукає усі назви предметів, що зустрічаються у файлі, та зберігає їх у змінній \id{codes}.



\begin{verbatim}
codes = set()
regex = re.compile(r'XXX')
with open('1.txt', encoding='utf8') as f:
    for s in f:
        res = re.fullmatch(regex, s)
        if not res: continue
        s = ''
        codes.add(r'YYY')
\end{verbatim}





Який рядок треба записати на місце \kw{XXX} ?\par
\vspace{5mm}


Який рядок треба записати на місце \kw{YYY} ?\par


\newpage

{\begin {center}\textbf{ 350}\end {center}}\par
Файл \code{'1.txt'} містить журнал з рядками, в яких через двокрапку записано поряд\-ко\-вий номер запису (ціле), ім'я користувача (складається з латиниці), час події (фор\-мат як у прикладі), тип події (ERROR, WARNING, INFO, STATUS), код події (ціле).

Білі символи навколо роздільників не допускаються.

Приклад рядка журналу



\begin{verbatim}
12:username:2022-05-05 13-26:ERROR:404
\end{verbatim}


Нижче наведено код, який шукає усі події, що відбулись\\ для користувача \code{anonymous}, та зберігає їх типи в змінній \kw{codes}.



\begin{verbatim}
codes = set()
regex = re.compile(r'XXX')
with open('1.txt', encoding='utf8') as f:
    for s in f:
        res = re.fullmatch(regex, s)
        if not res: continue
        s = ''
        codes.add(r'YYY')
\end{verbatim}





Який рядок треба записати на місце \kw{XXX} ?\par
\vspace{5mm}


Який рядок треба записати на місце \kw{YYY} ?\par


\newpage

